\programdayheader{Tuesday}
\daypalettebanner{Tuesday}

\hypertarget{poster-abstract-talk-1}{}\label{poster-abstract-page-talk-1}
\subsection*{Poster 1. On the Information Content of Ariel Transmission Spectra: Reassessing the Tier System}
\textbf{Presenter:} Michael Radica\\
\textbf{Poster number:} 1\\
\textbf{Fields:} Exoplanets $\cdot$ Next generation observatories

The European Space Agency's Ariel mission will conduct a survey of the atmospheric properties of exoplanets around bright stars. The mission is nominally divided into three Tiers. The Tier 1 survey will consist of low-precision observations of $\sim$1000 planets, with a subset of these included in the higher-precision Tier 2 survey expected to be necessary for atmospheric characterization. Tier 3 will be repeated observations of a small number of benchmark planets. Though previous studies have assessed the ability of Ariel to uncover population-level trends, they have generally presupposed a given Tier. Here we seek to interrogate those assumptions and assess the information content of Ariel spectra as a function of Tier for three benchmark planets: a hot-Saturn, warm-Neptune, and temperate sub-Neptune. In this talk I will outline the results from a comprehensive analysis of simulated Ariel spectra and assess which chemical species are potentially detectable as a function of Tier. We find that for a hot-Saturn, Tier 1-quality observations are sufficient for $\sim$1.5dex constraints on H2O and CO2, irrespective of the presence of clouds --- meaning important chemical insights are already obtainable in the Tier 1 survey. Moving to Tier 2 results in an incremental increase in precision and other molecules (e.g., SO2, CO) becoming detectable in certain scenarios. Tier 2 or better spectra are necessary to identify molecules in either a warm-Neptune or temperate sub-Neptune. However, the number of transits necessary to reach this precision may be prohibitive to include comparable targets in the Tier 2 survey. These results demonstrate the potential for Ariel to transform comparative exoplanetology and will help define the scientific objectives of the Tiered survey system.

\hypertarget{poster-abstract-talk-303}{}\label{poster-abstract-page-talk-303}
\subsection*{Poster 2. Exoplanet orbital evolution and connections with transit durations and variations}
\textbf{Presenter:} Aaron Boley\\
\textbf{Poster number:} 2\\
\textbf{Field:} Exoplanets

Short period exoplanets offer exciting science opportunities for exploring planetary interiors, planet-planet interactions, and planet-star dynamical evolution. Long-term (secular) transit timing trends, including Transit Timing Variations (TTVs) and Transit Duration Variations (TDVs), can reveal such interactions. While short-term and periodic TTVs have long been successfully used to explore near resonant planet-planet interactions and even discover non-transiting planets, observing secular variations has only relatively recently become possible, which reveal different dynamics, including the possibility of tidal decay and orbital precession. This in turn can be used to constrain, among other things, possible planetary interior structures and energy dissipation in the system. Different types of orbital evolution can lead to a range of timing trend possibilities, with various mechanisms corresponding to correlations between transit mid-times, eclipses, and durations. We present a general discussion of secular exoplanet TTVs and TDVs, highlight their observational signatures, and describe what such interactions can tell us about their host systems.

\hypertarget{poster-abstract-talk-237}{}\label{poster-abstract-page-talk-237}
\subsection*{Poster 3. Mapping the Atmosphere of TOI-561 b: First JWST/NIRSpec Phase Curve of Ultra-Short Period Planet with a Thick Atmosphere}
\textbf{Presenter:} Samuel Boucher\\
\textbf{Poster number:} 3\\
\textbf{Field:} Exoplanets

Ultra-short period rocky planets are expected to be stripped of gaseous envelopes due to intense stellar irradiation. Surprisingly, recent eclipse observations of TOI-561 b revealed a $\sim$3-5\,$\mu$m emission spectrum best explained by a thick volatile secondary atmosphere, thereby challenging our understanding of atmospheric outgassing and retention on lava planets. TOI-561 b stands out as the lowest-density lava world known, with a density analogous to the iconic 55 Cnc e, making it a unique laboratory for studying volatile retention under extreme conditions. However, the low dayside emission observed in eclipse-only data is degenerate: it could be explained by either high albedo or efficient heat recirculation to the nightside. To break this degeneracy and characterize this unexpected atmosphere, we present the first phase curve observations of a low-density lava world: multi-orbit JWST/NIRSpec G395H spectroscopic phase curve observations of TOI-561 b. Our phase curve maps the planet's thermal emission as a function of longitude, directly constraining the heat recirculation efficiency and Bond albedo. This thorough analysis confirms the high albedo inferred from previous eclipse-only results while also quantifying the planet's energy redistribution. Our observations have important implications for understanding atmosphere-interior interactions on rocky worlds in extreme environments and, more importantly, reinforce that volatile outgassing could play an important role in the presence of atmospheres on terrestrial worlds.

\hypertarget{poster-abstract-talk-53}{}\label{poster-abstract-page-talk-53}
\subsection*{Poster 4. TRAPPIST-1 in High Resolution: Constraining Rocky Exoplanet Atmospheres Amid Systematics}
\textbf{Presenter:} Mathis Bouffard\\
\textbf{Poster number:} 4\\
\textbf{Field:} Exoplanets

The TRAPPIST-1 system is a benchmark for the atmospheric characterization of rocky exoplanets through transit spectroscopy. However, no atmosphere has yet been detected with the JWST, largely due to the intense stellar activity of the M8 host. Ground-based high-resolution spectroscopy offers an alternative pathway to study these atmospheres, but it is made more difficult by the low orbital velocities and short transit durations of the planets and the weak transmission signals expected from high-mean-molecular-weight atmospheres. To understand the lower limits of detectable atmospheres with current ground-based facilities, we leverage an extensive dataset from the SPIRou spectrograph at the Canada-France-Hawaii Telescope. Our goal is to constrain the abundances of water and methane in the atmospheres of TRAPPIST-1 b and e, while mitigating stellar and instrumental systematics. SPIRou is affected by persistence: residual signal from preceding observations can contaminate data and introduce spurious spectral features. We analyze 11 transits of planet b and 2 of planet e using the open-source pipeline STARSHIPS. After correcting for persistence and telluric contamination, we cross-correlate the observations with atmospheric models. We perform Bayesian retrievals to extract constraints on atmospheric parameters. By comparing results with and without persistence correction, we demonstrate the importance of accounting for this instrumental effect. I will present upper limits on the abundances of H2O and CH4 in the atmospheres of TRAPPIST-1 b and e. These constraints will help eliminate specific atmospheric scenarios for the TRAPPIST-1 planets, advancing the search for habitable worlds and informing the design of future observing programs with the ELT.

\hypertarget{poster-abstract-talk-312}{}\label{poster-abstract-page-talk-312}
\subsection*{Poster 5. Modeling Detector Systematics in the Search for Rocky World Atmospheres with JWST MIRI}
\textbf{Presenter:} Nicholas Connors\\
\textbf{Poster number:} 5\\
\textbf{Field:} Exoplanets

JWST has enabled us to observe the compositions of rocky exoplanets in unprecedented detail, but no conclusive evidence has yet been found for an atmosphere on an Earth-sized planet outside of our solar system. Close-in rocky worlds around M dwarfs are ideal targets for atmospheric detection due to their favourable geometry for transit observations, assuming the activity from their host stars does not prevent atmospheric retention. To investigate this, measurement of thermal emission at 15 micron using JWST MIRI can be used to detect the presence of atmospheric CO2. This technique has been used in a number of previous JWST programs (GTO-1279, GO-2304, GO-3730) and in the on-going 500 hour Rocky Worlds DDT. Here, I present a uniform reanalysis of previously published data and of newly released DDT data using a new data-reduction technique called Frame-Normalized Principal Component Analysis (FN-PCA). This technique works to find and account for detector-level systematic effects inherent in the data without accidentally removing any relevant astrophysical signal. I will present benefits and potential limitations of this technique, and discuss the potential compositions of the rocky worlds (TRAPPIST-1b, TRAPPIST-1c, LHS 1478b, TOI-1468b, LHS 1140c, LTT 3780b, and GJ 3929b) we have analyzed.

\hypertarget{poster-abstract-talk-56}{}\label{poster-abstract-page-talk-56}
\subsection*{Poster 6. Population-Level Atmospheric Characterization of Ultra-Cool Dwarfs with SPHEREx}
\textbf{Presenter:} Louis-Philippe Coulombe\\
\textbf{Poster number:} 6\\
\textbf{Fields:} Exoplanets $\cdot$ Stars and Stellar cluster formation

Despite decades of surveys, our understanding of ultra-cool dwarf (UCD) atmospheres remains incomplete. Most studies have focused on individual objects, producing detailed but isolated results prone to degeneracies caused by the intrinsic diversity of these atmospheres. What is missing is a systematic, population-level approach that links atmospheric diversity to fundamental parameters such as age, mass, and temperature. We present an in-depth atmospheric analysis of UCD spectra from the SPHEREx mission, which provides continuous spectrophotometry spanning 0.755 microns and covers the bulk of these objects' spectral energy distribution down to late-T dwarfs. The data are reduced using our custom-built, open-source SPIFF pipeline, which has so far extracted spectra for over a thousand UCDs. Our spectral extraction and atmospheric retrieval methods are benchmarked using the well-studied brown dwarf SIMP 0136, which has been extensively characterized with both ground- and space-based observatories. We find that its SPHEREx spectrum is consistent with previous observations and enables constraints on key CNOS-bearing species (e.g., H$_2$O, CH$_4$, CO, CO$_2$, NH$_3$, and H$_2$S) as well as metal hybrides (CrH and FeH), with precision at the 0.010.1 dex level. Finally, we extend this analysis to a sample of L-, T-, and Y-type brown dwarfs to identify trends in thermal structure, elemental and molecular abundances, and cloud properties. These results demonstrate the potential of SPHEREx for population-level studies of UCD atmospheres, allowing us to map for the first time the physical, chemical, and cloud condensation transitions shaping the M/L/T/Y spectral sequence, and enabling a direct comparison with exoplanet and Solar System gas giant atmospheres.

\hypertarget{poster-abstract-talk-261}{}\label{poster-abstract-page-talk-261}
\subsection*{Poster 7. First Look at the James Webb Space Telescope Lava Planet Survey}
\textbf{Presenter:} Lina D'Aoust\\
\textbf{Poster number:} 7\\
\textbf{Field:} Exoplanets

Observational exoplanetary science is driven by incredible feats of engineering like the James Webb Space Telescope (JWST). The latter is revolutionising the field by going beyond binary detections, and rather characterizing exoplanet atmospheres' chemical inventories, and in some cases, providing two-dimensional climate maps of cloud coverage and temperature. Lava planets are small terrestrial planets which, due to their very short orbital periods and proximity to their star, have a molten rock surface or more eloquently, magma oceans. They are interesting study subjects since they probe the limit of exoplanet parameters by being the smallest known planets with the shortest orbital periods. Learning more about them thus tests the limits of theoretical models both about planet formation, and about the behaviour of materials and atmospheres in extreme conditions. Studying lava planets simultaneously sheds light on the makeup and early evolution of Earth-like planets, since it is believed that young terrestrial planets undergo a lava world phase. JWST MIRI observations of five lava planets are trickling in (from March to September 2026), and I report a first look at the programs data. Notably, white light curves showcasing these exoplanets' transits and eclipses were generated using both the Eureka! and exoTEDRF JWST data reduction pipelines. I present complete phase curves and fitted planetary parameters using Markov Chain Monte Carlo (MCMC) routines. The goal of this lava planet survey is to study their atmospheric makeup, their surface's thermal emission, and the interaction between the two. I present preliminary results, and how they can be interpreted in light of each other in order to make statements about evolutionary and cooling timescales of Earth-like planets.

\hypertarget{poster-abstract-talk-32}{}\label{poster-abstract-page-talk-32}
\subsection*{Poster 8. From Day to Night: High-Resolution Transmission and Emission Spectroscopy of Ultra-Hot Jupiters with the SPECTRE-GHOST Survey}
\textbf{Presenter:} Emily Deibert\\
\textbf{Poster number:} 8\\
\textbf{Field:} Exoplanets

The atmospheres of ultra-hot Jupiters (UHJs) are examples of physics and chemistry at the extreme. With high stellar irradiation levels, equilibrium temperatures exceeding 2200 K, and orbital periods of only a few days, these exoplanets have permanent, scorching daysides dominated by molecular dissociation and atomic ionization, and much cooler nightsides exhibiting recombination and rainout. Dynamical processes such as atmospheric re-circulation, vertical mixing, and global heat transport, as well as the unique properties of each system, dictate how nightside chemistry impacts the dayside, and vice versa. Today, despite progress in the field, major questions remain regarding the formation, evolution, and atmospheric physics of these extreme worlds. In this presentation, I will introduce the ongoing SPECTRE-GHOST (Spectroscopically Probing Extremely hot Climates of TRansiting Exoplanets with GHOST) survey, which aims to systematically compare chemical and dynamical processes in the daysides and terminator regions of a broad sample of UHJs using the high-resolution GHOST optical spectrograph at the Gemini South Observatory. UHJs are prime candidates for in-depth characterization via high-resolution optical spectroscopy due to the rich forest of atomic and ionic spectral features accessible at these wavelengths, and GHOST's high spectral resolution and broad optical coverage make it ideally suited for this work. I will present the first results from the survey, including time-resolved detections of multiple species in a range of UHJ atmospheres, and discuss the implications of these early results on our understanding of the UHJ population. I will also outline plans for the remainder of this multi-year survey.

\hypertarget{poster-abstract-talk-71}{}\label{poster-abstract-page-talk-71}
\subsection*{Poster 9. Exoplanet Atmospheres as Rosetta Stones: Modeling Exoplanetary Origins}
\textbf{Presenter:} James C. Dufresne\\
\textbf{Poster number:} 9\\
\textbf{Field:} Exoplanets

Every exoplanet system we have discovered looks vastly different from our own, and we currently do not know why. Understanding how planets form, starting from dust particles to massive planets is crucial to solve this problem. Since this process takes place over millions of years, we cannot watch systems evolve in real time and must infer their formation mechanisms from current-day observations.

However, there is a significant gap between what modern telescopes can measure and what current planet formation models can predict. The ARIEL telescope will launch in 2029 and provide measurements of atmospheric abundances for thousands of exoplanets with the aim of linking these atmospheric signatures to formation and evolution. However no existing model can yet connect atmospheric composition to the physics of the evolving disk in a self-consistent way. ARIELs maximal scientific return depends critically on a survey strategy informed by such predictive models.

We propose the development of a self-consistent planetary formation model which can predict the atmospheric composition of planets. Most current models treat the protoplanetary disks orbiting young stars as purely viscous and neglect the role of magnetically driven disk winds, which recent magnetohydrodynamical studies show can dominate angular-momentum transport in a way that aligns better with observations. Many models also tend to oversimplify the growth process by ignoring pebble accretion. Our model will bridge this gap by combining both viscous and wind-driven disk evolution along with pebble accretion and tracking the evolving chemical composition of solids as they grow. Population synthesis will allow us to quantify our model's sensitivity and predict ARIEL observables as a function of planetary properties.

\hypertarget{poster-abstract-talk-287}{}\label{poster-abstract-page-talk-287}
\subsection*{Poster 10. Combining Transits with Astrometry and RV: New and Updated Orbital Constraints for Multiple Transiting Systems}
\textbf{Presenter:} Andre Fogal\\
\textbf{Poster number:} 10\\
\textbf{Field:} Exoplanets

Constraining orbital parameters of exoplanets is critical to understand the dynamics of these systems and for informing follow-up work such as direct imaging. Exoplanet transits offer a powerful glimpse into the motion of the planets, allowing us to constrain the inclination and period. In combination with other data, namely radial velocity (RV) and astrometric motion of the star, we can perform Markov Chain Monte Carlo fits to determine orbital posterior predictions. In this work, we will present joint models for transiting planets detected by the Transiting Exoplanet Survey Satellite with our novel calibrated astrometry catalog G23H, which extracts maximum information from Gaia data releases 2 and 3 along with Hipparcos astrometric data. With the addition of archival RV data, these joint models allow us to constrain the full orbital properties of these systems. We will present new and updated orbit fits for multiple systems from the Bayesian inference code Octofitter, performing joint transit-astrometry-RV fits for our systems of interest. We will show results demonstrating the power of combining transit data with other common exoplanet data for determining orbital and planetary characteristics. Some systems of study, such as TOI-201, present the intriguing possibility of imaging follow-up with next-generation direct imaging instruments, marking the first system to be detected both via transmission and emission. These tools in combination with the upcoming Gaia Data Release 4 present an exciting new frontier in the detection and study of exoplanets.

\hypertarget{poster-abstract-talk-42}{}\label{poster-abstract-page-talk-42}
\subsection*{Poster 11. Atmospheric characterization of two low-density Jupiters with NIRPS transit spectroscopy: warm WASP-80 b and hot WASP-193 b}
\textbf{Presenter:} Frédéric Genest\\
\textbf{Poster number:} 11\\
\textbf{Field:} Exoplanets

High-resolution spectroscopy is a valuable tool for atmospheric characterization of exoplanets, complementary to JWST. Using the 5-year GTO program allotted to the NIRPS near-infrared spectrograph on ESOs 3.6 m telescope at La Silla, many dayside and transit observations are being collected for planets of interest.

We present atmospheric characterization efforts for two puffy hot/warm Jupiter planets: WASP-193 b (T $\sim$ 1250K) and WASP-80 b (T $\sim$ 825 K). In both cases, we work with the telluric-corrected spectra from the APERO pipeline and obtain transit spectra with the open-source STARSHIPS software. STARSHIPS includes stellar spectrum removal, residual telluric masking and stellar residual removal using a PCA-based method.

In the case of WASP-193 b, we present the ongoing analysis of 4 (up to 6) transits observed with NIRPS. We obtain a robust water detection in cross correlation. We conduct a full retrieval analysis on these transits using STARSHIPS and petitRADTRANS. As of now, this is the first in-depth atmospheric characterization work for this planet. We discuss how WASP-193 b compares to other hot Jupiters.

For WASP-80 b, we present the completed joint analysis of 4 NIRPS transits together with a previously published JWST/NIRCam transit spectrum from the literature. We performed separated and joint retrievals on these two datasets and find that by combining them, we can constrain the abundance of water (from the NIRPS data) and methane (from the NIRCam data) simultaneously. We find strict constraints for a cold temperature profile along the terminator. We discuss the significance of the methane non-detection in the NIRPS data, the unexpected PT profile constraints, and explore the additional constraints we can put on abundance ratios (C/O, metallicity).

\hypertarget{poster-abstract-talk-152}{}\label{poster-abstract-page-talk-152}
\subsection*{Poster 12. Drivers of protocluster-to-cluster evolution with the FLAMINGO simulations}
\textbf{Presenter:} Syeda Lammim Ahad\\
\textbf{Poster number:} 12\\
\textbf{Fields:} Galaxy $\cdot$ Galaxy Clusters $\cdot$ Next generation observatories

Galaxy clusters arise from overdense regions (protoclusters) at high redshift, yet simulations suggest that their progenitors do not always trace cosmic density peaks, implying that other large-scale environmental factors may drive their growth. We use the FLAMINGO simulations to study how the cosmic webs filamentary structure and large-scale environment of protoclusters shape galaxy cluster assembly across cosmic time. We measure filament properties and mass distribution around protoclusters and follow their descendants to assess their impact on the diversity of the cluster mass-accretion histories from z=6 to z=0. In this talk, I will show i) how filament properties vary across protocluster populations and epochs, ii) the impact of large-scale protocluster environment on the descendant cluster mass, and iii) how these trends enable simulation-anchored mappings from high-redshift observables to low-redshift cluster growth of protocluster candidates identified with JWST and Euclid.

\hypertarget{poster-abstract-talk-45}{}\label{poster-abstract-page-talk-45}
\subsection*{Poster 13. The Bright Future of the Dark and Dim Universe with the SKA}
\textbf{Presenter:} Nathan Deg\\
\textbf{Poster number:} 13\\
\textbf{Fields:} Galaxy $\cdot$ Next generation observatories

Dark and dim galaxies provide critical probes of dark matter, structure formation, and cosmology. In particular, gas-rich dark and dim galaxies provide the ability to constrain a variety of cosmological questions and measurements, such as the halo mass function, the formation of primordial halos, outliers from the baryonic Tully Fisher Relation (bTFR), and the inner dark matter density distribution of galaxies. Probing these on a statistical basis is challenging with current observatories due to limitations in sensitivity and resolution, although significant strides are being made. This will dramatically change with the advent of the SKA. In this talk I will present forecasts for the SKA applied to the dark and dim Universe, ranging from the prospects of detecting gas-rich primordial dark galaxies, to mapping the population of ultra diffuse dwarf galaxies in the field, to probing the full low mass star forming population of galaxies. I will also present recent work in this field, focused on diffuse dwarfs and the bTFR, highlighting current key results and how the SKA will illuminate our knowledge of the dark and dim Universe.

\hypertarget{poster-abstract-talk-297}{}\label{poster-abstract-page-talk-297}
\subsection*{Poster 14. Galaxy quenching mechanisms at cosmic noon}
\textbf{Presenter:} Laya Ghodsi\\
\textbf{Poster number:} 14\\
\textbf{Fields:} Galaxy $\cdot$ Gas, dust and the ISM

Understanding how massive galaxies shut down their star formation remains one of the central questions in galaxy evolution. Quiescent galaxies already exist in significant numbers by z$\sim$2, implying that powerful quenching mechanisms must operate rapidly at early cosmic times. Gravitational lensing provides a unique opportunity to study these systems in detail. MRG-M0138 is a strongly lensed, massive quiescent galaxy at z$\sim$2 that has been extensively studied through deep spectroscopy and imaging, revealing an old stellar population and unusually detailed chemical abundance constraints for a galaxy at this epoch. In this work, I present a new analysis of the radio and (sub)millimetre continuum emission from MRG-M0138 using observations from the VLA and ALMA. We detect a bright continuum source at 150\,GHz that is spatially coincident with the galaxy and exhibits a spectral energy distribution consistent with synchrotron emission. The radio luminosity and spectral properties suggest that the emission is likely powered by a low-level active galactic nucleus. Combined with recent dynamical measurements of an overmassive black hole in this system, these results point toward a scenario in which AGN feedback may play a key role in maintaining the quiescent state of the galaxy. MRG-M0138 therefore provides a rare window into the physical mechanisms responsible for quenching massive galaxies at cosmic noon. These observations highlight the importance of deep radio and millimetre data for identifying hidden AGN activity and constraining the role of feedback in the early evolution of massive galaxies.

\hypertarget{poster-abstract-talk-202}{}\label{poster-abstract-page-talk-202}
\subsection*{Poster 15. Identifying backsplash galaxies with machine learning}
\textbf{Presenter:} Roan Haggar\\
\textbf{Poster number:} 15\\
\textbf{Fields:} Galaxy $\cdot$ Galaxy Clusters

The evolution of galaxies is dependent on cosmic environment: whether the galaxies are isolated, or live in dense regions such as galaxy clusters. The subtlety of this is that both a galaxy's present-day environment, and its previous environments, all play a role in its evolution. Backsplash galaxies are a key example of this -- galaxies that have previously passed through the dense centre of a galaxy cluster, but now reside in the less-dense cluster outskirts. Identifying these galaxies is difficult, but is crucial in understanding whether a galaxy's properties are caused by pre-processing en-route to a cluster, or by a passage through the cluster centre.

We have developed a machine learning-powered model, trained on simulations of large galaxy clusters, that is able to distinguish between backsplash and infalling galaxies, with a significantly higher accuracy than existing methods. Crucially, this model only uses galaxy properties that are observationally measurable, such as their present-day line-of-sight velocities or 2D positions. This tool can therefore be used to identify backsplash galaxies in astronomical observations; we will demonstrate how identifying backsplash galaxies in the Virgo cluster can provide insights to the connections between galaxy evolution and environmental history. These insights will only get better in the near future, with upcoming observations of galaxy clusters from large spectroscopic surveys, like DESI and the wide-field cluster survey from WEAVE.

\hypertarget{poster-abstract-talk-252}{}\label{poster-abstract-page-talk-252}
\subsection*{Poster 16. VLA 3\,GHz Radio Continuum Mapping of Nearby PHANGS Galaxies: Tracing the Life Cycle of Star Formation with JWST and VLT}
\textbf{Presenter:} Hamid Hassani\\
\textbf{Poster number:} 16\\
\textbf{Fields:} Galaxy $\cdot$ Gas, dust and the ISM $\cdot$ Stars and Stellar cluster formation

Massive star clusters evolve from embedded, optically thick phases to exposed, optically bright stages shaped by stellar winds, supernovae, and cosmic rays, and understanding this evolution is key to studying how stellar feedback regulates future star formation. We present high-resolution 3\,GHz radio continuum observations of nearby galaxies obtained with the VLA. With physical resolutions of roughly 60\,pc, these data allow us to resolve individual star-forming regions. We identify 209 compact radio sources with a median radio luminosity of about $6.3 \times 10^{19}$ W/Hz. Using extinction-corrected H$\alpha$ maps as a tracer of free-free emission, we separate the radio continuum into thermal and non-thermal components and find that the thermal fraction varies strongly across galactic environments, ranging from roughly 40\% in galaxy centers to nearly 100\% in spiral arms. About 50\% of the radio sources have H$\alpha$ counterparts tracing exposed star formation, while roughly 75\% are associated with 21\,$\mu$m mid-infrared emission. In addition, we find that about 10\% of the radio detections correspond to embedded regions that lack optical counterparts but exhibit strong 3.3\,$\mu$m PAH emission. Stellar ages show that very young regions (13\,Myr) are dominated by thermal emission ($\sim$85\%). At 36\,Myr, the thermal fraction drops to $\sim$70\% and the radio luminosity per stellar mass decreases by about a factor of two. At later stages (610\,Myr), synchrotron emission increases, with the non-thermal fraction rising above 70\%, marking the onset of supernova activity. These results show that radio continuum observations trace the timeline of star-forming region evolution from embedded clusters to supernova-dominated synchrotron stages and constrain how stellar feedback regulates the energy balance in galaxies.

\hypertarget{poster-abstract-talk-40}{}\label{poster-abstract-page-talk-40}
\subsection*{Poster 17. Canadian Hydrogen Observatory and Radio-transient Detector Galaxy Search Pipeline}
\textbf{Presenter:} Hans Hopkins\\
\textbf{Poster number:} 17\\
\textbf{Fields:} Galaxy $\cdot$ Instrumentation $\cdot$ Next generation observatories

The Canadian Hydrogen Observatory and Radio-transient Detector (CHORD) is an under-construction redundant-array driftscan radio interferometer. One of its science aims is to produce a large catalogue of spectra of HI emitting galaxies. CHORD will have the sensitivity to detect objects with HI mass down to 10$^5$-10$^6$ solar masses because of its repeated baseline design. Its rectangular grid structure will allow for many beams to record the sky at once, but in exchange it introduces strong spatial aliasing, meaning multiple points on the sky will be candidate locations for a detected galaxy. This poster will explain the severity of this issue and ways to mitigate it through scan strategy. Another problem is that CHORD's spatial resolution is 3-5 arcminutes, which means the vast majority of the detected galaxies will provide no spatial information. The useful measurements such as redshift, total integrated HI mass, and velocity curve shape come from spectra. CHORD's base spectral resolution is 41\,km/s at the 21\,cm line rest frequency. To improve this, we can use a signal processing technique called up-channelization. This can help spectral resolution by a small power of 2, but introduces a new problem of frequency aliasing, where power is moved around to the wrong channels. This problem can also be solved with correct data processing. This poster will introduce the catalogue-producing pipeline, which is a matched filter with extra considerations to account for CHORD's unique physical and software design.

\hypertarget{poster-abstract-talk-212}{}\label{poster-abstract-page-talk-212}
\subsection*{Poster 18. Unlocking the evolution of star formation with CASTOR and Euclid}
\textbf{Presenter:} Gareth Jones\\
\textbf{Poster number:} 18\\
\textbf{Fields:} Galaxy $\cdot$ Next generation observatories

The Canadian-led Cosmological Advanced Survey Telescope for Optical and ultraviolet Research (CASTOR) mission will deliver deep, high-resolution ultraviolet (UV) imaging and complementary spectroscopy, opening a new window for studying star formation in galaxies out to intermediate-redshifts. CASTOR will enable transformative studies of galaxy growth, feedback, and environmental dependence. Rest-frame UV flux is a primary tracer of recent star formation, yet star formation rates (SFRs) derived from broadband photometry remain strongly model dependent, particularly with respect to dust attenuation, star formation histories, and stellar population prescriptions. We present preliminary results in preparation for CASTOR's Phase A+ study quantifying the reproducibility of photometrically derived SFRs in the CASTOR era. Using forward-modelled galaxy spectral energy distributions constructed with the Binary Population and Spectral Synthesis (BPASS) models, we evaluate the robustness of SFR estimates derived using CASTOR photometry, quantify systematic differences arising from varying modelling assumptions, and highlight the synergy with expected Euclid surveys at full depth.

\hypertarget{poster-abstract-talk-130}{}\label{poster-abstract-page-talk-130}
\subsection*{Poster 19. SED fitting of Globular Clusters at Intermediate Redshift}
\textbf{Presenter:} Jinoo Kim\\
\textbf{Poster number:} 19\\
\textbf{Fields:} Galaxy $\cdot$ Stars and Stellar cluster formation

We present a multi-band photometric study of globular cluster candidates in the galaxy cluster Abell 2744 at redshift z=0.308, using deep JWST/NIRCam imaging from 0.7 to 4.4 micron. Rather than relying only on colors, we fit the full spectral energy distributions of 69 bright candidates with simple stellar population models to estimate their masses, ages, and metallicities. The inferred masses are typically around 10$^7$ solar mass, placing these objects at the bright end of the globular cluster population. The best-fit solutions suggest a mix of populations, including older metal-poor and older metal-rich candidates, as well as a concentration at intermediate ages and moderate metallicity. As expected for broadband SED fitting, the results are affected by age-metallicity and mass-age degeneracies, which limit the precision of individual parameter estimates. Even so, the overall trends remain stable under alternative fitting choices. This work shows that JWST can begin to constrain the stellar populations of globular clusters at lookback times of several gigayears, while also highlighting the need for additional data to break remaining degeneracies.

\hypertarget{poster-abstract-talk-271}{}\label{poster-abstract-page-talk-271}
\subsection*{Poster 20. Forward modeling the starburst ring and YMCs of NGC 3351 in visibility plane}
\textbf{Presenter:} Peter Knowlton\\
\textbf{Poster number:} 20\\
\textbf{Field:} Galaxy

We investigate the morphology of the central starburst ring in NGC 3351, a nearby Milky Way analog galaxy, with archival high-resolution ALMA (sub-)millimeter observations. We focus on the young massive star clusters embedded within the disk, previously identified as continuum sources from the same dataset. Massive star clusters form in dense and high-pressure environments, and they seem able to efficiently convert gas into stars before the significant onset of stellar feedback from the newly formed stars. Multiwavelength observations trace the evolutionary stages of a massive clusters lifetime, with our ALMA observations being sensitive to the short infant phases, when the clusters are still deeply embedded within gas and dust. By forward modelling and fitting to our complex interferometric data in the visibility plane, we have access to the full resolution of our data, which can be lost when reconstructing the data in the image plane. We do this by utilizing the GALARIO library to generate models and sample the visibilities at the same baselines as our observations. This method should give an accurate description of the properties of the central ring and young massive clusters. Comparing these to shorter wavelength observations from HST and JWST, which trace the more evolved clusters, allows us to connect younger and older evolutionary stages. From these results, we can construct an evolutionary timeline for massive cluster formation and better constrain models of star formation and stellar feedback.

\hypertarget{poster-abstract-talk-63}{}\label{poster-abstract-page-talk-63}
\subsection*{Poster 21. The Balloon-borne VLBI Experiment: Instrument Overview and Upgrades}
\textbf{Presenter:} Mayukh Bagchi\\
\textbf{Poster number:} 21\\
\textbf{Field:} Instrumentation

In this talk, I will present BVEX, the Balloon-borne VLBI Experiment, where we are working to demonstrate the first-ever balloon to ground baselines in VLBI. In 2019, the Event Horizon Telescope (EHT) released the first image of the supermassive black hole M87* using Very Long Baseline Interferometry (VLBI). Ground-based VLBI is limited in resolution by the diameter of the Earth and in dynamic range by the number of suitable high frequency sites. A balloon-borne VLBI station operating above 99\% of the atmosphere at altitudes exceeding 35 km gives access to higher observing frequencies and better uv-coverage. BVEX is a 22 GHz, 0.9 m radio telescope built to demonstrate VLBI between a balloon-borne and ground-based telescope. The payload has a weld-free elevation mount, high-precision position tracking sensors, an RFSoC-based digital backend and data storage system, and OCXO timing with a custom drift monitoring system. While the first launch from Timmins in August 2025 was unsuccessful due to a balloon leak, we are preparing a reflight from Palmas, Brazil, in 2027. I will give a brief overview of the 2025 flight and then focus on three upgrades for 2027: a nitrogen-cooled cryogenic frontend at 22\,GHz to improve receiver sensitivity and calibration, a redesigned backend built around a mini-ITX board with 24 TB of NVMe storage, 100 GbE fibre optic interface, and CPU-bypass data streaming that brings power consumption down to approximately 100 W, and enhanced GPS data collection in RINEX format to supplement the OCXO timing for Precise Point Positioning post flight. I will wrap up by discussing mm-BVEX, a future W-band (88 GHz) telescope designed to observe alongside the GMVA with frequency phase transfer between 22 and 88 GHz.

\hypertarget{poster-abstract-talk-47}{}\label{poster-abstract-page-talk-47}
\subsection*{Poster 22. In-band spectral index techniques for extended radio sources}
\textbf{Presenter:} Brianna Ball\\
\textbf{Poster number:} 22\\
\textbf{Field:} Instrumentation

Spectral indices are an effective way of differentiating between thermal and non-thermal sources of radio emission. Most often, spectral indices are determined using flux measurements taken from multiple observations with different instruments. Since this data is not always available, finding indices from a single telescope observation is desirable. The reliability of these in-band spectral indices has not been well studied, particularly for extended sources. Here, I will demonstrate three different techniques for finding in-band spectral indices using EMU (Evolutionary Map of the Universe) data from ASKAP (the Australian Square Kilometre Array Pathfinder). These techniques include: integrating fluxes in sub-bands, temperature-temperature plots, and utilizing the Taylor terms produced in the imaging process. I apply these techniques to a selection of extended sources of radio emission, primarily supernova remnants (SNRs) and HII regions, and assess if the indices produced by the different techniques agree with the expected values. I examine how well the in-band techniques can differentiate between the thermal (HII regions) and non-thermal (SNRs) sources. In general, I find the T-T plot method to be the most reliable across a range of source sizes and brightnesses for both SNRs and HII regions.

\hypertarget{poster-abstract-talk-241}{}\label{poster-abstract-page-talk-219}
\subsection*{Poster 23. An Exposure Time Calculator for the CASTOR mission's UV Multi-Object Spectrograph}
\textbf{Presenter:} Jade Yeung\\
\textbf{Poster number:} 23\\
\textbf{Field:} Instrumentation

The Cosmological Advanced Survey Telescope for Optical and UV Research (CASTOR) is a proposed 1 meter space telescope under development by the National Research Council of Canada, the Canadian Space Agency, and partner institutions around the world. It will serve as a successor to the Hubble Space Telescope (HST) at short wavelengths, providing continued access to the ultra-violet (UV) and blue-optical at high spatial resolution ($\sim$0.15) over a field of view $\sim$100x larger than HST. The UV Multi-Object Spectrograph (UVMOS) is a medium resolution (R$\sim$1500) spectrograph based on the BATMAN spectrometer that is proposed as a parallel instrument onboard CASTOR. Here we present a pixel-based, user-friendly, extensible exposure time calculator (ETC) for the CASTOR UVMOS instrument, adaptable for other space-borne MOS missions. I will be presenting the current efforts outlining various key science cases, including M dwarf flares, white dwarf atmospheres, and next generation galactic surveys in the ultraviolet.

\hypertarget{poster-abstract-talk-105}{}\label{poster-abstract-page-talk-105}
\subsection*{Poster 24. Latest APERO news}
\textbf{Presenter:} Neil Cook\\
\textbf{Poster number:} 24\\
\textbf{Field:} Instrumentation

APERO (A Pipeline to REduce Observations) continues to evolve as the official data-reduction pipeline for SPIRou and as an additional reduction pipeline for NIRPS. We present an update on recent processing improvements and new features, most notably those enabling sub-m/s radial-velocity precision. Additionally, we highlight upcoming developments, including support for new instruments such as SPIP on the TBL and iLocater on the LBT. We also outline ongoing efforts to strengthen APEROs long-term computational and data-distribution ecosystem. Our team is currently working toward an Alliance allocation to accelerate data processing and, importantly, to enable timely data sharing across the Canadian community. In parallel, we are collaborating with the CADC to establish a dedicated collection for APERO and LBL data products, ensuring robust archival access and discoverability. Finally, we are exploring how CANFAR virtual machines can be leveraged to sustain and scale these capabilities, particularly in preparation for the upcoming CFHT community survey, where stable, community-wide processing infrastructure will be essential.

\hypertarget{poster-abstract-talk-146}{}\label{poster-abstract-page-talk-146}
\subsection*{Poster 25. The CHIME/FRB Heartbeat Service: Monitoring Uptime and Exposure of the CHIME/FRB+Outrigger Network}
\textbf{Presenter:} Maxwell Fine\\
\textbf{Poster number:} 25\\
\textbf{Fields:} Instrumentation $\cdot$ Transients

The CHIME/FRB experiment has discovered thousands of fast radio bursts (FRBs), providing an unprecedented sample for population studies. To enable precise interferometric localization, CHIME/FRB operates a network of three outrigger telescopes that record data following real-time triggers from the main instrument. Our primary scientific objective is to measure the effective exposure of the CHIME/FRB+outrigger system to FRB sources, enabling robust FRB rate calculations and population synthesis studies. Prior to this work, no systematic framework existed to quantify outrigger uptime or the joint exposure of the CHIME/FRB+outrigger network.

We developed the CHIME/FRB heartbeat service, a monitoring system that measures the operational state of the CHIME/FRB+outriggers system. The service injects synthetic triggers, called heartbeats, into the same end-to-end data acquisition pipeline used for real FRB events. These triggers prompt CHIME/FRB and the outrigger stations to record short segments of data, producing a record of whether each station was capable of writing data at any given time.

The Heartbeat service provides measurements of outrigger uptime and the effective exposure of the combined system, enabling robust FRB rate estimates and population synthesis studies. In addition, the heartbeat data reveal intermittent failures and infrastructure issues that would otherwise remain undetected, allowing these problems to be identified and addressed and thereby improving the reliability and scientific return of the CHIME/FRB+outrigger network.

\hypertarget{poster-abstract-talk-148}{}\label{poster-abstract-page-talk-148}
\subsection*{Poster 26. Exploring the trade-space of distributed aperture telescopes in spectroscopy}
\textbf{Presenter:} Theodore Grosson\\
\textbf{Poster number:} 26\\
\textbf{Field:} Instrumentation

Scientific programs targeting the faintest objects push modern telescopes to increasingly large sizes, simultaneously pushing their instrumentation to increasingly high costs and design complexities. Obtaining spectra of low surface brightness galaxies, for example, requires hours of observations using state-of-the art instruments on the largest telescopes. An innovative strategy for combating these costs is to use many small telescopes in place of a single large aperture. For observations in which high angular resolution is not necessary, this can be equivalent to a large diameter telescope in terms of sensitivity, so long as the noise associated with stacking many images can be minimized. Improvements in commercial off-the-shelf (COTS) components, especially low noise detectors, have made this concept a possibility, as demonstrated by instruments such as the Dragonfly Telephoto Array, the Huntsman Telescope, and the Condor Array Telescope. In this work, we discuss the applicability of so-called "distributed aperture telescopes" to spectroscopic observations of ultra-diffuse galaxies. Observations of a large sample of these poorly-understood objects are necessary to constrain their properties and formation pathways. We compare the cost and simulated performance of different configurations of apertures, detectors, and other components in observing these galaxies, finding where cost is minimized while performance is maximized. We proceed to describe a telescope that would be able to observe spectra of low surface brightness galaxies at a much lower cost than major facilities, and discuss the path to prototyping this concept.

\hypertarget{poster-abstract-talk-263}{}\label{poster-abstract-page-talk-263}
\subsection*{Poster 27. Programmable Slits for the CASTOR Ultraviolet Multi-Object Spectrograph to Unlock the History of Multiple Stellar Populations in Globular Clusters}
\textbf{Presenter:} Charles Lee\\
\textbf{Poster number:} 27\\
\textbf{Fields:} Instrumentation $\cdot$ Stars and Stellar cluster formation

The Ultraviolet Multi-Object Spectrograph (UVMOS) is a proposed UV (150 to 300 nm) instrument for the CASTOR mission. The UVMOS utilizes a digital micromirror device (DMD) to program object-selecting slits. The DMD consists of programmable mirrors that tilt $\pm$ 12 degrees to selectively project the incoming field to a collecting aperture, allowing the UVMOS to filter image fields spatially.

The DMD introduces additional light propagation effects from scatter and diffraction. We previously developed an electromagnetic simulation to understand these effects. We introduce new results derived from improvements to our simulation model, quantifying the DMD efficiency, contrast, and stray light in the UVMOS operating wavelengths (150 to 300 nm). As an illustration of CASTOR UVMOSs planned capabilities, we incorporate these parameters inside the FORECASTOR exposure time calculator and build a simulation model for observations. This simulator is used to draft a proposed survey of multiple stellar populations in globular clusters to be carried out by the CASTOR mission upon launch in the next decade.

\hypertarget{poster-abstract-talk-50}{}\label{poster-abstract-page-talk-50}
\subsection*{Poster 28. Inferring cluster assembly from present day star cluster dynamics}
\textbf{Presenter:} Isabella Armstrong\\
\textbf{Poster number:} 28\\
\textbf{Field:} Stars and Stellar cluster formation

Individual stars are born from gas clouds. Star clusters, on the other hand, can be assembled over time. Developing methods to probe these different timescales is crucial for constraining the formation and evolution of nearby star clusters.

Observations provide a single snapshot of cluster properties. Cluster ages are typically derived using stellar evolution models (e.g. isochrone fitting), while 6D kinematics can be used to compute the trajectories of the stars (e.g. Gaia astrometry). Recent works construct dynamical ages from the kinematics to determine when the cluster was formed. However, different methods use different properties of the stellar motion and different assumptions of the cluster assembly process. Because methods are derived from a single set of cluster properties, it is unknown how accurately they trace cluster formation and assembly.

In this work, we compare different dynamical ages for a simulated cluster. We find that some methods evolve similarly to isochronal ages whereas others better trace the expansion of the cluster. Methods which trace the size evolution of a cluster evolve differently from methods tracing the origin of individual stars, producing different dynamical age estimates. The final dynamical age depends on assumptions about the initial cluster geometry, stellar velocities, and boundness. The differences between different dynamical ages and the isochronal age can be used to investigate the early kinematics of star clusters. Star cluster assembly timescales have direct implications for understanding how dense environments affect binary and planetary properties. Our work shows that the concept of a dynamical age is highly method dependent, with different methods being appropriate for tracing different formation timescales.

\hypertarget{poster-abstract-talk-190}{}\label{poster-abstract-page-talk-190}
\subsection*{Poster 29. Rapid Rotation in Magnetic A/B Stars: New Constraints from Photometry and Spectropolarimetry}
\textbf{Presenter:} Noah Budz\\
\textbf{Poster number:} 29\\
\textbf{Field:} Stars and Stellar cluster formation

Magnetic chemically peculiar (mCP) A and B stars host stable, globally organized magnetic fields that are thought to be fossil remnants from earlier evolutionary stages. These stars are typically slow rotators due to magnetic braking early in their lifetimes, making the discovery of rapidly rotating magnetic stars particularly intriguing. The existence of a small population of rapidly rotating magnetic stars, therefore presents an important challenge for understanding the rotational evolution of stars hosting fossil magnetic fields.

We investigate a sample of five candidate fast-rotating A/B stars using a combination of high-resolution spectropolarimetry and space-based photometry. Longitudinal magnetic field measurements derived from spectropolarimetric observations are used to search for magnetic signatures and characterize magnetic variability. In parallel, TESS photometry is analyzed to measure rotational modulation associated with surface abundance structures. Light curves are extracted and detrended before applying Lomb-Scargle periodogram analysis to identify periodic variability and refine stellar rotation periods.

We report magnetic detections in two stars within the sample and present updated rotational constraints for all targets based on the combined spectropolarimetric and photometric analysis. These results provide new observational constraints on the rotational properties of rapidly rotating magnetic intermediate-mass stars and contribute to understanding the origin of this rare population of renegade magnetic A/B stars.

\hypertarget{poster-abstract-talk-302}{}\label{poster-abstract-page-talk-302}
\subsection*{Poster 30. Evidence of Dust Grain Growth in Protoplanetary Disk IRS 63 Found Using Polarized Light}
\textbf{Presenter:} Audrey Burggraf\\
\textbf{Poster number:} 30\\
\textbf{Field:} Stars and Stellar cluster formation

To understand how planets form, we must study how dust grains in protoplanetary disks stick together and grow in size.

One of the biggest challenges in understanding this process is determining the sizes of dust grains themselves. Fortunately, polarized light from self-scattering in disks provides us with a rare opportunity to constrain dust grain sizes. Self-scattering occurs when thermally emitted light from dust is scattered by other grains in the disk, with the resulting polarization pattern depending on dust size, disk geometry, and optical depth.

This project analyzes multi-wavelength ALMA observations (0.87 mm, 1.3 mm, 1.7 mm, and 2.0 mm) to study polarized emission in the Class I protostellar disk IRS 63.

We detect wavelength-dependent changes in both polarization morphology and polarization fraction. By looking at slices of polarized intensity, we find evidence that the disk is geometrically thick and flared. Focusing on the central region of the disk and assuming spherical dust grains, our modelling finds the grains are compact with a constrained maximum dust grain size of approximately a few hundred microns. These grains are much larger than the typical interstellar medium dust grains, providing evidence that dust grain growth has occurred in the disk. To explore the effects of non-spherical dust grains, we utilize a neural network to model the polarized emission from grains to place improved constraints on grain shapes and sizes in the centre of the disk.

Constraining grain sizes in young disks is essential for understanding the early stages of planet formation and how micron-sized dust grains evolve into planets.

\hypertarget{poster-abstract-talk-220}{}\label{poster-abstract-page-talk-220}
\subsection*{Poster 31. All that (seemingly) flares is not magnetic: the cautionary tale of HD 149834}
\textbf{Presenter:} Alexandre David-Uraz\\
\textbf{Poster number:} 31\\
\textbf{Field:} Stars and Stellar cluster formation

Recently, the massive pulsating B-type primary of the eclipsing binary system HD 149834 was claimed to be magnetic based on two main indirect pieces of evidence. First, a low-amplitude, but significant, peak in the periodogram of its TESS (Transiting Exoplanet Survey Satellite) light curve was interpreted as being due to rotational modulation. Second, an X-ray flare was observed by the Chandra X-ray Observatory, at an orbital phase deemed to be incompatible with it originating from the low-mass companion. In this work, we present new spectropolarimetric observations of HD 149834, which do not allow us to detect any magnetic field. We also revisit the phasing and energetics of the observed X-ray flare to reassess its potential origin. Finally, we discuss the broader context of this claim, and why flares from magnetic massive stars are not expected in the first place.

\hypertarget{poster-abstract-talk-331}{}\label{poster-abstract-page-talk-331}
\subsection*{Poster 32. Investigating the Influence of Magnetic Fields on Bubble Morphology in Star-Forming Regions}
\textbf{Presenter:} Raina Irons\\
\textbf{Poster number:} 32\\
\textbf{Field:} Stars and Stellar cluster formation

Magnetic fields play a fundamental role in the process of star formation, influencing the structure and evolution of molecular clouds. These fields can become warped and shaped into shell-like structures from the energy released by massive stars, referred to as bubbles. This process leads to higher densities along the edges of bubbles, and can influence the morphology of surrounding magnetic fields to align tangentially as it expands. Despite their importance, the precise role of magnetic fields in bubble evolution remains an open question, and a comprehensive statistical analysis of a large sample of bubbles is necessary in order to work towards understanding this question. I will introduce a novel technique developed to search for and study the impact of bubbles on magnetic field lines that detect circular structures in dust polarization maps will be presented. Thermal emission from dust grains exhibit a linear polarization due to grain alignment with magnetic fields, and can be used to infer the magnetic field morphology through their polarization. This method uses ideal circular models, called kernels, and compares the kernel to the polarization map in a pixel-by-pixel comparison. Circular alignment maps are then produced which signal where the maps exhibit either circular and radial features. A sample of molecular clouds from the BISTRO survey are studied, and I will present results of candidate bubbles for each of these clouds, as well as confirmed bubbles from the WISE and Simpson catalogs. In doing so, we determine broader trends in how bubbles influence their environment and establish a clearer relationship between polarization efficiency and bubble properties.

\hypertarget{poster-abstract-talk-143}{}\label{poster-abstract-page-talk-143}
\subsection*{Poster 33. AGN Contamination Biases in Resolved Galaxy Scaling Relations}
\textbf{Presenter:} Aayushmathi Ashok\\
\textbf{Poster number:} 33\\
\textbf{Fields:} Astrostatistics, Astroinformatics, and AI/Machine Learning $\cdot$ Galaxy $\cdot$ Gas, dust and the ISM

Galaxy scaling relations are empirical connections between physical properties such as stellar mass, gas-phase metallicity, and star formation rate, and are fundamental tools for understanding how galaxies grow and evolve. These relations are built from nebular emission-line diagnostics of the interstellar medium, which are only physically meaningful in star-forming regions. However, Active Galactic Nuclei (AGN) and Low-Ionization Emission-line Regions (LIERs) produce many of the same emission lines as star-forming regions but with different relative intensities. Even in spatially resolved surveys, where individual spatial pixels (spaxels) capture emission from distinct regions within a galaxy, the signal from star-forming and non-star-forming processes can overlap. Standard selection criteria, such as diagnostic diagrams and equivalent width cuts, attempt to isolate star-forming spaxels, but both discard real star-forming emission in excluded regions and fail to fully remove contamination in spaxels that pass the cuts, thereby introducing possible systematic biases into scaling relations. We apply a spectral decomposition method built on a deep neural network to separate each spaxels spectrum into star-forming and non-star- forming components. By fitting the decomposed star-forming spectra, we recover emission-line fluxes and physical parameters even in composite and AGN-dominated regions, where the star- forming contribution may be small but remains present. Using these recovered measurements, we construct key-resolved scaling relations from the isolated star-forming component, directly assessing whether non-star-forming contamination affects traditional scaling relations and how that effect depends on galaxy and spectral properties

\hypertarget{poster-abstract-talk-162}{}\label{poster-abstract-page-talk-162}
\subsection*{Poster 34. Modelling the CMB Optical Depth from Reionization Physics: A Symbolic Regression Approach}
\textbf{Presenter:} Marek Detière-Venkatesh\\
\textbf{Poster number:} 34\\
\textbf{Fields:} Astrostatistics, Astroinformatics, and AI/Machine Learning $\cdot$ Cosmology and Early Universe

The CMB optical depth, $\tau$, is a key cosmological parameter that encodes the integrated ionization history of the Universe. In current CMB analyses, $\tau$ is typically treated as a free parameter and fitted jointly with other cosmological quantities. In principle, however, $\tau$ should be predictable from the astrophysical processes that govern reionization. Developing accurate, physically motivated models for $\tau$ could therefore enable its independent inference and help reduce parameter degeneracies in CMB analyses.

In this work, we use three cosmological simulations to study how $\tau$ depends on the astrophysical parameters that regulate reionization. Because these simulations exhibit highly nonlinear parameter dependencies, we employ symbolic regression (SR), a machine-learning approach that produces interpretable, closed-form analytic expressions. Unlike traditional black-box ML methods, the resulting equations make the parameter dependence explicit while still achieving percent-level accuracy. In addition, equations derived for different simulations provide a transparent way to compare alternative reionization modelling frameworks on equal footing.

Our results reveal notable differences among simulations in their predicted dependence of $\tau$ on cosmological parameters, particularly the scalar spectral index and the matter fluctuation amplitude. These compact SR relations also provide a practical route for propagating reionization-informed estimates of $\tau$ into cosmological inference, forecasting improvements of up to 20\% in parameter constraints. As one of the first applications of SR to reionization, this work shows that complex simulation outputs can be distilled into compact equations, and it establishes a general framework for comparing simulations through shared observables.

\hypertarget{poster-abstract-talk-34}{}\label{poster-abstract-page-talk-34}
\subsection*{Poster 35. Kinematics of HI-rich Ultra-diffuse Galaxies}
\textbf{Presenter:} Lawrence Faria\\
\textbf{Poster number:} 35\\
\textbf{Fields:} Astrostatistics, Astroinformatics, and AI/Machine Learning $\cdot$ Galaxy

Ultra-Diffuse Galaxies (UDGs) are an extreme population of low surface brightness galaxies that challenge our understanding of galaxy formation and evolution. Despite their faint stellar component, many are rich in neutral hydrogen (HI), making them ideal targets for kinematic studies. However, the limited spatial resolution, low rotation speeds, and uncertain inclinations typical of UDGs place these systems in a regime where traditional tilted-ring kinematic models are poorly constrained, complicating efforts to determine their dynamical properties.

Using new VLA HI observations of gas-rich UDGs together with archival data, we are compiling a large sample of resolved HI maps and examining their velocity structure. Preliminary tilted-ring modelling allows an initial placement of these systems on the baryonic TullyFisher relation, while highlighting the limitations of conventional methods in this regime. To address these challenges, we are developing diffusion-based inference models trained on synthetic HI data cubes to infer kinematic parameters and associated uncertainties. I will present the VLA observations, preliminary kinematic results, and ongoing efforts to develop more robust methods for interpreting the HI kinematics of ultra-diffuse galaxies.

\hypertarget{poster-abstract-talk-294}{}\label{poster-abstract-page-talk-294}
\subsection*{Poster 36. Revealing Hidden Structure in Large Spectroscopic Datasets with Machine Learning: From SDSS to DESI}
\textbf{Presenter:} Stéphanie Juneau\\
\textbf{Poster number:} 36\\
\textbf{Fields:} Astrostatistics, Astroinformatics, and AI/Machine Learning $\cdot$ Galaxy $\cdot$ Supermassive black hole

Understanding how galaxies evolve requires studying both their stellar populations and the growth of their central black holes. Spectroscopic surveys provide a powerful window into this co-evolution, encoding information about star formation, chemical enrichment, gas kinematics, and nuclear activity within galaxy spectra. However, the complexity of spectral data means that traditional analyses often rely on a limited set of diagnostic features, potentially missing broader patterns embedded across the full spectrum. I will present recent work applying machine learning (ML) techniques to large spectroscopic datasets to uncover previously missed trends. Using $>$200,000 galaxy spectra from the Sloan Digital Sky Survey, we applied an Autoencoder and dimensionality-reduction methods to construct low-dimensional representations of spectra. We found coherent trends and groupings that track galaxy quenching but also revealed surprising morphological sequences typically associated with imaging data, highlighting how ML/AI can expose relationships beyond traditional spectral diagnostics. I will then discuss how this approach is being expanded to the much larger Dark Energy Spectroscopic Instrument (DESI) survey, with spectra for more than 35 million stars, galaxies and quasars. Beyond scaling analyses to this unprecedented sample, we are also exploring new ML frameworks to capture the full complexity of astronomical data. In particular, I will introduce a new multimodal astronomical foundation model, the Astro Vision model for Images and Spectra (AVIS), trained jointly on pairs of images and spectra spanning stars, galaxies and quasars. By learning shared representations across images and spectra, AVIS enables new forms of discovery and interpretation in large astronomical surveys.

\hypertarget{poster-abstract-talk-98}{}\label{poster-abstract-page-talk-98}
\subsection*{Poster 37. Probing foreground galaxy haloes in the nearby Universe using CHIME/FRB Baseband Catalog 2}
\textbf{Presenter:} Affan Khadir\\
\textbf{Poster number:} 37\\
\textbf{Fields:} Astrostatistics, Astroinformatics, and AI/Machine Learning $\cdot$ Galaxy $\cdot$ Transients

The circumgalactic medium (CGM) is a reservoir of gas around galaxies that can condense onto galactic halos and disks, and drive star formation. In quenched galaxies, the CGM is heated by AGN and supernovae feedback, as well as the hot intracluster or intragroup medium, suppressing active star formation. Therefore, understanding the distribution of matter and the baryonic content of the CGM is important to obtain a full picture of star formation, galaxy evolution, and locating the Universes missing baryons. Since most of the gas in the CGM is composed of ionized baryons, fast radio burst (FRB) dispersion measures (DMs) are ideal to probe the CGM. Using $\sim$2100 FRBs (with positional accuracies of $\sim$10') from the upcoming CHIME/FRB Baseband Catalog 2, we present the largest statistical study to date of the CGMs baryon content based on FRB DMs, obtained by stacking sightlines as a function of foreground galaxy impact parameter. We explore the possible presence of a negative correlation between DM and impact parameter of foreground galaxies, and use this DM profile to constrain models of the average electron density for galaxy haloes in our sample. We also conduct a closer analysis of large nearby galaxies like M31 and M33.

\hypertarget{poster-abstract-talk-159}{}\label{poster-abstract-page-talk-159}
\subsection*{Poster 38. Tracing Star Formation Histories and Evolution in $\rho$ Ophiuchi}
\textbf{Presenter:} Sharone Elsa Abhilash\\
\textbf{Poster number:} 38\\
\textbf{Fields:} Gas, dust and the ISM $\cdot$ Stars and Stellar cluster formation

The $\rho$ Oph complex is one of the nearest and best-studied sites of ongoing star formation, hosting a rich population of young stellar objects (YSOs) and represents an archetypal region for studying the earliest stages of stellar and molecular cloud evolution. Part of the Upper Scorpius (USco) subgroup of the Sco-Cen association, the region lies near several feedback sources that shape the surrounding cloud. Stellar feedback can drive sequential star formation across molecular clouds, producing clusters that trace the propagation of feedback through gas. As the youngest region within the USco subgroup, $\rho$ Oph allows us to study early imprints of feedback-driven star formation. We therefore construct a multiwavelength census of pre and protostellar structures in $\rho$ Oph by combining Gaia YSOs identified by the SPYGLASS survey with infrared protostellar catalogs, dense core surveys, Herbig-Haro objects, and 13CO molecular gas, together with kinematic data including IR proper motions and multiple RV surveys. We recover known subregions using HDBSCAN clustering and estimate ages through isochrones, protostellar abundances and lifetimes, and dynamical traceback. Our results reveal internal age gradients across the complex, with the youngest populations concentrated in L1688 (Oph A-F), L1689, and adjacent structures such as L1709 and L1712. Older, unbound populations are found near the $\rho$ Oph star and among the broader USco population. Using dynamical traceback analyses, we are investigating signatures of expansion and possible feedback-driven acceleration. These analyses explore how stellar feedback drives gas dispersal and the dynamical evolution of young stars, and how the star formation history of $\rho$ Oph is linked to nearby feedback sources.

\hypertarget{poster-abstract-talk-89}{}\label{poster-abstract-page-talk-89}
\subsection*{Poster 39. Dynamics of Dense Gas Filaments in the Serpens South Star-Forming Region}
\textbf{Presenter:} Julian Caza\\
\textbf{Poster number:} 39\\
\textbf{Field:} Gas, dust and the ISM

Characterizing the dynamics of dense gas in hub-filament systems is crucial to understand the reservoir of mass from which young stellar objects can accrete. In this study, we investigate relative orientations between velocity gradients, filamentary structures, and gravitational forces in the Serpens South region both globally and by filament. Using NH$_3$ (1,1) emission from combined VLA+GBT observations as a dense gas tracer, we perform a kinematic analysis of velocity-coherent filaments. We then compare the orientations of local velocity gradients to the orientation of filamentary structures and to the plane-of-sky gravitational acceleration field derived from Herschel H$_2$ column density maps. We observe some filaments where velocity gradients are preferentially oriented perpendicular to their structure, which may suggest mass accretion onto filaments from their surroundings, however we also observe localized regions where gradients are aligned parallel to filaments which could be indicative of longitudinal mass flow. We also find areas where the velocity gradient field is oriented parallel to the inferred plane-of-sky gravitational field, suggesting that gas motion in these areas could be the result of gravitational infall. The significance of these relative orientation trends are quantified using the Histogram of Relative Orientations (HRO) method and the Projected Rayleigh Statistic (PRS). We look to identify trends in orientation based on global properties, such as column density and NH$_3$ (1,1) intensity, and filamentary properties, such as mass and distance to the central star-forming hub.

\hypertarget{poster-abstract-talk-239}{}\label{poster-abstract-page-talk-239}
\subsection*{Poster 40. Galactic Center at Sub-Parsec Scales: Studying Gas Ionization and Kinematics Around SgrA* with JWST}
\textbf{Presenter:} Nicole Ford\\
\textbf{Poster number:} 40\\
\textbf{Fields:} Gas, dust and the ISM $\cdot$ Supermassive black hole

SgrA* is the nearest quiescent massive black hole (BH), and its proximity offers a unique opportunity to study its surrounding fuel supply. Using JWST/MIRI observations, we peer through galactic dust to measure the spatial, kinematic, and intensity distribution of ionized gas within the central $\sim$0.1 parsec of SgrA*. We use forbidden lines to derive gas radial velocities, line ratios, and abundances for multiple distinct kinematic components that overlap along the line of sight. The three dominant emitting structures that we identify are: the minispirals Northern Arm, the Bar region, and high-velocity ($\geq$ 600 km/s) gas clumps likely produced by colliding WolfRayet stellar winds or SgrA* outflows. We infer the radiation field hardness and metallicity of each of these emitting gas structures. This study sheds light on the origins of the galactic fuel reservoir and how it is physically shaped and ionized by stellar as well as SgrA*'s activity. Our findings contextualize accretion and feedback processes in low luminosity AGN and quiescent SMBH systems where the central parsec is unresolvable.

\hypertarget{poster-abstract-talk-77}{}\label{poster-abstract-page-talk-77}
\subsection*{Poster 41. Unwinding the Magnetic Wind: A Polarimetric Study of the Pulsar Wind Nebula G295.6+0.5 with ASKAP}
\textbf{Presenter:} Gabrielle Graham\\
\textbf{Poster number:} 41\\
\textbf{Field:} Gas, dust and the ISM

Magnetic fields make a significant but poorly understood contribution to the energetics of supernova remnants. Directly probing magnetic fields is a notoriously difficult problem, but new wideband radio continuum images from SKA precursor facilities offer a promising new direction. Here we present an in-depth radio study of the magnetic field across the pulsar wind nebula G295.6+0.5. Polarization and continuum data are acquired with the Australian Square Kilometre Array Pathfinder (ASKAP) telescope as part of the Polarization Sky Survey of the Universes Magnetism (POSSUM) and Evolutionary Map of the Universe (EMU) surveys. A rotation measure (RM) map is produced numerically via Faraday rotation synthesis and Stokes Q and U fitting routines. The RM map, in conjunction with the POSSUM data, allows us to measure the intrinsic polarization angle and internal Faraday rotation of G295.6+0.5. For the first time, we extract the PWNs magnetic field configuration using these methods, providing a case study for future application.

\hypertarget{poster-abstract-talk-199}{}\label{poster-abstract-page-talk-199}
\subsection*{Poster 42. Tracing PAH populations with JWST using weak 5--6\,$\mu$m bands in the Orion Bar}
\textbf{Presenter:} Alice Guerras\\
\textbf{Poster number:} 42\\
\textbf{Field:} Gas, dust and the ISM

Polycyclic aromatic hydrocarbons (PAHs) produce a set of infrared emission bands known as aromatic infrared bands (AIBs) and are widespread in the interstellar medium. These bands are widely used as tracers of star formation and interstellar conditions in both Galactic regions and external galaxies. The unprecedented sensitivity and spatial resolution of the MIRI-MRS on the James Webb Space Telescope (JWST) now allow the Orion Bar photodissociation region (PDR) to be probed in exceptional detail. Here we investigate the weak and rarely studied AIBs at 5.25, 5.75, and 5.90\,$\mu$m across the key zones of this strongly UV-illuminated PDR to constrain their carriers and their evolution across the PDR.

The 5.25 and 5.75\,$\mu$m AIBs show a spatial distribution like that of the main AIBs, with emission peaking in the atomic PDR. In contrast, the 5.90\,$\mu$m AIB exhibits a distinct spatial behavior intermediate between the main AIBs and bands associated with aliphatic material or PAH clusters and very small grains (VSGs). Spectrally, the 5.90\,$\mu$m AIB has a symmetric profile, whereas the 5.75 and 5.25\,$\mu$m AIBs display more complex structures. In particular, the 5.75\,$\mu$m AIB can be decomposed into three components whose relative intensities vary across the Orion Bar.

Correlation analyses and principal component analyses indicate that these components trace different PAH populations: one is associated with fragile PAHs and VSGs in shielded molecular zones, while the others are linked to stable neutral and cationic PAHs in more strongly UV-irradiated regions. These results highlight the diagnostic potential of the weak 5--6\,$\mu$m AIBs for probing the evolution of PAH populations across photodissociation regions and for interpreting JWST observations of star-forming environments in nearby galaxies.

\hypertarget{poster-abstract-talk-207}{}\label{poster-abstract-page-talk-207}
\subsection*{Poster 43. Disentangling Multiwavelength Emission from X-ray Binaries through Colour-Magnitude Diagram Modelling}
\textbf{Presenter:} Emily Carver\\
\textbf{Poster number:} 43\\
\textbf{Fields:} Compact Objects $\cdot$ Transients

Black hole X-ray binaries in our Galaxy are interacting systems comprised of a stellar-mass black hole and a donor star that accretes matter onto the black hole. These systems undergo outburst cycles wherein they brighten across the electromagnetic spectrum, and relativistic jet outflows are launched. The emission patterns of X-ray binaries in outburst are complex, with contributions from multiple system components and multiple radiation mechanisms. In particular, the optical/infrared emission may be composed of thermal components from the accretion disc and donor star, X-ray reprocessing in the disc, or non-thermal synchrotron emission from relativistic jet outflows. In this talk, we present a new technique for disentangling these complex emission patterns through building and modelling colour-magnitude diagrams, using two of the largest community-driven archives of X-ray binary outburst observations available: the Las Cumbres Observatory Network (optical/infrared) and the Neil Gehrels Swift Observatory (X-ray). This unique toolkit provides a new avenue to understand how changing optical/infrared emission properties translate to changing physical properties of the accretion disc and jet components in outbursting X-ray binaries.

\hypertarget{poster-abstract-talk-277}{}\label{poster-abstract-page-talk-277}
\subsection*{Poster 44. Investigating the spectral volatility of RRAT J1541+47 with CHIME/Pulsar}
\textbf{Presenter:} Alyssa Cassity\\
\textbf{Poster number:} 44\\
\textbf{Fields:} Compact Objects $\cdot$ Transients

The Canadian Hydrogen Intensity Mapping Experiment (CHIME) conducts a daily survey of the entire northern sky, enabled by its cylindrical design and large field of view. CHIME is equipped with a novel single-pulse survey that conducts a rapid, all-sky daily search for pulsars using the CHIME/Fast Radio Burst (FRB) backend. As the largest survey of its kind, it probes pulsar parameter space largely inaccessible to traditional periodicity searches, enabling the discovery of highly intermittent Galactic transients. In particular, the CHIME/FRB single-pulse survey has discovered over 100 new pulsars and Rotating Radio Transients (RRATs). RRATs are a type of neutron star that emit bright, sporadic radio bursts, and are therefore only detectable via their individual pulses.

In this talk, I will present a study of J1541+47, a unique RRAT discovered by the CHIME/FRB single-pulse survey in 2020. Using over 200 hours of observations with CHIME/Pulsar, we detected $\sim$1000 single pulses. I will present a timing solution and characterize the interesting single-pulse properties of the source, including double-peaked profiles and extremely narrow pulse widths ($<$0.025\% of the pulse period). Most notably, I will show that it undergoes extreme spectral variability on timescales of daysweeks. I will discuss the implications of these properties for neutron star emission physics, and explore connections between RRATs and other radio transient phenomena.

\hypertarget{poster-abstract-talk-169}{}\label{poster-abstract-page-talk-169}
\subsection*{Poster 45. The discovery of a hyper-luminous rotating radio transient with CHIME/FRB}
\textbf{Presenter:} Fengqiu Adam Dong\\
\textbf{Poster number:} 45\\
\textbf{Fields:} Compact Objects $\cdot$ Transients

The CHIME/FRB instrument has revolutionized the landscape of ms-duration radio transients in the northern hemisphere. Along with extragalactic Fast Radio Bursts (FRB), CHIME/FRB has discovered 111 new Galactic radio transients. These are primarily rotating radio transients (RRAT), a subtype of pulsars that share many commonalities with FRBs. However, a problem remains: the luminosities of RRATs are orders of magnitude fainter than those of FRBs. In this talk, I will review the CHIME/FRB Galactic transient survey. I will pay particular attention to a new RRAT that we have discovered, PSR J2109+50. This source has been detected in the far sidelobes of CHIME, implying kJy peak flux densities. We find that, for the first time, we are forming a bridge between the energetics of the RRAT and FRB populations. I will subsequently discuss the future of our survey and its potential in the era of CHIME/FRB outriggers.

\hypertarget{poster-abstract-talk-150}{}\label{poster-abstract-page-talk-150}
\subsection*{Poster 46. Orbital Modelling of Fast-Moving Stars in Omega Centauri: New Constraints on the Central Intermediate-Mass Black Hole}
\textbf{Presenter:} Mackenzie Hayduk\\
\textbf{Poster number:} 46\\
\textbf{Fields:} Compact Objects $\cdot$ Stars and Stellar cluster formation

Intermediate-mass black holes (IMBHs), with masses between $10^2$ and $10^5$ M$_\odot$, represent a missing link between stellar and supermassive black holes. Characterizing their population is key to understanding the evolution of galactic centres and interpreting future gravitational wave signals detectable with the upcoming Laser Interferometer Space Antenna (LISA). The Milky Way globular cluster Omega Centauri ($\omega$ Cen) is a prime candidate for hosting an intermediate-mass black hole. Recently, Haberle et al. (2024) identified several stars near the cluster core that exceed the central escape velocity, constraining the IMBH mass to be $>$8,200 M$_\odot$. In this work, we analyse the central mass by fitting the full orbits of the fast-moving stars using the Octofitter package and all available astrometric data. We obtain new constraints on the mass and position of the IMBH and show that the fast stars admit bound Keplerian solutions around a common mass. We predict the line-of-sight velocities and future orbits of the fast-moving stars, which will help to refine this analysis when compared with upcoming James Webb Space Telescope data.

\hypertarget{poster-abstract-talk-96}{}\label{poster-abstract-page-talk-96}
\subsection*{Poster 47. Sharing the Sky: Connecting Canadian FRB Discoveries to the Global Transient Community}
\textbf{Presenter:} Daniel Amouyal\\
\textbf{Poster number:} 47\\
\textbf{Field:} Education and Public Outreach

Rapid, standardized transient reporting is essential for enabling follow-up across the astronomy community, including observatories with limited dedicated alert infrastructure. The Canadian Hydrogen Intensity Mapping Experiment's (CHIME) Fast Radio Burst (FRB) backend accounts for the majority of known FRBs, making event circulation essential for maximizing the scientific return of its detections. We describe efforts to integrate CHIME/FRB discoveries into established international systems: real-time alert distribution through Virtual Observatory Event (VOEvent) alerts and NASA's Global Coordinates Network (GCN), and standardized naming and cataloguing through the IAU Transient Name Server (TNS). Automating transient alerts introduces a fundamental tension between speed and reliability. We discuss where automation is appropriate, where human supervision remains necessary, and how those boundaries were negotiated in practice to balance rapid dissemination with the responsibility of contributing trustworthy data to shared global databases. By distributing discoveries through channels that observatories worldwide already monitor, this work aims to lower the barrier to FRB follow-up and broaden participation in time-domain astronomy. We reflect on lessons relevant to other Canadian facilities seeking to connect with global transient networks.

\hypertarget{poster-abstract-talk-270}{}\label{poster-abstract-page-talk-270}
\subsection*{Poster 48. AI: Amanuensis Intelligence in Astronomy}
\textbf{Presenter:} Martin Connors\\
\textbf{Poster number:} 48\\
\textbf{Field:} Education and Public Outreach

The now-obscure term amanuensis may come back into common use since artificial intelligence systems permit us many of the functions exceeding those of ancient high-level scribes. The best scribes went far beyond merely copying what their employers said, adding detailed knowledge to their texts. Julius Caesar dictated his books to scribes who had likely accompanied his army and viewed the battles, while Winston Churchill employed many secretaries who were highly trusted and knowledgeable. Most Artificial Intelligence is based on advanced correlation which can mimic knowledge. Much of what we want to impart to students is how to do correlative tasks which help them understand the universe. In some cases, this would be supplemented by analytical skills at which present AI implementations do not excel but nonetheless can frame. I have found that the knowledge of popular AIs on general astronomy has improved greatly in scope over the past two years. The level at which answers start to be incomplete or incorrect has also increased. The tendency to hallucinate has decreased, particularly if one specifies use of correct references. The ability to look things up on the web if updating past the formal training set is needed, has also improved markedly. For specific reasons, AIs are strong in programming in Python, to the point that vibe coding, basically using AI as a 6th generation language, makes sense if one is equipped to evaluate the results critically. All of these developments call for a critical yet positive look at how AI can be used in astronomy education. Anecdotally, Aristotelians refused to look through Galileos telescope through fear of what it might show. We must view AI in a broader context than as a facilitator of student cheating.

\hypertarget{poster-abstract-talk-109}{}\label{poster-abstract-page-talk-109}
\subsection*{Poster 49. Education and Public Outreach at the Trottier Space Institute}
\textbf{Presenter:} Eesha Das Gupta\\
\textbf{Poster number:} 49\\
\textbf{Field:} Education and Public Outreach

The Trottier Space Institute at McGill University has an active and robust outreach arm. The TSIs outreach operations run in conjunction with other units at McGill University as well as other astronomy institutes and centers within the Greater Montreal Area. Additionally, there are programs that the TSI offers which are unique to our operations as an interdisciplinary space institute. Our collaborative efforts make it possible to offer some distinctive space focused educational and outreach opportunities in Montreal including school programs, observing nights, public talks, and more. Moreover, our ethos of providing hybrid programming or an event recording on our YouTube channel where possible, as well as an active social media account provides us with a high impact virtual platform to extend our outreach efforts to a wider audience. This presentation will highlight TSIs education and public outreach initiatives within the Greater Montreal Area, and its continued public impact, both in-person and online.

\hypertarget{poster-abstract-talk-206}{}\label{poster-abstract-page-talk-206}
\subsection*{Poster 50. Supermassive Black Hole Binaries in ZTF AGN Light Curves: Toward Rubin-LSST Searches}
\textbf{Presenter:} Sara Babic\\
\textbf{Poster number:} 50\\
\textbf{Field:} Supermassive black hole

The Vera C. Rubin Observatory Legacy Survey of Space Time (LSST) will deliver densely-sampled, multi-band light curves for millions of active galactic nuclei (AGN), with exceptional cadence in the Deep Drilling Fields. This data will provide an unprecedented opportunity to search for supermassive black hole binaries (SMBHBs), which are predicted in hierarchical galaxy formation but are still observationally elusive.

Motivated by the challenges of upcoming LSST AGN light curves, we present a pre-cursor analysis on archival Zwicky Transient Facility (ZTF) AGN light curves using EztaoX, a scalable algorithm designed to search for SMBHBs via quasi-periodic variability in stochastic AGN light curves. EzTaoX uses multi-band Gaussian Process models to detect evolving, non-sinusoidal quasi-periodic behaviours, as predicted by simulations of circumbinary accretion in SMBHBs. We apply this method to previously-reported SMBHB candidates from the literature, and reassess the evidence for periodic variability and demonstrate EzTaoX's improved recovery of complex quasi-periodic signatures in AGN light curves.

\hypertarget{poster-abstract-talk-185}{}\label{poster-abstract-page-talk-185}
\subsection*{Poster 51. Broad Absorption Line Quasar Spectra Modelling Using SimBAL}
\textbf{Presenter:} Nicolas Frisella\\
\textbf{Poster number:} 51\\
\textbf{Field:} Supermassive black hole

Most massive galaxies contain supermassive black holes, and those that are actively accreting matter and producing light that outshines the host galaxy by many magnitudes are classified as quasars. Quasars with outflows of relativistic speeds directed toward the observer will appear blue-shifted from the rest frame and exhibit broad absorption features. These objects are categorized as broad absorption line (BAL) quasars, and the spectra of their outflows have been observed to vary over time for a select number of these objects. The source of this variability is not yet fully understood, though it is theorized to depend on a set of physical parameters such as the ionization of the gas, column density, and outflow velocities. The spectral fitting algorithm SimBAL has proven successful in modelling the physical properties of quasar winds using a forward-modelling and Markov-Chain-Monte-Carlo approach. This software package fits synthetic spectra to observed spectra for a given quasar, producing detailed physical and kinematical descriptions of these outflows. Using data from the Sloan Digital Sky Survey, we selected a sample of quasars with multiple epochs of observation that demonstrated variability in their absorption lines (particularly C IV and Si IV). We modelled these objects with SimBAL using the local supercluster at Western University to obtain a set of best fit parameters that characterized the spectra we observed.

\hypertarget{poster-abstract-talk-120}{}\label{poster-abstract-page-talk-120}
\subsection*{Poster 52. JWSTs First Detections of Mid-IR Flares from Sgr A*}
\textbf{Presenter:} Daryl Haggard\\
\textbf{Poster number:} 52\\
\textbf{Field:} Supermassive black hole

Sgr A*, the black hole at the heart of the Milky Way Galaxy, has long been observed as a variable source. It displays high contrast X-ray flares and near-continuous variability at near IR wavelengths. Recently, weve detected the first flares at mid-IR wavelengths with the James Web Space Telescopes MIRI IFU. I will present these new MIR flare detections and place them in the context of variability observed at other wavelengths, including by the EHT/SMA/ALMA (sub-mm), Chandra (X-ray), NuSTAR (hard X-ray) and others. Together these observations offer a deep, multi-wavelength, time-domain view of the closest supermassive black hole and allow us to constrain the physical and radiative processes that are driving variability near the event horizon. This JWST/MIRI program has also led to excellent Galactic Center science at larger scales and Ill briefly describe improved solutions for Galactic extinction at MIR wavelengths, constraints on an IMBH in the Galactic Center, new abundance measurements for diffuse gas near Sgr A*, and new detections of a GC X-ray transient, MAXI J1744-294.

\hypertarget{poster-abstract-talk-123}{}\label{poster-abstract-page-talk-123}
\subsection*{Poster 53. Constraining the UV Luminosity of Sgr A* Flares}
\textbf{Presenter:} Mayura Balakrishnan\\
\textbf{Poster number:} 53\\
\textbf{Fields:} High Energy Astrophysics $\cdot$ Supermassive black hole

Direct measurements of ultraviolet (UV) emission from Sagittarius A* are effectively impossible because of the extreme extinction toward the Galactic Center (Av $>$ 30 magnitudes). Nevertheless, the UV flux is a key constraint on accretion-flow models and flare energetics, linking near-infrared emission to X-ray processes. We present a new approach to constraining the UV flux density of Sgr A* and its immediate environment using JWST/MIRI observations and mid-infrared spectral diagnostics sensitive to UV excitation. Leveraging JWST/MIRIs spatial and spectral resolution, we search for temporal variability in UV-excitable lines that would be expected if Sgr A* flares intermittently enhance the local radiation field. We then use CLOUDY simulations to model the irradiated gas in the vicinity of the black hole and identify which mid-infrared lines respond most strongly to changes in the UV field. By comparing the line intensities and spectra predicted by the model with observational constraints derived from the absence of variability in the UV-sensitive lines, we place upper limits on the UV luminosity of Sgr A* in its flaring state.

\hypertarget{poster-abstract-talk-36}{}\label{poster-abstract-page-talk-36}
\subsection*{Poster 54. A cooler look at the environment of Cygnus X-1: searching for dynamical interactions within cold molecular gas}
\textbf{Presenter:} Pau Bosch-Cabot\\
\textbf{Poster number:} 54\\
\textbf{Field:} High Energy Astrophysics

Outflows from compact objects and massive stars play a key role in shaping the interstellar medium. Through injecting significant amounts of energy and momentum, outflow impact sites act as natural calorimeters, providing a way to quantify the energy required to form and maintain them. Stellar-mass black holes existing in binary systems in our Galaxy can produce multiple outflow mechanisms, including relativistic jets and stellar winds, making it essential to disentangle their respective feedback signatures. We present IRAM-30m observations of the stellar mass black hole system Cygnus X-1, which hosts persistent relativistic jets and a strong donor star stellar wind, aimed at probing and characterizing the interaction between these outflows and the surrounding interstellar medium. Our analysis reveals evidence that both the stellar wind and the relativistic jets are actively shaping the environment around Cygnus X-1, and demonstrates that molecular line imaging can effectively disentangle complex feedback processes in environments where multiple outflows take place.

\hypertarget{poster-abstract-talk-87}{}\label{poster-abstract-page-talk-87}
\subsection*{Poster 55. Improved Modeling for Moving Sources of Radio Frequency Interference and Impact on Flagging Strategies}
\textbf{Presenter:} Jade Ducharme\\
\textbf{Poster number:} 55\\
\textbf{Field:} Cosmology and Early Universe

The Epoch of Reionization (EoR) represents a rich and informative period in the early Universe's history. Accessing it through the 21-cm signal of neutral hydrogen remains an immense observational challenge due to the sheer faintness of the signal relative to the radio foreground. In particular, radio frequency interference (RFI), often several orders of magnitude brighter than astronomical foregrounds, poses a significant obstacle to 21-cm experiments. We introduce a new modeling scheme for moving sources of RFI based on the characteristic imprint they leave in the uv-plane. We demonstrate that, given a known source trajectory in image space, one can analytically solve for and efficiently simulate the corresponding sinc-like pattern this source traces in uv-space. Fast and accurate RFI models are essential for ultimately subtracting these contaminants from interferometric data directly rather than simply flagging and discarding affected visibilities, thereby preserving a higher fraction of usable data. In addition, we show that the sinc-like structure of the pattern, which leads to depression of the RFI signal on some uv-modes, prevents per-baseline flagging algorithms from identifying contamination on some baselines. We find the maximum RFI attenuation from this effect only approaches three orders of magnitude between the most and least contaminated baselines, which is likely insufficient to prevent unflagged residual RFI from biasing an EoR analysis. We therefore advocate for per-observation flagging techniques, due to their ability to aggregate visibility information across all baselines, allowing information from highly contaminated uv-modes to guide RFI identification on the suppressed modes.

\hypertarget{poster-abstract-talk-311}{}\label{poster-abstract-page-talk-311}
\subsection*{Poster 56. Searching for Fast Radio Transients using Polarization Parameters: Initial simulations and results}
\textbf{Presenter:} Audrey Bernier\\
\textbf{Poster number:} 56\\
\textbf{Field:} Transients

In recent years, the discovery of Galactic long period transients (LPTs) with burst durations ranging from seconds to minutes has revealed a new portion of the radio transient phase space. However, due to their longer durations and sporadic nature, identifying them with radio telescopes remains a major challenge. Similarly, identifying longer duration fast radio bursts can be difficult with current pulse detection algorithms.

Compared to an approach that uses only the total intensity of the signal, searching for these bursts using their polarization parameters can potentially identify astrophysical transients that are currently being missed in ongoing searches. Detecting more of these events could offer a valuable opportunity to study the nature of highly polarized sources and probe the otherwise invisible ionized medium and magnetic field between stars and galaxies.

In this poster, I will first motivate the need for a polarization-based search. I will then discuss early results from searches in Q and U space of simulated bursts, focusing specifically on the computational cost of these routines. I will end by discussing next steps for this work and implementation for instruments such as CHIME and CHORD.

\hypertarget{poster-abstract-talk-84}{}\label{poster-abstract-page-talk-84}
\subsection*{Poster 57. Inferring the spins of merging black holes in the presence of data-quality issues}
\textbf{Presenter:} Rhiannon Udall\\
\textbf{Poster number:} 57\\
\textbf{Field:} Gravitational waves

Gravitational waves from black hole binary mergers carry information about the component spins, but inference is sensitive to analysis assumptions, which may be broken by terrestrial noise transients known as glitches. Using a variety of simulated glitches and gravitational wave signals, we study the conditions under which glitches can bias spin measurements. We confirm the theoretical expectation that inference and subtraction of glitches invariably leaves behind residual power due to statistical uncertainty, no matter the strength (signal-to-noise ratio; SNR) of the original glitch. Next we show that low-SNR glitches --- including those below the threshold for flagging data-quality issues --- can still significantly bias spin inference. Such biases occur for a range of glitch morphologies, even in cases where glitches and signals are not precisely aligned in phase. Furthermore, we find that residuals of glitch subtraction can result in biases as well. Our results suggest that joint inference of the glitch and gravitational wave parameters, with appropriate models and priors, is required to address these uncertainties inherent in glitch mitigation via subtraction.

\hypertarget{poster-abstract-talk-182}{}\label{poster-abstract-page-talk-182}
\subsection*{Poster 58. Chasing the dragon: forecasting the Draconid meteor shower for spacecraft safety}
\textbf{Presenter:} Auriane Egal\\
\textbf{Poster number:} 58\\
\textbf{Field:} Solar System

Predicting the meteoroid environment is essential for the safety of spacecraft operating in Earth orbit and across the Solar System. Meteoroid impacts remain a significant hazard for satellites, deep-space probes, and human spaceflight. In 1993, the communications satellite Olympus-1 was lost following damage consistent with a meteoroid impact during the Perseid meteor shower. More recently, a micrometeoroid strike damaged the primary mirror of the James Webb Space Telescope in May 2022. Even millimetre-sized particles can degrade sensitive instruments, while larger meteoroids could threaten astronauts during extravehicular activities. Accurate meteor shower forecasts are therefore a key component of space situational awareness and mission risk mitigation.

The Draconid meteor shower, produced by dust trails from the comet 21P/GiacobiniZinner, is renowned for generating intense, short-lived outbursts and storms. These events exhibit strong mass segregation, causing large variations in the particle sizes encountered by Earth and making their intensity difficult to predict. Because the Draconids can reach storm-level activity, they represent a potential hazard for spacecraft.

We present a new numerical model of the Draconid meteoroid stream designed to improve predictions of its particle environment near Earth. Combining updated dust ejection scenarios, refined dynamical integrations, and recent meteor observations, this work aims to better constrain the flux and size distribution of Draconid particles and improve hazard assessment for current and future space missions.

\programdayheader{Wednesday}
\daypalettebanner{Wednesday}

\hypertarget{poster-abstract-talk-27}{}\label{poster-abstract-page-talk-27}
\subsection*{Poster 1. Simulating transients for the next generation of surveys}
\textbf{Presenter:} Connor Stone\\
\textbf{Poster number:} 1\\
\textbf{Fields:} Cosmology and Early Universe $\cdot$ Transients

Here we present the Cosmographi simulation suite for accurate simulation of transients (Supernovae, TDEs, AGNs, Kilonovae, and lensed variants) in next generation survey programs (the LSST). The LSST will generate orders of magnitude more transients than previously observed, in principle allowing for vastly tighter constraints on inferred parameters. However this statistical power is inaccessible without detailed and accurate simulation suites to pair with the data. Using JAX acceleration we may run simulations on GPUs to gain orders of magnitude acceleration over traditional sampling techniques. Allowing us to access techniques that were previously infeasible in realistic timescales. In this talk we will demonstrate how Cosmographi can generate large simulated light curve datasets in record time, and how it can integrate with other cosmological analysis tools like caustics and AstroPhot to perform realistic image simulations of lensed transients.

\hypertarget{poster-abstract-talk-285}{}\label{poster-abstract-page-talk-285}
\subsection*{Poster 2. Corotating Interaction Regions: A polarizing topic [polarimetry]}
\textbf{Presenter:} Luca Fabiani\\
\textbf{Poster number:} 2\\
\textbf{Field:} Stars and Stellar cluster formation

WolfRayet stars, the evolved descendants of stars initially more massive than $>$20 M$_\odot$, undergo intense mass loss that drives dense, optically thick stellar winds. These outflows are highly structured across a wide range of spatial and temporal scales. Among the mechanisms proposed to shape them, Co-rotating Interaction Regions (CIRs) provide a compelling framework for explaining several forms of large-scale variability, particularly those revealed by polarization diagnostics. WR 6 (EZ CMa) is one of its most prominent candidates for hosting CIRs, displaying strong time-dependent variability modulated on a 3.76-day period. Using spectropolarimetric datasets obtained with ESPaDOnS at the CFHT, we detect strong linear polarization variations across several He II lines as well as lines of other ions such as He I, C IV, N IV, and N V, probing different wind regions. For He II lines in particular, the polarization variations trace loops in the QU plane that appear on the blue side of the line, while the red side shows mostly linear excursions. The loops can be traced in both clockwise and counter-clockwise directions from night to night, indicating a complex and evolving wind geometry. In the framework of CIR-induced asymmetries, we find that the observed behavior cannot be explained solely by dilution of continuum polarization by unpolarized line emission (the classical line effect). We therefore investigate models in which the polarization variations observed on the blue side of the lines arise from partial occultation of the polarized Thomson photosphere by CIR structures. This approach provides a natural explanation for the blue-side variability seen in the He II lines and forms the basis of ongoing efforts to reconcile the continuum and line polarization variations.

\hypertarget{poster-abstract-talk-286}{}\label{poster-abstract-page-talk-286}
\subsection*{Poster 3. Corotating Interaction Regions: Youd have to see it to believe it}
\textbf{Presenter:} André-Nicolas Chené\\
\textbf{Poster number:} 3\\
\textbf{Field:} Stars and Stellar cluster formation

WolfRayet (WR) stars, the evolved descendants of stars initially more massive than $>$20 Msol, undergo intense mass loss that drives dense, optically thick stellar winds. These outflows are highly structured across a wide range of spatial and temporal scales. Among the mechanisms proposed to shape them, Co-rotating Interaction Regions (CIRs) provide a compelling framework for explaining several forms of large-scale variability, particularly those revealed by polarization diagnostics.

In O-type stars, CIRs account for the Discrete Absorption Components observed in ultraviolet resonance lines. The same mechanism has been invoked in WR winds, although direct observational evidence remains limited. The best candidate is WR 6 (EZ CMa), which exhibits persistent spectroscopic and photometric variability with a 3.76-day period widely interpreted as CIR modulation. Similar behavior has been reported in WR 1, WR 110, and WR 134 who share comparable characteristics in photometry, spectroscopy, polarimetry and/or X-ray emission.

Until now, the CIR interpretation in WR stars has relied primarily on indirect diagnostics. Resolving this ambiguity requires probing the geometry of the inner wind itself. Long-baseline interferometry offers precisely this capability: rotating spiral structures are predicted to imprint characteristic signatures in spectrally dispersed visibility and phase. With modern facilities now able to resolve WR winds on sub-au scales, interferometry provides a direct test of the CIR hypothesis. In this context, WR 6 is an ideal target.

Here we present multi-epoch GRAVITY interferometric observations aimed at directly probing the morphology of its inner wind.

\hypertarget{poster-abstract-talk-308}{}\label{poster-abstract-page-talk-308}
\subsection*{Poster 4. Corotating Interaction Regions : a roll model}
\textbf{Presenter:} Guillaume Lenoir-Craig\\
\textbf{Poster number:} 4\\
\textbf{Field:} Stars and Stellar cluster formation

WolfRayet (WR) stars, the evolved descendants of stars initially more massive than $>$20 M$_\odot$, undergo intense mass loss that drives dense, optically thick stellar winds. These outflows are highly structured across a wide range of spatial and temporal scales. Among the mechanisms proposed to shape them, Co-rotating Interaction Regions (CIRs) provide a compelling framework for explaining several forms of large-scale variability.

Hydrodynamical modelling coupled with radiative transfer calculations is the ideal theoretical approach to constrain both the causes and observational consequences of these large-scale structures. While numerical investigations of the CIR phenomenon have been conducted in the context of optically thin O-star winds in the past, their formation and observational signatures in the dense, optically thick winds of WolfRayet stars remain poorly explored.

In this work, we present three-dimensional radiationhydrodynamical simulations of CIRs in WR winds. The simulations incorporate parametrized radiative driving to model the wind acceleration. Surface perturbations are introduced to emulate localized brightness variations that generate CIRs in the outflow. The resulting density and velocity structures are then post-processed with a 3D radiative transfer code to compute synthetic observables, including spectral-line variability and continuum emission.

Our results aim to establish the link between surface perturbations, large-scale wind structures, and observable variability in WR stars. This approach provides a framework to interpret spectroscopic variability and offers new constraints on the physical mechanisms responsible for structure formation in the dense winds of massive stars.

\hypertarget{poster-abstract-talk-7}{}\label{poster-abstract-page-talk-7}
\subsection*{Poster 5. TESS Occurrence Rates Reveal the Disappearance of the Radius Valley Around mid-to-late M Dwarfs}
\textbf{Presenter:} Erik Gillis\\
\textbf{Poster number:} 5\\
\textbf{Field:} Exoplanets

Mid-to-late M dwarfs are prolific planet hosts and represent the most common type of star in the galaxy, but due to their faintness, transit surveys have yet to robustly characterize their planet population. The all-sky coverage of TESS has enabled the deepest systematic characterization of the mid-to-late M dwarf planet population to date. We survey 8134 mid-to-late M dwarfs ($>$0.4 M\_sun) observed by TESS with a custom built transit pipeline and recover 77 vetted transiting planet candidates. We characterize the sensitivity of our survey via injection--recovery and measure the occurrence rate of planets as a function of orbital period, instellation, and planet radius. In particular, we find that mid-to-late M dwarfs host an average of 1.100.16 planets with radii $<$ 1 R\_earth orbiting within 30 days. Unlike the bimodal Radius Valley exhibited by close-in planet populations around FGK and early M dwarfs, we recover a unimodal planet radius distribution peaking at 1.250.05 R\_earth. The plentiful population of sub-Neptunes which form the second peak of the FGK radius valley is practically absent around stars in our sample, with super-Earths out numbering sub-Neptunes 5.5:1. In essence, our work shows that the Radius Valley disappears around the lowest mass stars, and that the most common planets in our galaxy do not frequently exist around its most common stars. This dearth of sub-Neptunes around mid-to-late M dwarfs is consistent with water-rich pebble accretion models, supporting the emerging idea that the M Dwarf sub-Neptune population is composed of water-rich worlds.

\hypertarget{poster-abstract-talk-186}{}\label{poster-abstract-page-talk-186}
\subsection*{Poster 6. Searching for Transiting Exoplanets in Romans GBTDS}
\textbf{Presenter:} Kelsey Hoffman\\
\textbf{Poster number:} 6\\
\textbf{Field:} Exoplanets

The Nancy Grace Roman Space Telescope is on track for a Fall 2026 launch. A core component of its five-year mission is the Galactic Bulge Time-Domain Survey (GBTDS). The Transits in the Roman galactic Exoplanet Survey (TRExS) team has been developing infrastructure to leverage these observations in order to detect transiting exoplanets. The infrastructure for this program is based on proven techniques developed for exoplanet missions such as, Kepler, K2 and TESS. A key component of the program's success will be the transit search software pipeline and validation toolset, being developed by Canadian researchers, which will enable the identification and confirmation of new exoplanet candidates. Initial estimates suggest the potential for the discovery of 100,000 new transiting planets, the Roman Space Telescope Transit Program is poised to revolutionize our knowledge of planetary systems beyond our own.

\hypertarget{poster-abstract-talk-276}{}\label{poster-abstract-page-talk-276}
\subsection*{Poster 7. Synthetic modelling of exoplanet transit light curve using state-of-the-art model stellar atmosphere limb darkening}
\textbf{Presenter:} Ishan Kaushal\\
\textbf{Poster number:} 7\\
\textbf{Field:} Exoplanets

Variations in intensity, and flux from the center to the limb of the host stars play crucial roles in the photometric study of exoplanetary transits. Flux variations occur when a planet passes in front of a star, blocking a portion of the star's total light. Additionally, the planets transit across the host star breaks the circular symmetry, resulting in the change in linear polarization. In this project, I developed an exoplanet transit code that calculates the flux for a given normalized planetary radius (p) using the stellar limb-darkening model. The primary goal of this project is to explore the variations in flux during transit for different exoplanetary samples and to probe the properties of stellar and exoplanetary atmospheres. In this project, I find that a model assuming a planet with no atmosphere i.e a planet having a constant radius can yield a varying transit spectrum. I also calculated the spectra for a given limb darkening model as a function of wavelength.

\hypertarget{poster-abstract-talk-102}{}\label{poster-abstract-page-talk-102}
\subsection*{Poster 8. Using flare timing in TESS light curves to uncover the starspot distributions on fully and partially convective M dwarfs}
\textbf{Presenter:} Rohan Kumar\\
\textbf{Poster number:} 8\\
\textbf{Fields:} Exoplanets $\cdot$ Stars and Stellar cluster formation

The surface distribution of starspots on M dwarfs remains an open question in stellar astrophysics. Constraining the locations of starspots gives us insight into how stellar magnetic dynamos are generated while improving our ability to accurately model precision radial velocity data and atmospheric transmission spectra, both of which are contaminated by starspots. Starspot mapping techniques, such as Zeeman Doppler Imaging, are observationally expensive, making them poorly-suited to obtaining a population-level understanding. Alternatively, stellar flares, the result of magnetic reconnection events, are spatially correlated with starspots observed on the Sun, thus making flare observations a powerful diagnostic of starspot positions on distant stars. We have developed and rigorously tested a framework to detect stellar rotation and energetic flares in light curves from NASAs Transiting Exoplanet Survey Satellite, which we are using to correlate flare times with rotational phase, a proxy for starspot positions. In our preliminary sample of 100 partially convective and 27 fully convective M dwarfs, fast rotators (0.3d $<$ Prot $<$ 6d) show enhanced flaring when starspots face the observer, consistent with equatorial spots. Conversely, ultra-fast rotators (Prot $<$ 0.3d) appear to show no rotational phase dependence, suggesting persistent polar spots that remain in view as the star rotates. The differing distributions of starspots between the two populations may imply a fundamental transition in the underlying magnetic dynamo.

\hypertarget{poster-abstract-talk-79}{}\label{poster-abstract-page-talk-79}
\subsection*{Poster 9. Rotation Periods of Candidate Single Late-M Dwarfs in TESS}
\textbf{Presenter:} Samantha Lambier\\
\textbf{Poster number:} 9\\
\textbf{Fields:} Exoplanets $\cdot$ Stars and Stellar cluster formation

Recent work has shown that the rotational evolution of late-M and brown dwarfs differs from the well-known spin-down behaviour of hotter, Sun-like stars. Characterizing the distribution of rotation periods of these objects in the solar neighbourhood provides insight into the evolutionary pathway near the stellar-substellar boundary. Rotation periods are also important for exoplanet searches, as they help inform stellar inclination and identify signals caused by stellar activity rather than planets. In this presentation, we examine 1789 candidate single late M dwarfs ($\geq$M6) using light curves from the Transiting Exoplanet Survey Satellite (TESS). Our sample includes some of the faintest late-M dwarfs for which rotation periods have been measured with TESS, extending to T $\sim$ 18 mag and pushing the mission to its sensitivity limits. Rotation periods are identified using Lomb-Scargle periodograms and refined with Gaussian Process modelling. We measured 174 rotation periods, ranging from 2 hours to 6 days, with variability amplitudes between 0.05\% and 3\%. The fraction of stars exhibiting detectable variability increases with the number of available TESS sectors, approaching an apparent ceiling of $\sim$50\%, likely reflecting geometric viewing constraints. This supports the idea that spot-driven variability is widespread among late-M dwarfs, similar to what is observed in cooler brown dwarfs. When combined with published rotation periods for a broader range of spectral types, we find a lower envelope on the rotation period decreasing from 5 hours in early-M dwarfs to one hour in L, T, and Y dwarfs. These findings help refine our understanding of low-mass stellar rotation and provide practical inputs for future exoplanet surveys targeting late-M dwarf systems.

\hypertarget{poster-abstract-talk-232}{}\label{poster-abstract-page-talk-232}
\subsection*{Poster 10. Improving Stellar Modelling to Reduce Uncertainties in Atmosphere Characterisation for Ongoing JWST Observations of Five Lava Planets}
\textbf{Presenter:} Owen Lammert\\
\textbf{Poster number:} 10\\
\textbf{Field:} Exoplanets

Understanding the spectra of a host star is of critical importance if spectral analysis of a planets atmosphere is being performed. This is especially important as the James Webb Space Telescope (JWST) obtains more and more mid-infrared instrument (MIRI) observations of planetary systems whose host stars are poorly characterised in the mid-infrared (mid-IR) wavelengths, or might not have been robustly characterised in all their parameters. We aim to assist ongoing observations using JWST that will give an unprecedented level spectrographic precision of five stars hosting lava planets by characterising their stellar hosts properties. By comparing this highly detailed data to models for stellar spectra in the mid-IR, such as those described by the PHOENIX models, stellar parameters such as metallicity, effective temperature, and surface gravity can be better constrained. These constrained values can then be applied to improve the accuracy of stellar models and used to produce theoretical models of lava planets which will be compared with the JWST observations for interpretation. For these reasons, it is critical to have a well understood stellar spectra when analysing a planets atmosphere. With the JWST data for these five systems arriving over the next year, these stars are in need of careful analysis so that their accompanying planets can also be more accurately characterised, as well as to avoid any uncertainties in composition and temperatures that could be avoided by simply having more complete stellar models.

\hypertarget{poster-abstract-talk-299}{}\label{poster-abstract-page-talk-299}
\subsection*{Poster 11. Disentangling Stellar Activity and Planetary Features with New Models of M Dwarf Flares in the NIR}
\textbf{Presenter:} Élise Leclerc\\
\textbf{Poster number:} 11\\
\textbf{Field:} Exoplanets

Since its launch, the James Webb Space Telescope (JWST) has enabled transformative advances in the characterization of exoplanetary atmospheres. In particular, it has provided unprecedented insights into rocky exoplanets located in the habitable zone, where liquid water can exist on the surface given the presence of an atmosphere. Many of these exoplanets orbit small and cool stars called M dwarfs, which are key targets for transmission spectroscopy. However, M dwarfs are highly magnetically active. Their powerful eruptions, known as flares, introduce space-, time- and wavelength-dependent contamination in the transmission spectra of orbiting exoplanets, masking or mimicking key molecular absorption features such as H$_2$O, CO$_2$, and CH$_4$. To date, these unpredictable bursts of radiation remain largely unconstrained in the near-infrared (NIR), a spectral regime of interest for exoplanet atmospheric studies with JWST. Flares have already been detected in multiple JWST observing programs and have significantly hindered the interpretation of planetary features in transmission spectra. Accurate characterization of M dwarf flares in the NIR is therefore essential to better probe exoplanetary atmospheres and fully exploit JWSTs potential. In this talk, I will present new constraints and models of flare activity across the NIR regime derived from JWST/NIRISS observations of highly active M dwarfs. I will characterize the energetic, temporal and spectral properties of these flares and compare them with existing optical and ultraviolet models. These results will help constrain the NIR flare continuum of M dwarfs and allow the scientific community to more reliably correct stellar contamination in the transmission spectra of potentially habitable exoplanets.

\hypertarget{poster-abstract-talk-90}{}\label{poster-abstract-page-talk-90}
\subsection*{Poster 12. Pushing Coherent Differential Imaging to its Limits with SPIDERS}
\textbf{Presenter:} Christopher Mann\\
\textbf{Poster number:} 12\\
\textbf{Fields:} Exoplanets $\cdot$ Instrumentation

The Subaru Pathfinder Instrument for Detecting Exoplanets and Recovering Spectra (SPIDERS) is an exploratory instrument testing new techniques and procedures for the next generation of exoplanet direct-imaging instrumentation. Here we discuss the Coherent Differential Imaging (CDI) post-processing technique as a way to substantially boost imaging contrast: including theoretical performance, real-world limitations due to chromatic fringe blurring, lab-bench performance, as well as first steps toward a machine learning algorithm to overcome certain limitations.

\hypertarget{poster-abstract-talk-218}{}\label{poster-abstract-page-talk-218}
\subsection*{Poster 13. Toward Directly Imaging Earth-like rocky planets with Focal Plane Wavefront Sensing: SPIDERS On-Sky Results at Subaru and CAL2 Gemini Planet Imager Project Update}
\textbf{Presenter:} Christian Marois\\
\textbf{Poster number:} 13\\
\textbf{Fields:} Exoplanets $\cdot$ Instrumentation

The quest to detect and characterize exoplanets and debris disks around nearby stars will be transformed in the coming years with the development of the next generation of high-contrast imagers. These new/upgraded instruments aim to achieve up to a 100 gain in sensitivity at small $\sim$$\lambda$/D angular separations around bright stars compared to current facilities such as Gemini Planet Imager and SPHERE. Over the past decade, significant R\&D at the National Research Council of Canada has focused on developing focal-plane wavefront sensing techniques that enable measurements of the residual electric field directly in the science focal plane. These measurements can be used in closed loop with the instrument deformable mirror to create deep dark holes in the quasi-static dominated stellar halo, enabling the detection of faint disks and exoplanets. Using this approach, the science data can also be enhanced through coherent differential imaging, which leverages the coherence of the residual starlight to further suppress speckle noise in addition to standard techniques such as angular and spectral differential imaging.

I will present on-sky results from the NRC SPIDERS pathfinder experiment obtained at Subaru Telescope in late 2025 and early 2026, validating these breakthrough concepts. I will also discuss progress toward deploying a similar facility-class system, CAL2, for the Gemini Planet Imager 2. Finally, I will outline how SPIDERS is being used to explore extensions of these techniques, including their combination with very high-resolution spectroscopy at extreme contrasts for future instruments such as ANDES and PCS on the Extremely Large Telescope, and the NASA Habitable Worlds Observatory.

\hypertarget{poster-abstract-talk-318}{}\label{poster-abstract-page-talk-318}
\subsection*{Poster 14. A Panoramic View of the Milky Way Galaxy with the Australian SKA Pathfinder}
\textbf{Presenter:} Roland Kothes\\
\textbf{Poster number:} 14\\
\textbf{Fields:} Galaxy $\cdot$ Gas, dust and the ISM

With the Evolutionary Map of the Universe (EMU) and the Polarization Sky Survey of the Universe's Magnetism (POSSUM), the touchstone radio continuum surveys of the southern hemisphere are now under way. EMU and POSSUM use the Australian SKA Pathfinder (ASKAP) telescope to image the southern sky to an unprecedented sensitivity and resolution. Covering the southern hemisphere, EMU is ideal for observing the Milky Way Galaxy. I will summarize recent developments in processing procedures for linear polarization radio data, such as RM synthesis and RM clean and the separation of individual Faraday depth components along the line of sight. With these techniques we will produce a sensitive atlas of discrete radio nebulae and a spectacular data set to study the diffuse magneto-ionic medium of the Milky Way Galaxy at unprecedented detail.

\hypertarget{poster-abstract-talk-197}{}\label{poster-abstract-page-talk-197}
\subsection*{Poster 15. Bars in Nearby Galaxies: Decoding their Influence on Gas Concentrations and Star Formation}
\textbf{Presenter:} Jennifer Laing\\
\textbf{Poster number:} 15\\
\textbf{Fields:} Galaxy $\cdot$ Gas, dust and the ISM

More than half of star forming galaxies are barred, but studies of the gas properties in the bar have been limited to small sample sizes. KILOGAS is a survey of the molecular gas in 450 mass-selected galaxies covering the full range of morphologies and environments. Using the Atacama Large Millimeter/Sub-millimeter Array (ALMA), CO(2-1) measurements have unprecedented spatial resolution of 1 kpc for galaxies between redshifts 0.015 and 0.06. This study expands on work by Chown et al. (2019) who found that galaxies with higher molecular gas concentrations and centrally enhanced star formation tend to be barred galaxies. We expect gas concentrations to be higher due to gas inflows along the bar which have been found to suppress star formation efficiency (Hogarth et al. 2024). I will present preliminary results of my analysis using radial profiles to quantify the central concentration of molecular gas and star formation rate in KILOGAS galaxies.

\hypertarget{poster-abstract-talk-189}{}\label{poster-abstract-page-talk-189}
\subsection*{Poster 16. MAGAZ3NE: Environmental effects on ultra-massive galaxies in the early Universe}
\textbf{Presenter:} Han Lei\\
\textbf{Poster number:} 16\\
\textbf{Field:} Galaxy

Ultra-massive galaxies (UMGs; $M_\star \geq 10^{11} M_\odot$) represent the extreme end of galaxy formation and evolution, and provide a natural laboratory for studying stellar mass assembly and galaxy quenching mechanisms. Although it has been well established that, in the local Universe, such massive systems prefer to reside in dense environments and that most are quenched, whether similar environmental dependencies exist at higher redshift remains unclear. Using more than 7,000 galaxies with spectroscopic detections and extensive photometric data in the same sky region, we construct an overdensity map with the Voronoi Monte Carlo (VMC) technique to study the environments of UMGs at $z\sim 3$.

For each galaxy, we characterize its environment by the normalized overdensity at its position, tracing the local environment, and by the distance to the nearest overdensity peak, tracing the large-scale environment. Our sample consists of 10 UMGs $3<z<3.4$ from the MAGAZ3NE survey, all with spectroscopic redshifts and well-sampled photometric spectral energy distributions. Comparing these 10 MAGAZ3NE UMGs with a massive cluster-galaxy sample, a UMG-candidate sample, and the non-UMG galaxies in the same redshift range, we find that at $z\sim 3$: (1) UMGs preferentially reside in denser environments, regardless of how environment is characterized; and (2) quenching in UMGs appears to depend more on the large-scale environment than on the local density.

We further examine our results with simulations. We construct a mock sample from TNG300-1 that observationally matches our survey data, generate corresponding VMC maps, and identify simulation-based UMG samples. The simulation results are consistent with our observational findings and support the same two conclusions above.

\hypertarget{poster-abstract-talk-179}{}\label{poster-abstract-page-talk-179}
\subsection*{Poster 17. Searching for Tidal Structure in the Fornax Globular Clusters with LSST}
\textbf{Presenter:} Alyson Loney\\
\textbf{Poster number:} 17\\
\textbf{Fields:} Galaxy $\cdot$ Local Group $\cdot$ Stars and Stellar cluster formation

Dwarf galaxies are essential to our understanding of galaxy evolution, as their shallow gravitational potential wells preserve dynamical signatures of early formation. The Fornax dwarf spheroidal galaxy is especially valuable in this context due to its proximity within the Local Group and its unusually rich population of six globular clusters (GCs). Dynamical friction models predict that these clusters should have sunk into the galaxy centre, yet they remain at large radii, forming the Fornax dynamical puzzle. Using deep, wide-field imaging from the Legacy Survey of Space and Time (LSST) Data Preview 1, we search for faint outer structures in the Fornax GCs that may indicate ongoing interaction with the galaxy's gravitational potential. We derive surface brightness profiles of the clusters and examine their outer regions for evidence of tidal distortion or faint tidal tails.

\hypertarget{poster-abstract-talk-122}{}\label{poster-abstract-page-talk-122}
\subsection*{Poster 18. A Multi-wavelength Census of Star Formation in Nuclear Rings at 10\,pc Scales, with JWST, HST, VLA and ALMA}
\textbf{Presenter:} Anan Lu\\
\textbf{Poster number:} 18\\
\textbf{Field:} Galaxy

Nuclear rings in barred galaxies are among the most visually striking and physically important star-forming structures in the nearby Universe. They mark the sites of dense gas reservoirs that fuel intense star formation and potentially influence the growth of central supermassive black holes. The associated complicated gas dynamics, high gas surface density, and clustered star formation makes them ideal laboratories for testing star formation diagnostics and theories in extreme circumnuclear environments. With modern facilities, we can now resolve the fundamental star formation units in nearby nuclear rings at unprecedented sensitivity and spatial resolution. We present a multi-wavelength star formation rate study of 10 nearby nuclear rings using JWST Pa$\alpha$ imaging, HST H$\alpha$ observations, and VLA 33\,GHz radio continuum maps, all matched to a common angular resolution of 0.2'' ($\sim$10\,pc). At this resolution, individual H II regions, molecular cloud complexes, and young stellar structures can be directly identified and compared. Although both the recombination-line emission and 33\,GHz free-free emission are expected to originate from massive young stars, we find substantial discrepancies between these tracers in both the brightness and morphology of resolved H II regions. To investigate the origin of these mismatches, we incorporate archival and new observations on: dust properties from JWST PAH and continuum emission, stellar populations from combined HST+JWST spectral energy density fitting, and gas dynamics from ALMA CO observations. One of the key findings is a strong correlation between 33\,GHz surplus and PAH+dust emission, while H$\alpha$ bright and 33\,GHz faint regions tend to have less small grain PAHs and older stars.

\hypertarget{poster-abstract-talk-225}{}\label{poster-abstract-page-talk-225}
\subsection*{Poster 19. Starbursting Galaxy Cluster Cores at Cosmic Noon in TNG-Cluster: Insights into SpARCS1049+56}
\textbf{Presenter:} Bhuvan Manojh\\
\textbf{Poster number:} 19\\
\textbf{Fields:} Galaxy $\cdot$ Galaxy Clusters

Cosmic noon (z$\sim$1.5-2.5) marks a transitional phase for the earliest galaxy clusters to form from protoclusters, yet the behaviour of cluster cores in this era remains poorly understood. The massive galaxy cluster SpARCS1049+56 at z=1.709 provides a unique window into high-redshift cluster core phenomena. It hosts a starburst of $\sim$ 860 solar masses/year spatially offset from the central Brightest Cluster Galaxy (BCG), which has been interpreted as a sloshing event triggering a runaway cooling flow in the absence of radio mode feedback. Using the TNG-Cluster simulations, we investigate whether SpARCS1049-like analogues naturally appear during cosmic noon. Tracing the progenitors of 352 clusters between z=1.5-2.44, we find 284 progenitors with extreme starbursts within the cluster core. While 44\% of all clusters experience extreme starbursts at cosmic noon, centrally-offset star formation is a rare and short-lived event. Centrally-offset systems are characterized by increased dynamical disturbances, ongoing merger activity, and bursts of in-situ stellar growth in the BCG rather than the intracluster light. Our results suggest that extreme star formation is common at the epoch of cluster formation, but centrally-offset events are rare and arise from star-forming satellites rather than runaway cooling flows. SpARCS1049 may therefore be an extreme example within current observed and simulated high-redshift cluster populations.

\hypertarget{poster-abstract-talk-272}{}\label{poster-abstract-page-talk-272}
\subsection*{Poster 20. Cosmic Ray Electron Escape in Dwarf Galaxies: The RCSFR Relation from LoTSS DR3 and the DESI Dwarf Catalog}
\textbf{Presenter:} Natalie Oosterman\\
\textbf{Poster number:} 20\\
\textbf{Fields:} Galaxy $\cdot$ Stars and Stellar cluster formation

We present a radio continuum study of star formation in dwarf galaxies from the DESI Extragalactic Dwarf Galaxy Catalog, the largest dwarf catalog to-date with 647,241 samples. The Third Data Release of the LOFAR Two-Meter Sky Survey (LoTSS DR3) provides 144 MHz radio flux densities which yield a dust-free analysis of star formation with the corresponding DESI galaxies. This study measures the radio continuum-star formation rate (RC-SFR) dependence in dwarf galaxies, with Klein et al. (1991) being the foundational work on $\sim$20 blue compact dwarfs. To estimate the relation, we follow the method outlined in Shenoy et al. (2026) and fit $L\_\mathrm{150} = L\_c\,\mathrm{SFR}^{\beta}\,(M\_\star/10^{10}\,M\_\odot)^{\gamma}$ (Shenoy 2026) via MCMC to our dwarf sample, 10 596 galaxies after cross-matching and implementing redshift cuts (z $<$ 0.2). As few literature RCSFR measurements for dwarf galaxies exist, we first validate our method by comparing the non-dwarf population from the same catalogs (213 345 regular galaxies) against the literature under identical conditions. We find pre-bias-correction values of log L\_c = 22.491 0.011, $\beta$ = 0.300 0.001 and $\gamma$ = 0.226 0.007 for dwarfs, compared to log L\_c = 21.992 0.004, $\beta$ = 0.644 0.007 and $\gamma$ = 0.571 0.008 for a non-dwarf comparison sample fit identically (Shenoy et al. 2026). The dwarf galaxies are systematically radio-faint compared to the regular galaxies, suggesting they are inefficient synchrotron emitters per unit star formation, consistent with predicted cosmic ray electron escape (Klein et al. 1991; Condon 1992). The modestly higher $\gamma$ in dwarfs relative to the non-dwarf sample suggests a steeper stellar mass dependence at low masses. These results provide an empirical anchor for radio-derived SFRs in the dwarf regime, ahead of SKA and LOFAR 2.0 all-sky surveys.

\hypertarget{poster-abstract-talk-208}{}\label{poster-abstract-page-talk-208}
\subsection*{Poster 21. Anisotropic effects in supernova blown bubbles}
\textbf{Presenter:} Nicholas Owens\\
\textbf{Poster number:} 21\\
\textbf{Fields:} Galaxy $\cdot$ Gas, dust and the ISM

Superbubbles are the vast bubbles of hot gas surrounding clustered supernovae. Frequently, they break out of galactic disks, driving outflow. They are responsible for feedback that regulates star formation and changes the baryonic content of the ISM. However, idealised models of supernovae result in bubbles that are too hot or outflows that are under-massive. This implies that thermal conduction, cosmic rays, or both, may play a strong role in the evolution of these bubbles. Both processes must be aligned to magnetic fields, which are ubiquitous at all scales in galaxies. Therefore, the physics governing both must necessarily by anisotropic. In this presentation, we examine the impact of anisotropic conduction and cosmic rays on superbubbles. We also examine how the orientation and of magnetic fields can impact the size, and morphology of the bubbles. All of these processes are inherently small scale, but can have galaxy wide impact.

\hypertarget{poster-abstract-talk-216}{}\label{poster-abstract-page-talk-216}
\subsection*{Poster 22. Inference of Star Formation and Metallicity Histories from Galaxy Spectra with Score-Based Machine Learning Models}
\textbf{Presenter:} Sacha Perry-Fagant\\
\textbf{Poster number:} 22\\
\textbf{Field:} Galaxy

Star formation histories (SFHs) are a key component for understanding galaxy evolution. Because SFHs are not directly observable, they must be inferred from galaxy spectra. However, the mapping from spectra to SFHs is highly degenerate, making the recovery of SFHs an ill-posed inverse problem. As a result, most approaches rely on simple parametric models that impose strong assumptions about the functional form of the history. To generate more realistic histories, we train a score-based diffusion model to act as a prior over SFHs and metallicity histories (MHs). The model is trained on histories from the TNG50 cosmological simulation, where ground-truth values are available. We perform Bayesian posterior sampling by drawing (SFH, MH) pairs from the diffusion model, generating spectra through stellar population synthesis, and comparing them to observed galaxy spectra. Using mock observations from TNG, the method successfully recovers histories consistent with the ground truth, as verified through posterior coverage tests. The approach also generalizes to out-of-distribution histories from other simulations. We apply the method to SDSS galaxy spectra and obtain stellar mass estimates. These results demonstrate that diffusion-based priors provide a flexible framework for inferring realistic galaxy formation histories from spectroscopic observations.

\hypertarget{poster-abstract-talk-12}{}\label{poster-abstract-page-talk-12}
\subsection*{Poster 23. From jet to filaments: a multi-scale and multi-wavelength view of AGN activity in M87}
\textbf{Presenter:} Camille Poitras\\
\textbf{Poster number:} 23\\
\textbf{Fields:} Galaxy $\cdot$ Galaxy Clusters $\cdot$ Supermassive black hole

Active galactic nuclei (AGN) can launch relativistic jets that inject and redistribute energy in their surroundings. As the nearest brightest cluster galaxy in a cool-core cluster, M87 provides a unique laboratory to trace these processes in detail. In this talk, I will first present new constraints on the expansion of its well-known jet in X-rays with Chandra HRC-I, where emission from short-lived synchrotron electrons traces regions near ongoing particle acceleration. The sub-arcsecond imaging capability of HRC-I, combined with a 13-year baseline, enables measurements of proper motions and flux variability of individual knots, constraining their velocities, cooling behavior, and temporal evolution in a multi-wavelength context. On larger scales, jet-driven energy influences the extended warm filamentary nebula in M87. Using new integral-field spectroscopy observations from MEGARA (GTC) and SITELLE (CFHT), we provide the first complete views of its kinematics and excitation across the entire network. The gas exhibits turbulent motions and AGN-dominated excitation and shows strong coupling to both the cold molecular phase and the cooler component of the hot intracluster medium. Together, these results provide new insight into AGN feedback across spatial scales, wavelengths, and timescales.

\hypertarget{poster-abstract-talk-155}{}\label{poster-abstract-page-talk-155}
\subsection*{Poster 24. CFHT News and Updates}
\textbf{Presenter:} Nadine Manset\\
\textbf{Poster number:} 24\\
\textbf{Field:} Instrumentation

We will present updates on recent and ongoing initiatives at the Canada-France-Hawaii Telescope (CFHT) aimed at enhancing its scientific capabilities, including the upcoming Community / Legacy Surveys, community-driven scientific surveys expected to start in 27B, and Wenaokeao, which will allow simultaneous spectropolarimetric observations with ESPaDOnS and SPIRou, from 370\,nm to 2.44 microns. In addition, CFHT users now benefit from new tools and services: a redesigned Python-based MegaCam Exposure Time Calculator called pyDIET and programmatic access tools provided through the Astronomical Event Observatory Network (AEON) starting in 2026B. These AEON tools will be available to PIs of accepted standard Target-of-Opportunity programs using MegaCam, ESPaDOnS, or SPIRou. Together, these developments expand CFHTs flexibility, responsiveness, and scientific reach, particularly in time-domain astronomy.

\hypertarget{poster-abstract-talk-195}{}\label{poster-abstract-page-talk-195}
\subsection*{Poster 25. TH\'ES\'EE and the Crab Nebula: a new full visible bandpass hyperspectral imager and a not that old supernova remnant}
\textbf{Presenter:} Thomas Martin\\
\textbf{Poster number:} 25\\
\textbf{Field:} Instrumentation

The family of unique astronomical instruments known as hyperspectral imagers, developed at Universit Laval by Prof. Laurent Drissen, continues to grow with TH\'ES\'EE. This prototype, currently under development, will offer dramatically improved observational efficiency compared to its two predecessors: SpIOMM (the original prototype installed at the Mont-Mégantic Observatory) and SITELLE (the enhanced version of SpIOMM, installed at the Canada-France-Hawaii Telescope). We are relying primarily on the improved control system implemented this year and, above all, on the use of EMCCD cameras to revolutionize the observational process, enabling us to collect nearly three times more data than before in the same amount of time. On anther note, and thanks to its exceptional capabilities, SITELLE has also provided a renewed view of many astrophysical objects, including the Crab Nebula. After achieving its 3D reconstruction and measuring its expansion in three dimensions, two new datasets that no other instrument in the world could have obtained, we will also present a precise estimate of the 3D distribution of the ionized material mass, and discuss some of the many questions raised by these new data.

\hypertarget{poster-abstract-talk-247}{}\label{poster-abstract-page-talk-247}
\subsection*{Poster 26. Improving Localization of Far-Sidelobe FRBs Detected by CHIME/FRB}
\textbf{Presenter:} Isabella Rivera\\
\textbf{Poster number:} 26\\
\textbf{Fields:} Instrumentation $\cdot$ Transients

Fast radio bursts (FRBs) are millisecond-duration radio transients of extragalactic origin whose physical progenitors and emission mechanisms remain mostly unknown. The Canadian Hydrogen Intensity Mapping Experiment Fast Radio Burst Project (CHIME/FRB) has discovered more than 4500 FRBs and now dominates the detection rate of these sources. Determining whether all FRBs repeat and identifying their host galaxy environments are key steps toward understanding their underlying astrophysical origins. CHIMEs wide field of view and long daily exposure times enable the detection of exceptionally bright FRBs outside the telescopes primary beam, in the far sidelobes. Because the sidelobes have reduced sensitivity, only the brightest events are detected in this region, implying that these sources are typically nearby compared to the broader FRB population. As a result, sidelobe FRBs are promising targets for precise localization and multiwavelength follow-up observations. However, precise localization of sidelobe FRBs is complicated by poorly understood systematics and dearth of calibration sources due to reduced sensitivity arising from the complex beam response in these regions. In this work, I present the development of analysis techniques to improve the localization of far sidelobe FRBs by observing bright transient and steady astronomical calibrators to characterize and correct localization systematics.

\hypertarget{poster-abstract-talk-83}{}\label{poster-abstract-page-talk-83}
\subsection*{Poster 27. Designing a millimeter Balloon-borne VLBI telescope.}
\textbf{Presenter:} Aarchi Shah\\
\textbf{Poster number:} 27\\
\textbf{Field:} Instrumentation

In this poster, I will present the initial design of a balloon-borne 3-mm (88 GHz) Very Long Baseline Interferometry (VLBI) station that would operate as part of a Global VLBI array. VLBI combines simultaneous observations from widely separated telescopes to synthesize an Earth-sized virtual telescope capable of producing extremely high-resolution images of astrophysical objects. This technique was used by the Event Horizon Telescope to produce the first image of the supermassive black hole Sagittarius A* at the center of our galaxy. One of the primary limitations of ground-based VLBI is the atmospheric absorption of high-frequency radio waves. Balloon-borne VLBI stations can address this challenge by operating in the stratosphere, above approximately 99.5\% of Earths atmosphere. Our group has already built the Balloon-borne VLBI Experiment (BVEX) a 22 GHz, 0.9 m radio telescope which will demonstrate that a high-altitude balloon-borne telescope can function as a VLBI station during a stratospheric flight in 2027. This poster presents an initial design for a successor millimeter VLBI telescope. I will outline the engineering requirements and trade-offs involved in developing a sensitive millimeter VLBI telescope suitable for a balloon platform. I will compare several possible receiver designs, including a liquid nitrogen cooled receiver, as well as receivers designed for cryocoolers or liquid helium dewars, discussing their advantages and limitations. I will also discuss how a balloon borne millimeter VLBI telescope can improve global arrays like the Global Millimeter VLBI Array (GMVA) and therefore improve our ability to image the regions around super massive black holes.

\hypertarget{poster-abstract-talk-267}{}\label{poster-abstract-page-talk-267}
\subsection*{Poster 28. Technology Development for Space-based Far-Infrared Spectroscopy and Interferometry}
\textbf{Presenter:} Locke Spencer\\
\textbf{Poster number:} 28\\
\textbf{Field:} Instrumentation

Many fundamental questions leading current astrophysics research require observations with both high resolution spectroscopy and imaging at high spatial resolution, i.e., hyper-spectral imaging, in the Far-Infrared (Far-IR). Topics include: the origin of the universe, the requisite conditions for star and planet formation, and, ultimately, the conditions for the emergence of life. Over half of the energy emitted by the Universe appears in the relatively unexplored Far-IR spectral region, most of which is opaque from ground-based sites necessitating space-borne instrumentation. The European Space Agency (ESA) Herschel space observatory (2009-2013) has provided the first unfettered views of the universe in the Far-IR. These advances in both sensitivity and broad spectral coverage have served to revolutionize our understanding of the formation and evolution of planetary systems, stars, galaxies, and the universe as a whole. Herschel has shown what can be accomplished with a relatively large aperture single dish Far-IR observatory-class space telescope. Herschel observations have also served to highlight the Far-IR gap, i.e., the dramatically poorer angular resolution and sensitivity of Far-IR facilities compared with either side of the spectrum. Many Herschel discoveries are waiting on enhanced spectral and spatial resolution follow-up observations to address the questions raised in this new window on the Universe. Improvements in high spectral resolution and/or sub-arcsecond angular resolution require additional technology development to further exploit this discovery space. We discuss the motivating science and technology development supporting Far-IR instrumentation and components requisite for such advances.

\hypertarget{poster-abstract-talk-168}{}\label{poster-abstract-page-talk-168}
\subsection*{Poster 29. VROOMM: Development Update on a High-Precision Radial Velocity Spectrograph for the Observatoire du Mont-Mégantic}
\textbf{Presenter:} Jonathan St-Antoine\\
\textbf{Poster number:} 29\\
\textbf{Field:} Instrumentation

We present a development update on VROOMM, an optical (360\,nm - 930 nm) high-resolution chelle spec- trograph designed for precision radial velocity measurements at the 1.6-meter telescope of the Observatoire du Mont-Mégantic (OMM) in Qubec, Canada. Since our initial design presentation, the project has progressed through several key milestones. The frontend optical system design is now finalized, with fabrication and assem- bly underway. The backend spectrograph optical design has been completed and optical elements are currently under construction, while the mechanical design of the spectrograph is in progress. A detector engineering cryostat is being built to evaluate the performance of the unique 4K photon-counting EMCCD by e2v for the innovative photon-counting mode that will leverage electron multiplication to achieve 100-200 m/s precision on extremely faint targets (rsdss = 19-20). To accommodate VROOMM as the primary instrument at OMM, comprehensive facility preparations are in progress at the observatory. These include planning for frontend and backend integration, guiding system testing and modifications to meet VROOMMs unprecedented precision requirements, and implementation of a dedicated scheduling system for instrument operations

\hypertarget{poster-abstract-talk-58}{}\label{poster-abstract-page-talk-58}
\subsection*{Poster 30. Using Young Stars to Map the Forces Driving Local Star Formation}
\textbf{Presenter:} Ronan Kerr\\
\textbf{Poster number:} 30\\
\textbf{Field:} Stars and Stellar cluster formation

Young stellar populations provide a fossil record of star formation spanning tens of millions of years, revealing the evolving gas structures that produced them. Through the SPYGLASS program, we are using Gaia astrometry and photometry to expand our knowledge of these populations by computing their demographics and ages, uncovering star formation patterns, and placing them within the broader galactic context. By tracing nearby stellar populations back through the galactic potential to their star formation sites, we find that many young populations share common origins, corroborating and expanding on the cluster families proposed by Swiggum et al. 2024. Two of these sites match the expansion origins of the nearby Local and Orion-Eridanus Bubbles and show velocities oriented outward from the bubble centres, consistent with star formation in the edges of those bubbles. Our analysis also reveals that out of a sample of 16 nearby low-mass stellar populations, only two lack a plausible relationship with a cluster family, suggesting that the influence of expanding bubbles may dominate local star formation. However, we also find evidence for more unusual star formation modes, such as in the Leo Association, which appears to have formed out of an intermediate velocity cloud colliding with material in the galactic midplane. By expanding our record of young stellar populations to increasingly low masses, we are revealing the full range of features and processes that guide local star formation, from expanding bubbles, to spiral overdensities, to colliding gas clouds.

\hypertarget{poster-abstract-talk-39}{}\label{poster-abstract-page-talk-39}
\subsection*{Poster 31. Investigating Binary Stars as a Source of Multiple Populations in Globular Clusters}
\textbf{Presenter:} Zayna Khadour\\
\textbf{Poster number:} 31\\
\textbf{Field:} Stars and Stellar cluster formation

Globular clusters (GCs) are dense stellar systems that show star-to-star chemical abundance variations known as multiple populations (MPs). These stars are enhanced in He, N, and Na and depleted in C and O, consistent with material processed by hot hydrogen burning. Observations, particularly those from the Hubble Space Telescope (HST) show that MPs are ubiquitous in GCs. However, the source of this processed material remains unclear. My research tests whether interacting massive binary stars can contribute to this enrichment. During non-conservative mass transfer, massive binaries eject chemically processed material that matches the observed abundance trends. Furthermore, globular clusters are also dense environments where close stellar encounters can modify binary systems. These interactions can change binary fractions and alter orbital periods, mass ratios, and primary masses. I track how massive binary properties evolve in dynamical cluster models to measure how cluster density reshapes the interacting binary population. I developed a statisti- cal framework to separate true dynamical changes from effects caused by small-number sampling and stellar mass loss. Ultimately, I will compute cluster-informed chemical abundances to test how the environment modifies the predicted enrichment signatures and timescales. This work will show that binary evolution and cluster dynamics must be studied together when evaluating enrichment scenarios in globular clusters.

\hypertarget{poster-abstract-talk-81}{}\label{poster-abstract-page-talk-81}
\subsection*{Poster 32. Advantages of simultaneous monitoring of magnetic chemically peculiar stars with spectropolarimeters ESPaDOnS and SPIRou at the CFHT}
\textbf{Presenter:} Viktor Khalack\\
\textbf{Poster number:} 32\\
\textbf{Field:} Stars and Stellar cluster formation

High-resolution and high signal to noise Stokes I and V spectra of magnetic chemically peculiar (mCP) stars acquired at the Canada-France-Hawaii Telescope (CFHT) with spectropolarimeters SPIRou in infrared and ESPaDOnS in visual spectral diapasons are most suitable to study horizontal and vertical stratification of chemical abundance and to recover a structure of large-scale magnetic field in their stellar atmospheres. Launch of the Wenaokeao instrument in 2027 at the CFHT will allow to use SPIRou and ESPaDOnS simultaneously in the spectropolarimetric mode and accumulate high-quality Stokes I and V spectra for cool mCP targets in the spectral domain from 400nm to 2.4$\mu$m. Application of the Least Square Deconvolution (LSD) method allows to measure the mean longitudinal magnetic field ($<$Bz$>$) with high precision (sometimes few Gauss) even if the field is relatively weak. LSD analysis of line profiles in IR spectra provides even higher precision for $<$Bz$>$ measurements because of much higher wavelengths. These IR line profiles usually are originated at higher levels of stellar atmosphere that increase significantly diapason of optical depths to study vertical stratification of elements abundance and probable variation of $<$Bz$>$ with height. Study of Zeeman and Paschen-Back splitting of line profiles in IR domain for targets with relatively strong magnetic fields will help to improve atomic data (log(gf), configuration of energy levels, etc.) as well. The advantages of using the instrument Wenaokeao for observation of mCP stars are illustrated with examples of several studied targets.

\hypertarget{poster-abstract-talk-65}{}\label{poster-abstract-page-talk-65}
\subsection*{Poster 33. Implementing Binary Yields in Simulations of Star Cluster Formation to Explain the Multiple Populations Problem}
\textbf{Presenter:} Edwin Laverde Villarreal\\
\textbf{Poster number:} 33\\
\textbf{Field:} Stars and Stellar cluster formation

One of the open problems of star cluster formation is that of multiple populations (MPs), where massive clusters have distinct chemical populations, with enhancements in He, N, and Na and depletions in C, and O. Many mechanisms have been proposed to explain this behavior, but a main responsible mechanism is not yet agreed upon. Nguyen \& Sills (2024) investigated the possibility of massive binary stars as one of the necessary ingredients to explain MPs. They found that interacting massive binaries which undergo non-conservative mass transfer produce the right abundances over sensible timescales. To better assess the influence of these sources on the signal of MPs as a cluster and its stars evolve, it is necessary to run simulations that take into account binary evolution and feedback processes, as well as a proper prescription for the formation of binary systems. Thankfully, recent contributions to the star cluster formation code Torch make it an ideal framework with which to test how these binary yields may contribute to MPs. By implementing the yields from binary stars into the stellar feedback routines of Torch, injecting the correct abundances of materials at the expected times, we will be able to determine if sufficient gas mass is retained and enough enrichment is ejected to produce the expected populations. Such simulations will help us understand better the role that binary systems have in the MPs problem, and contribute towards a long-sought explanation for this open problem. In this presentation, I will describe the new implementation of binary yields in Torch and present the results from validation simulations.

\hypertarget{poster-abstract-talk-60}{}\label{poster-abstract-page-talk-60}
\subsection*{Poster 34. Study of stellar pulsations in $\delta$ Scuti star HD68725}
\textbf{Presenter:} Weymè Maiga\\
\textbf{Poster number:} 34\\
\textbf{Field:} Stars and Stellar cluster formation

$\delta$ Scuti stars are variable stars whose rapid variations in brightness provide important information about their internal structure. However, analyzing these variations remains complex due to the presence of instrumental noise and systematic effects in the photometric data provided by TESS. The applied methodology relies on the use of Python-based tools and libraries dedicated to the analysis of astronomical data. The main steps include data preprocessing, and frequency analysis using the Lomb-Scargle periodogram and the Discrete Fourier Transform (DFT). Preliminary results highlight several dominant frequencies associated with the stellar pulsations in HD68725 selected for this study. The obtained results contribute to a better understanding of the variability of $\delta$ Scuti stars and pave the way for more in-depth analyses of their internal structure.

\hypertarget{poster-abstract-talk-166}{}\label{poster-abstract-page-talk-166}
\subsection*{Poster 35. Illuminating the Possible Origins of Six Contact Binaries}
\textbf{Presenter:} Lorne Nelson\\
\textbf{Poster number:} 35\\
\textbf{Field:} Stars and Stellar cluster formation

Contact binaries are ubiquitous, accounting for almost 1\% of the visible stars in our Galaxy. According to the 'standard model', contact binaries contain two stars that are both filling their critical Roche lobes (i.e., sharing the same atmosphere). Unfortunately, it is not clear what pathways led to their formation, and how long they can persist in this phase of contact.

Here we describe a photometric investigation of six contact binary systems of the W Ursae Majoris (W UMa)-type. We extracted the primary and secondary minima from the Transiting Exoplanet Survey Satellite observations of these systems. Using observations previously reported in the literature, linear fits in the O C diagrams were carried out, and new ephemerides were determined. Light curve solutions were obtained using the legacy version of PHysics Of Eclipsing BinariEs (PHOEBE). We found that five of these systems required the presence of cold starspots, thus implying the presence of the O'Connell effect. The absolute parameters for all of these systems were estimated, we show that they are all consistent with the properties of W UMa-type stars.

Using MESA (binary stellar evolution code), we show that the progenitors of these binaries were likely low-mass, short-period, primordial binaries (total mass $\sim$2M$_\odot$). We find that mass transfer can be arbitrarily non-conservative and still produce short-period contact binaries. Other formation scenarios are also briefly examined.

\hypertarget{poster-abstract-talk-213}{}\label{poster-abstract-page-talk-213}
\subsection*{Poster 36. LSST Star Cluster Analysis via CanDIAPL}
\textbf{Presenter:} Sneha Nair\\
\textbf{Poster number:} 36\\
\textbf{Field:} Astrostatistics, Astroinformatics, and AI/Machine Learning

The large volume of imaging data from the Vera C. Rubin Observatory requires new, scalable tools for analysis. The Canadian Data-Intensive Astrophysics Platform (CanDIAPL) provides the cloud-based infrastructure and software needed to support this work, specifically for the Legacy Survey of Space and Time (LSST). I present two software projects designed to improve how we find and study star clusters in Rubin data. The first is an automated spectral energy distribution (SED) fitting pipeline that calculates physical properties like age, mass, and metallicity by comparing cluster photometry to stellar population models. This tool is highly flexible, allowing users to customize initial guesses, adjust plots, define specific priors, and choose between different fitting methods to suit their research needs. The second project is a modular star cluster injection pipeline built to test and benchmark detection algorithms. This framework creates synthetic clusters and adds them to survey images using realistic noise and light profiles. As a plug and play system users can easily swap in different detection methods to see how well they perform across different cluster types. Together, these tools support reproducible and scalable analysis of stellar cluster populations in LSST data products and demonstrate the role of infrastructure like CanDIAPL in enabling collaborative, data intensive astrophysical research.

\hypertarget{poster-abstract-talk-95}{}\label{poster-abstract-page-talk-95}
\subsection*{Poster 37. HERMES: Hierarchical Modelling for Exoplanet Science}
\textbf{Presenter:} Wasi Naqvi\\
\textbf{Poster number:} 37\\
\textbf{Fields:} Astrostatistics, Astroinformatics, and AI/Machine Learning $\cdot$ Exoplanets $\cdot$ Next generation observatories

The Ariel mission will usher in a transformative era in exoplanetary science by characterizing the atmospheres of $\sim$1000 exoplanets. Statistical tools capable of inferring robust population-level trends are needed. We present HERMES (HiERarchical Modelling for Exoplanet Science), a multidimensional hierarchical modelling framework designed to probe population-level constraints across multiple parameters such as planetary mass, planetary metallicity, and stellar metallicity. Starting from the Ariel Mission Candidate Sample (MCS), we inject plausible multidimensional trends and demonstrate successful parameter recovery. Simulated Survey samples are generated such that either survey leverage or sample size is fixed, in the presence of measurement noise and intrinsic scatter. We run our HERMES models on these surveys, and by constraining the a posteriori uncertainty parameters, assess whether survey leverage remains a reliable predictor of trend precision under intrinsic variance. For parameters with significant uncertainties (e.g., stellar metallicity, stellar age, etc.), we further investigate whether survey leverage alone is sufficient to probe multidimensional trends. Moreover, via hierarchical modelling, we seek to more precisely constrain the uncertainty on parameters by distinguishing intrinsic variance from the underlying multidimensional trends. As Ariel prepares to take flight and unlock the era of planetary population studies, through HERMES, we aim to deliver novel statistical tools such that will enable deeper insight into comparative exoplanetary climatology.

\hypertarget{poster-abstract-talk-256}{}\label{poster-abstract-page-talk-256}
\subsection*{Poster 38. Joint posterior inference of pixelated sources and foreground mass distributions in strong gravitational lenses using diffusion and recurrent inference machines}
\textbf{Presenter:} Guillaume Payeur\\
\textbf{Poster number:} 38\\
\textbf{Field:} Astrostatistics, Astroinformatics, and AI/Machine Learning

Modeling galaxy-galaxy strong gravitational lenses to infer the brightness of the lensed source and the mass distribution of the foreground lens is computationally challenging, particularly for high-resolution, high signal-to-noise observations. In these regimes, high-dimensional representations of both the source and the lens mass distribution are necessary to model the data down to the noise level. We present a machine-learning method to generate joint posterior samples of the source galaxy and foreground mass distribution as pixelated images conditioned on observations. The method combines diffusion models with recurrent inference machines and can model realistic gravitational lensing simulations with background and foreground galaxies drawn from cosmological hydrodynamical simulations down to the noise level.

\hypertarget{poster-abstract-talk-59}{}\label{poster-abstract-page-talk-59}
\subsection*{Poster 39. Classifying astrophysical signals and instrumental glitches in LISA data with machine learning}
\textbf{Presenter:} Aryamann Rao\\
\textbf{Poster number:} 39\\
\textbf{Fields:} Astrostatistics, Astroinformatics, and AI/Machine Learning $\cdot$ Gravitational waves $\cdot$ Transients

The Laser Interferometer Space Antenna (LISA) is a space-based gravitational wave detector scheduled for launch in the 2030s by the European Space Agency. Operating in the 0.1 mHz0.1 Hz band, LISA will observe a variety of astrophysical sources including supermassive black hole mergers, extreme mass ratio inspirals, and white dwarf binaries. This will yield new insights into supermassive black hole formation, new tests of general relativity and population studies of white dwarfs.

Like all gravitational wave detectors, LISA data is expected to contain non-astrophysical transients, commonly referred to as glitches. Accurate, rapid identification of these transients is essential, as they can contaminate true gravitational wave signals and impact parameter estimation. I will introduce a machine learning pipeline trained on time-frequency images to automatically detect glitches in LISA data streams. Initial results demonstrate a promising ability to identify overlapping signal and glitch events. I will also discuss challenges encountered during training, including limitations of the current simulation framework, and outline future improvements to enhance the robustness of the model.

\hypertarget{poster-abstract-talk-138}{}\label{poster-abstract-page-talk-138}
\subsection*{Poster 40. Machine learning for 1/f noise correction in JWST NIRISS spectral transit data}
\textbf{Presenter:} Salma Salhi\\
\textbf{Poster number:} 40\\
\textbf{Fields:} Astrostatistics, Astroinformatics, and AI/Machine Learning $\cdot$ Exoplanets

Transit spectroscopy for exoplanet atmospheric characterization is very sensitive to various sources of noise, and this is especially true when using the Single Object Slitless Spectroscopy (SOSS) mode on the NIRISS instrument aboard the JWST. Current methods to deal with 1/f (correlated) noise, such as approximating the 1/f signal as a constant across a column, leave noise residuals that are almost double that of the expected readout noise. Deep learning models could be a way to mitigate this problem, as they have already been shown to be very efficient at a wide variety of tasks in astrophysical data reduction, including denoising. We constructed a score-based generative diffusion model to learn the structure of the noise in dark images, including bad pixels, hot pixels, cosmic rays, and 1/f noise to create a noise model. We then used this noise model as the data likelihood to analyze mock spectral trace observations in a Bayesian framework, allowing us to produce noise-free posterior samples of the pixel values of the underlying trace. This showed an improvement in precision of spectral extractions from the trace of up to 2.9x of when applied to simulated data, allowing us to probe much closer to the photon noise limit than current methods allow. We will now present preliminary results when this method is applied to real data and show how machine learning techniques could be used to improve pre-processing techniques in practice.

\hypertarget{poster-abstract-talk-158}{}\label{poster-abstract-page-talk-158}
\subsection*{Poster 41. Probing Dust Opacity in Orion Protostellar Cores with New ALMA Band 1 Observations}
\textbf{Presenter:} Parisa Nozari\\
\textbf{Poster number:} 41\\
\textbf{Field:} Gas, dust and the ISM

Previous studies of the OMC 2/3 filament in the Orion Molecular Cloud have reported conflicting measurements of the dust opacity index, a key property of dust grains. We present new ALMA Band 1 ($\sim$7 mm) observations toward six protostellar cores in OMC 2/3 to investigate the origin of these discrepancies. Our preliminary results reveal extended emission structures toward five of the six cores that were undetected at shorter wavelengths with ALMA (1.52.1 mm) and NOEMA (2.73.3 mm), along with widespread molecular line contamination across the band, highlighting the chemical complexity of these environments. We combine the new Band 1 data with our previous NOEMA and ALMA observations to construct full spectral energy distributions (SEDs) for each core. Unlike earlier analyses that assumed a characteristic temperature of 25 K, we directly measure core temperatures using rotational diagram analysis of methanol (CH$_3$OH) and propyne (CH$_3$CCH) transitions in the NOEMA 3 mm data. Incorporating these temperature constraints reduces degeneracies in the SED fitting and enables us to better disentangle non-thermal emission, such as free-free contamination, from intrinsic variations in dust opacity. Our initial results show that several cores exhibit consistent opacity indices between the 3 mm NOEMA and 7 mm ALMA Band 1 data when the measured temperatures are adopted. However, one protostellar core shows a clear discrepancy in the Band 1 slope relative to the other wavelengths, suggesting that a different physical emission mechanism may contribute at longer wavelengths. These findings have important implications for dust grain evolution and for improving mass estimates derived from thermal dust emission in Orion protostellar cores.

\hypertarget{poster-abstract-talk-175}{}\label{poster-abstract-page-talk-175}
\subsection*{Poster 42. Modeling Periodic Methanol Maser Flares with MaxwellBloch Equations}
\textbf{Presenter:} Toktam Rashidi Shahri\\
\textbf{Poster number:} 42\\
\textbf{Field:} Gas, dust and the ISM

Periodic variability in astrophysical masers is observed in several high-mass star-forming regions, yet the physical mechanisms responsible for the diverse flare profiles and timescales remain poorly understood. We analyze long-term monitoring data of the 6.7\,GHz methanol maser in several periodic sources, including G9.62+0.20E and G22.356+0.066, confirming known periodicities and identifying additional cycles in multiple velocity components. The observed flares can be reproduced within a unified framework based on the MaxwellBloch equations, in which the masing regions operate in the fast transient superradiance regime driven by periodic pump variations. Modeling these events allows us to constrain key physical parameters of the masing gas and investigate how variations in pumping conditions and inverted column density shape the observed variability. Extending this analysis to a larger sample of periodic maser sources will provide new insight into the physical environments and pumping mechanisms responsible for periodic maser flaring.

\hypertarget{poster-abstract-talk-258}{}\label{poster-abstract-page-talk-258}
\subsection*{Poster 43. A Sharper View of Neutral Gas in Nearby Galaxies}
\textbf{Presenter:} Erik Rosolowsky\\
\textbf{Poster number:} 43\\
\textbf{Fields:} Gas, dust and the ISM $\cdot$ Stars and Stellar cluster formation

Star formation is a driver of galaxy evolution and the process of star formation out of neutral gas is understood at a general level. However, theoretical predictions for star formation efficiency do not match with observations across extragalactic observations at $\sim$100\,pc scales. Progress toward understanding the evolution of star forming regions in galaxies requires a higher resolution view of the gas and stellar feedback to resolve the detailed physics of star formation. Here, we present new mapping campaigns of molecular and atomic gas across the Local Group and other nearby galaxies ($<$5 Mpc) made with the VLA, MeerKAT and ALMA. These maps show we are able to expand a joint view of atomic and molecular gas to nearby disk galaxies (M31, M33, NGC 253 and M83) to resolve the connection between atomic and molecular gas with a focus on star forming regions. These maps show that the atomic and molecular gas are strongly kinematically linked, providing further evidence for short-lived dynamic molecular clouds.

\hypertarget{poster-abstract-talk-119}{}\label{poster-abstract-page-talk-119}
\subsection*{Poster 44. Multi-Component Fitting and Dense Gas Motions in Giant Molecular Clouds}
\textbf{Presenter:} Sarah Sadavoy\\
\textbf{Poster number:} 44\\
\textbf{Fields:} Gas, dust and the ISM $\cdot$ Stars and Stellar cluster formation

The KEYSTONE survey observed ammonia (NH3) gas in 11 giant molecular clouds (GMCs) using the Green Bank Telescope. Initial studies of these data performed a single velocity component fit, but these GMCs have complex velocity structures that include hub-filament systems, overlapping clumps, and feedback from massive stars. Here, we present results from multi-velocity component fits to the NH3 observations that also include radiative transfer. We find that 15\% of the NH3 gas requires two or more velocity components, and that these regions are primarily coincident with the most dense and active zones of the GMCs. By fitting multiple components, the NH3 linewidths decrease by a factor of 1.5, resulting in lower contributions of turbulence to cloud support. We also find variations between the clouds, indicating that local influences such as infall, outflow, and massive star feedback are impacting the dense gas.

\hypertarget{poster-abstract-talk-107}{}\label{poster-abstract-page-talk-107}
\subsection*{Poster 45. Hunting Down Fugitives: Observational Constraints on the Cluster Ejection Mechanisms of X-ray Binaries in Spiral Galaxies}
\textbf{Presenter:} Dev Khullar\\
\textbf{Poster number:} 45\\
\textbf{Fields:} Compact Objects $\cdot$ High Energy Astrophysics $\cdot$ Stars and Stellar cluster formation

X-ray binaries are astrophysical systems formed by a compact object accreting material from a luminous donor star. These systems serve as some of the brightest X-ray lighthouses in the Universe, offering a way to probe the population of both stellar-mass black holes and neutron stars. Most X-ray binaries are believed to have formed in the dense cores of star clusters, however, observations of spiral galaxies show that the majority are found within the galactic disk, away from their parent clusters. Current models suggest that these X-ray binaries may be ejected into the disk either due to dynamical interactions or due to the natal kick received from the supernova explosion that gives birth to the compact object. While this relationship is well researched within elliptical galaxies, the complex morphologies, ongoing star formation, and the significant dust content make it extremely challenging to study this X-ray binary-cluster connection in spiral galaxies. In this talk, we present a new technique to constrain the formation and evolutionary history of X-ray binaries in spiral galaxies. In particular, we use a unique combination of multiwavelength observations with the Chandra X-ray Observatory, the Hubble Space Telescope, and the James Webb Space Telescope, to compile a comprehensive catalog of X-ray binaries, tracing them back to their parent clusters, and setting constraints on their ejection scenarios.

\hypertarget{poster-abstract-talk-226}{}\label{poster-abstract-page-talk-226}
\subsection*{Poster 46. How To Design a Deep Daily Pulsar Search - CHAMPSS Technical Overview}
\textbf{Presenter:} Lars Künkel\\
\textbf{Poster number:} 46\\
\textbf{Fields:} Compact Objects $\cdot$ Transients

The recently launched CHIME All-Sky Multiday Pulsar Stacking Search (CHAMPSS) is a pulsar survey using the Canadian Hydrogen Intensity Mapping Experiment (CHIME) operating in the 400800 MHz frequency band. CHIME is a drift telescope which is able to observe the full Northern sky daily which presents a unique opportunity for designing a pulsar search.

Most pulsar surveys are limited to only a small number of observations per sky position and their overall sensitivity is usually constrained by the sensitivity of individual observations. CHAMPSS allows us to tackle these limitation in two ways. Firstly the data derived from the CHIME/FRB experiment allows us to observe about a quarter of the sky with daily cadence, with this limitation only arising from our current computing resources. Secondly we developed an analysis pipeline which uses incoherent summing of power spectra to slowly build up our sensitivity across time.

I will present our survey pipeline and the current status of the project. To date we were able to detect more than 25 new pulsars, which resulted both from our daily observations and our stacked power spectra from multiple days.

\hypertarget{poster-abstract-talk-324}{}\label{poster-abstract-page-talk-324}
\subsection*{Poster 47. MeV gamma-ray observational signatures of individual heavy isotopes in kilonovae}
\textbf{Presenter:} Maude Larivière\\
\textbf{Poster number:} 47\\
\textbf{Fields:} Compact Objects $\cdot$ Transients

LIGO's detection of gravitational waves has pointed to kilonovae as crucial events to examine in order to unravel the details of the processes responsible for the synthesis of the heavy elements in the Universe, such as the rapid (r) neutron capture process. But details such as the nucleosynthesis reach of such events as well as isotopic composition of kilonovae are still not well understood. The emission spectra of MeV gamma rays from the decays of unstable nuclei could lead to strong insights if individual lines can serve to identify single isotopes. With recent LIGO observations suggesting that neutron star mergers are rarer than previously thought and the launch of COSI in 2027, precise predictions of possible signatures of individual isotopes have become more pressing. In this talk, I will present an exhaustive list of MeV signals associated with individual isotopes that could be used to probe kilonovae in real time and as remnants ($>$ 10 years post-merger). I will also discuss how these calculations are sensitive to nuclear uncertainties and astrophysical variations.

\hypertarget{poster-abstract-talk-191}{}\label{poster-abstract-page-talk-191}
\subsection*{Poster 48. VLBI Observations of SN 2012au Reveal a Compact Radio Source a Decade Post Explosion}
\textbf{Presenter:} Mattias Lazda\\
\textbf{Poster number:} 48\\
\textbf{Fields:} Compact Objects $\cdot$ Transients

Three leading models have been put forth to justify the observed radio re-brightening associated with stripped-envelope supernovae (SESNe) years post explosion: radiation from an emerging pulsar wind nebula (PWN), shock interaction with a dense circumstellar medium, or emission from an off-axis, relativistic jet. SN 2012au is a particularly intriguing SESN in this regard as observations $>$6 years post-explosion have shown both (i) optical emission features consistent with a young PWN and (ii) a radio re-brightening. In this talk, I will present the results of an observational campaign to image SN 2012au using very-long-baseline-interferometry (VLBI) 8 to 13 years post-explosion. We find that our VLBI measurements can be readily explained by a $\sim$decade-old PWN, consistent with independent inferences based on late-time optical spectroscopy of SN 2012au. I will close by discussing the implications of our observations on the properties of the central neutron star, as well as how future observations can be leveraged to gain novel insight into the formation properties of, what is likely, a decade-old pulsar.

\hypertarget{poster-abstract-talk-309}{}\label{poster-abstract-page-talk-309}
\subsection*{Poster 49. Star-Disk Collisions in Quasi-Periodic Eruptions: Global Simulations}
\textbf{Presenter:} Stephanie Lee\\
\textbf{Poster number:} 49\\
\textbf{Fields:} Compact Objects $\cdot$ High Energy Astrophysics $\cdot$ Supermassive black hole $\cdot$ Transients

Quasi-periodic eruptions (QPEs) are X-ray flashes observed near galactic nuclei whose mechanisms remain mysterious. As QPEs occur near the centers of galaxies, understanding their progenitors allows us to probe the extreme physics of the environment around supermassive black holes (SMBHs). We investigate the possibility of QPEs being caused by main sequence companions of SMBHs interacting with the remnant disk of a stellar tidal disruption event. Using the 3D moving mesh hydrodynamics code Arepo, we present the first global simulations of a non-rotating, solar-mass star interacting with a dense, thin disk around a one million solar mass SMBH at several times its own tidal disruption radius. We observe repeated collisions between the stripped stellar material and disk, producing a plume of ejecta and rapidly heating the colliding material, and a complicated interaction between the star, the stream, and the disk disturbance left behind by previous collisions.

\hypertarget{poster-abstract-talk-104}{}\label{poster-abstract-page-talk-104}
\subsection*{Poster 50. CASCA Postdoctoral Committee 2025/2026 Activities}
\textbf{Presenter:} Emily Deibert\\
\textbf{Poster number:} 50\\
\textbf{Fields:} Education and Public Outreach $\cdot$ Equity, Diversity and Inclusion

The CASCA postdoctoral committee serves the interests of Canadian postdoctoral researchers in Canada, including both CASCA members and non-members. In this poster, we share an update on the committee activities from the 2025/2026 academic year. This includes updates on the recently revitalized CANVAS seminar series, our new postdoc-grad student mentorship program, and details from the previous year's GAAP lecture series. We will also share ways for postdocs to get involved in the committee going forward.

\hypertarget{poster-abstract-talk-115}{}\label{poster-abstract-page-talk-115}
\subsection*{Poster 51. Star Finder: Expanding Remote Observatory Access and Accessible Astronomy Curriculum}
\textbf{Presenter:} Tiffany Fields\\
\textbf{Poster number:} 51\\
\textbf{Field:} Education and Public Outreach

We introduce Star Finder - The David Lane Astronomy Outreach Program based at Saint Mary's University and the Burke-Gaffney Observatory. The program supports science educators in bringing astronomy-focused curriculum into high school classrooms through free, ready-to-use lesson plans, hands-on activities, and remote telescope access. Lesson plans link directly to the Nova Scotia Grade 9 Space Exploration unit, and connect with Math 9, Astronomy 12, and other grade levels.

Through Star Finder, teachers and students can request their own astronomical observations, analyze real archival data, and learn through inquiry-driven projects on topics including lunar exploration, exoplanet detection, galaxy imaging, and star clusters. The program aims to reduce barriers to astronomy education by providing accessible curriculum-based resources, remote observatory access, and ongoing virtual support to teachers.

An inaugural two-day teacher workshop in August 2025 held at Saint Mary's University brought science teachers from across Nova Scotia together for hands-on activities, demonstrations of Star Finder projects, presentations by professional astronomers, and evening observing sessions. Building on this foundation, we continue to connect with science educators through professional development workshops across the province and promoting national use of Star Finder resources and the Burke-Gaffney Observatory.

In this talk, we present program goals, outcomes, and current initiatives for expanding across Nova Scotia and Canada. We encourage collaboration from educators and outreach coordinators to expand equitable access to astronomy education.

\hypertarget{poster-abstract-talk-334}{}\label{poster-abstract-page-talk-334}
\subsection*{Poster 52. Teaching the thrill of discovery with real astronomical data: From Interactive Notebooks to AI-Assisted Exploration}
\textbf{Presenter:} Stéphanie Juneau\\
\textbf{Poster number:} 52\\
\textbf{Field:} Education and Public Outreach

Modern astronomy increasingly relies on large datasets, coding, and machine learning, making data literacy and computational skills essential for the next generation of scientists. Education and public outreach programs can play an important role in introducing these concepts early while connecting students to authentic scientific data and workflows. In this talk, I will highlight several activities that use real astronomical datasets and interactive coding environments to engage students from middle school through graduate levels. One example is the NSF NOIRLab Teen Astronomy Cafe program, where short sessions combine a brief introduction to an astrophysical concept, a hands-on demonstration, and an interactive Python notebook that students can modify and execute. These notebooks are hosted on the Astro Data Lab and the Google Collab platforms, both of which provide cloud-based environments with pre-installed software, allowing students to immediately explore data without the barrier of local setup.

Astro Data Lab also hosts more advanced tutorials, including machine learning and AI notebooks used in summer schools and had been used in both undergraduate and graduate university courses. In addition to lowering technical barriers, the platform provides access to large astronomical datasets, enabling students to work directly with the same data used in professional research. Partly building on these efforts, new initiatives within the NSFSimons CosmicAI Institute aim to expand the development and integration of AI tools into education and research. I will briefly discuss these developments and outline our vision for using modern data platforms to accelerate both learning and discovery.

\hypertarget{poster-abstract-talk-282}{}\label{poster-abstract-page-talk-282}
\subsection*{Poster 53. Active Galactic Nuclei Survey Simulation Pipeline for CASTOR}
\textbf{Presenter:} Nicholas Harries\\
\textbf{Poster number:} 53\\
\textbf{Fields:} Supermassive black hole $\cdot$ Transients

Active galactic nuclei (AGN) are supermassive black hole systems located at the centres of massive galaxies. Accretion onto the central black hole releases tremendous energy from X-ray to radio. Although AGN have been studied for decades, our understanding of how such giant black holes grow over cosmic time remains unclear. AGN vary on many timescales. A change in the accretion-disk continuum travels to the surrounding broad-line gas (orbiting close to the black hole) that then responds with a finite time delay. Reverberation mapping (RM; a time-domain technique) converts this time delay into a distance measurement, probing AGN size scales. AGN sizes and broad emission-line widths from spectra yield a virial estimate of the black hole mass the quantity of interest to study black hole growth. The Cosmological Advanced Survey Telescope for Optical and Ultraviolet Research (CASTOR) is a proposed Canada-led UV/optical telescope that will conduct wide-field imaging and spectroscopic surveys. We are designing an AGN RM Legacy Survey to measure black hole masses for a 1000 AGN with redshifts up to $\sim$2.3 to understand supermassive black hole growth over cosmic time a key science goal of CASTOR. My project involves developing a pipeline to simulate RM surveys and extract continuum time lags (accretion-disk RM). We simulate multi-epoch CASTOR scenes containing AGN and propagate these data through the pipeline to ultimately recover cross-correlation time lags. The AGN survey simulation will help in planning large-scale variability surveys with CASTOR by tuning survey observables to determine optimal observing strategies for specific science goals. The pipeline framework is adaptable beyond CASTOR, crucial for crafting ambitious surveys with next-generation facilities.

\hypertarget{poster-abstract-talk-295}{}\label{poster-abstract-page-talk-295}
\subsection*{Poster 54. Preparing for CASTOR Active Galactic Nuclei Science with a CFHT Time-Domain Community Survey}
\textbf{Presenter:} Viraja Khatu\\
\textbf{Poster number:} 54\\
\textbf{Field:} Supermassive black hole

Every massive galaxy hosts a central supermassive black hole. Some black hole systems powered by active accretion active galactic nuclei (AGN) release significant energy over a broad range of the electromagnetic spectrum. AGN are excellent laboratories to study how supermassive black holes grow over cosmic time a long-standing question in AGN science. In AGN, high-energy ionizing photons from the continuum are reprocessed to lower energies or reemitted as emission-line photons in the nearby broad-line region, with a time lag. The time-domain technique of reverberation mapping (RM) translates this time lag into a distance, yielding AGN accretion-disk and BLR sizes. Combined with widths of broad emission lines from AGN spectra, we can measure the central black hole masses in these systems. Deriving black hole masses for AGN over a wide redshift range is crucial to understand supermassive black holes growth also a key goal of the Cosmological Advanced Survey Telescope for Optical and ultraviolet Research (CASTOR; a proposed Canadian flagship, UV satellite mission). With a CASTOR AGN RM Legacy Survey, we aim to obtain black hole masses for a 1000 AGN in the deep drilling fields (DDFs) of the Vera C. Rubin Observatorys Legacy Survey of Space and Time. To prepare for AGN RM science with CASTOR, we propose a multi-year AGN monitoring campaign with the Canada-France-Hawaii Telescope (CFHT) to survey three DDFs in the MegaCam/u, g, r, and Halpha bands under CFHTs upcoming Community Survey (open for broader community participation). Our program will take advantage of the excellent seeing on Maunakea, MegaCam sensitivity in u, and complimentary cadences to LSST. I will present our CFHT AGN monitoring program, highlighting its goals and survey planning strategies.

\hypertarget{poster-abstract-talk-141}{}\label{poster-abstract-page-talk-141}
\subsection*{Poster 55. Optical and Infrared Signatures of Collapsar Disk Outflows}
\textbf{Presenter:} Nicolás Concha-Marroquin\\
\textbf{Poster number:} 55\\
\textbf{Fields:} High Energy Astrophysics $\cdot$ Transients

The leading framework for explaining long-duration gamma-ray bursts and their associated supernovae is the collapsar model: the failed core collapse of a rapidly rotating massive star that forms an accretion disk around a newly formed black hole. This system powers both a relativistic jet and a slower disk wind. Observations suggest these supernovae are too energetic to be produced by the jet alone, pointing to the disk wind as the primary engine powering the supernova emission. I will present synthetic light curves and spectra obtained by post-processing ejecta from global viscous hydrodynamic simulations of collapsars evolved to homologous expansion. Using Monte-Carlo radiative transfer, we directly compare the optical and infrared signatures of these disk outflows with observed transients. Our suite of models, which varies progenitor and disk parameters while excluding a jet, predicts a diverse range of signatures consistent with observed events, ranging from standard Type Ic to superluminous supernovae.

\hypertarget{poster-abstract-talk-192}{}\label{poster-abstract-page-talk-192}
\subsection*{Poster 56. Stellar shock emergence in compact aspherical supernovae: the boundary between relativistic and oblique branches}
\textbf{Presenter:} Christopher Matzner\\
\textbf{Poster number:} 56\\
\textbf{Fields:} High Energy Astrophysics $\cdot$ Transients

In stellar explosions, shock emergence is crucial for the creation of high-energy transients, as it produces both a photon flash and the highest-velocity ejecta. An exception are the oblique shocks in compact, aspherical supernovae, which fail to emit light or fast ejecta. But becasue the same explosions produce relativistic flows, a key question is how a relativistic outflow transitions to a oblique one. I will diagnose this question by choosing some useful limits that admit self-similar solutions.

\hypertarget{poster-abstract-talk-44}{}\label{poster-abstract-page-talk-44}
\subsection*{Poster 57. Probing Time-Dependent Physics with Phase-Folding CMB Maps}
\textbf{Presenter:} Yilun Guan\\
\textbf{Poster number:} 57\\
\textbf{Fields:} Cosmology and Early Universe $\cdot$ Gravitational waves

Time-resolved observations of the Cosmic Microwave Background (CMB) and the millimeter sky offer a powerful probe of time-dependent cosmological signals, such as a stochastic gravitational wave background passing through Earth, which imprints a time-varying deflection on the CMB, and time-dependent cosmic birefringence, which induces an oscillating polarization rotation. However, analyses based on time-division maps are fundamentally limited by the mapmaking cadence, restricting sensitivity to frequencies below $10^{-5}$ Hz. In this work, we present a phase-folding mapmaking framework to enable targeted searches of such periodic cosmological signals with frequencies up to the detector sampling rate. We demonstrate the power of this framework with two cosmological applications: (1) constraining a stochastic gravitational wave background via its time-dependent lensing signature, and (2) the search for an oscillating polarization rotation from axion-like particles. We show that this technique can transform CMB experiments into broadband probes of the oscillating sky, extending their constraining power from the microhertz regime up to $\sim$100 Hz, an expansion of over seven orders of magnitude in frequency. This provides a new observational window, complementary to other approaches by probing under-explored frequency ranges for gravitational waves and axion-like particles.

\hypertarget{poster-abstract-talk-320}{}\label{poster-abstract-page-talk-320}
\subsection*{Poster 58. Stacking Methods for Line Intensity Mapping Experiments with Bayesian Statistics}
\textbf{Presenter:} Patrick Horlaville\\
\textbf{Poster number:} 58\\
\textbf{Field:} Cosmology and Early Universe

Through the past decades, cosmology has become a precision science thanks to the advent and the triumphs of CMB and galaxy survey experiments, probing with great detail the Universe 400,000 years and 13.8 billion years after the Big Bang, respectively. However, the evolution of structures and galaxies in between these two epochs remains elusive, and Line Intensity Mapping (LIM) experiments are new observatories that aim to map large cosmological volumes during that time period by looking at the unresolved emission from specific spectral lines. LIM has its own set of observational challenges (including the faintness of the targeted signals and the presence of foreground contaminants), such that clever statistical methods in the analysis of its data are invaluable towards a signal detection. Examples of such techniques includes stacking the LIM data around the positions of matter tracers, such as galaxies or quasars drawn from an object catalog. This emerging technique has already proven both its potential and difficulties, with most recent results showing both detections and non-detections depending on the chosen LIM data, objects catalog and analysis methods. Here, I will present the first results of a new framework we have developped for stacking objects in LIM data, which utilizes bayesian statistics to improve on the current landscape of stacking statistics for LIM. This work not only allows a more rigorous treatment of error propagation, but is also more inclusive in the number of objects to include in the stacking analysis, which helps to build up the signal-to-noise ratio required for a detection. My presentation will center on the pipeline we have developed as well as predictions for current/upcoming LIM experiments which might benefit from this method.

\hypertarget{poster-abstract-talk-201}{}\label{poster-abstract-page-talk-201}
\subsection*{Poster 59. Constraining the evolution of dynamic environments surrounding actively repeating fast radio bursts}
\textbf{Presenter:} Nicole Mulyk\\
\textbf{Poster number:} 59\\
\textbf{Field:} Transients

Fast radio bursts (FRBs) are short-duration (ms), bright (Jy-kJy) bursts of radio emission, generally originating from cosmological distances. While most FRBs appear to be one-off events, $\sim$3\% of FRBs have been observed to repeat. One of the most vital unanswered questions regarding FRB origins is whether all FRBs are capable of repeating. The Canadian Hydrogen Intensity Mapping Experiment (CHIME) radio telescope has detected thousands of FRBs, monitoring a large number of repeaters over multi-year baselines with a daily cadence. Radio spectra of FRBs are uniquely able to probe the properties of the plasma in the FRBs' immediate vicinity. CHIME/FRB's high time resolution (2.56\,$\mu$s) and broadband, low-frequency coverage make CHIME/FRB ideal for constraining propagation effects and monitoring them over time. Recent works have discovered secular evolution of the magnetized plasma surrounding a few peculiar repeating FRBs, linking them to young stellar remnants of massive stars. Our project will build on this work by measuring changes in the magnetoionic plasma over a timescale of months to years, for a diverse sample of repeaters. By examining changes in the magnetoionic plasma for a larger population of repeaters, we aim to establish boundaries on the potential progenitor or environment for the general population of repeating FRBs.

\hypertarget{poster-abstract-talk-338}{}\label{poster-abstract-page-talk-338}
\subsection*{Poster 60. A complete census of Northern Milky Way satellites and streams using UNIONS}
\textbf{Presenter:} Andrew Li\\
\textbf{Poster number:} 60\\
\textbf{Field:} Local Group

Milky Way (MW) satellites are expected to tidally interact with the gravitational field of the Galaxy. This can result in extended tidal features that encode valuable information about their interaction history and consequently about the accretion history of the Galaxy. However, identifying these extremely faint features proves to be difficult due to their large spatial extent, very low stellar density and high degree of contamination from foreground Milky Way stars and even faint background galaxies. The Ultraviolet Near-Infrared Optical Northern Survey (UNIONS) provides a new avenue of exploration for these features in the northern hemisphere. Roughly equivalent in depth to year 1 of LSST, UNIONS consists of ugriz photometry over 6250 degree squared of the Northern sky. It potentially offers the widest, deepest and most complete ground-based view of our local universe until it is eventually superseded (albeit in the south) by LSST. Here we investigate the faint outskirts of every known MW satellite and stellar stream in the footprint, describe and validate our stellar membership selection methodology, and present our latest results searching for faint tidal features while conducting a comprehensive characterisation of all known tidal structures in the northern Galactic halo.

\hypertarget{poster-abstract-talk-257}{}\label{poster-abstract-page-talk-257}
\subsection*{Poster 61. Unveiling the Origin of a Historical Supernova}
\textbf{Presenter:} Amilia Petryk\\
\textbf{Poster number:} 61\\
\textbf{Field:} Supernova remnant

The nature of historical supernovae are often not well-known, but supernova remnants and their environments can serve as powerful tools for probing their progenitor systems. RCW 86 is an interesting supernova remnant, thought to be associated with the historical supernova SN 185 as recorded by ancient Chinese astronomers. If so, that would make RCW 86 the oldest supernova remnant with surviving records. To reconcile its age with its large size, astronomers have determined it likely exploded into a low-density cavity. This theory allows for the rapid expansion necessary for RCW 86 to have gotten as large as we see it today. However, this leads to questions about the nature of the progenitor system, including its temperature and luminosity. To answer these questions, longslit spectroscopic data were gathered in 2025 using the 8.1-meter Gemini South telescope, both slightly ahead of the expanding shock front as well as on the shock itself. The data spanned 400-750\,nm across two different observations, allowing for measurement of He I lines, He II lines, and the Balmer series. The data were reduced via the Gemini DRAGONS pipeline, with the goal of measuring or constraining the ionized H and He fluxes. From He/H flux ratios alongside CLOUDY and stripped star models, estimates of the temperature and luminosity of the progenitor object will be determined.

\programdayheader{Thursday}
\daypalettebanner{Thursday}

\hypertarget{poster-abstract-talk-80}{}\label{poster-abstract-page-talk-80}
\subsection*{Poster 1. Constraining Multi-Generational Planet Formation in Pressure Bumps: The Role of Viscosity, Fragmentation, and Gap Morphology}
\textbf{Presenter:} Ramadilen Moodelly\\
\textbf{Poster number:} 1\\
\textbf{Fields:} Exoplanets $\cdot$ Gas, dust and the ISM $\cdot$ Solar System

Current models of protoplanetary disk evolution suggest that dust trapping within pressure bumps is a necessary precursor to planet formation via the Streaming Instability. To understand the physical boundaries of this process, we investigate the parameter space that governs planet formation efficiency using a dynamically coupled DustPy and SyMBAp framework. By introducing an artificial Gaussian gap to trigger an initial pressure bump, we track the continuous evolution from sub-micron dust grains to massive planetary bodies. We systematically vary the disk alpha-viscosity, fragmentation velocity thresholds, and the amplitude and width of the initial gap to identify the regimes that either enable or suppress planetary growth. Specifically, we examine the threshold conditions under which initial massive planets carve deep gaps that efficiently dam the inward flow of dust to create secondary dust reservoirs. By mapping the viability of these newly induced dust traps, we constrain the parameters necessary to achieve multi-generational planetary growth and explore the physical limits of spawning more than two distinct generations of planets within a single disk.

\hypertarget{poster-abstract-talk-221}{}\label{poster-abstract-page-talk-221}
\subsection*{Poster 2. High Resolution Emission Phase Curve of the Ultra Hot Jupiter WASP-33 b}
\textbf{Presenter:} Georgia Mraz\\
\textbf{Poster number:} 2\\
\textbf{Field:} Exoplanets

Ultra-hot Jupiters (UHJs) are the most extremely irradiated exoplanets, receiving up to 100 times more stellar flux than typical hot Jupiters and reaching dayside temperatures exceeding 2200 K. These tidally locked planets have short orbital periods and exhibit strong daynight temperature contrasts often exceeding 1000 K. Intense dayside irradiation drives molecular dissociation and thermal inversions, while the cooler nightside allows recombination and non-inverted temperature profiles. Although global atmospheric views traditionally require space-based phase curves, recent work shows that ground-based near-infrared high-resolution spectroscopy (HRS) can recover similar information. WASP-33 b is the only exoplanet with confirmed detections of nightside thermal emission using SPIRou at the CanadaFranceHawaii Telescope, making it a key benchmark for HRS phase-curve studies. We obtained new SPIRou observations at quadrature phases (0.130.37 and 0.680.85), totaling 16 hours. Combined with archival dayside and nightside data, these observations nearly complete the planets orbital phase curve. This is the first time all emission phases of an exoplanet have been assembled using high-resolution spectroscopy from a single instrument, providing a uniform global view of the atmosphere. Preliminary analysis near phase 0.25 shows stable systematics and tentative CO emission, consistent with the hot dayside becoming visible at quadrature and with predictions from global circulation models. Our ongoing analysis aims to map the longitudinal distribution of CO and probe atmospheric winds via Doppler shifts. This work demonstrates the growing potential of high-resolution spectroscopic phase curves for atmospheric characterization in the era of ELTs.

\hypertarget{poster-abstract-talk-332}{}\label{poster-abstract-page-talk-332}
\subsection*{Poster 3. Low-Cost InGaAs Short-Wavelength Infrared Detector for Astronomy: On-Sky Validation with a 0.5 m Telescope}
\textbf{Presenter:} Matthew Muir\\
\textbf{Poster number:} 3\\
\textbf{Field:} Exoplanets

Recent advancements in Indium Gallium Arsenide (InGaAs) detector technology offer a high-performance, cost-effective alternative to traditional Mercury Cadmium Telluride (HgCdTe) sensors for short-wave infrared (SWIR) astronomy. We test the performance of a commercially available, thermoelectrically cooled InGaAs camera - a flight-ready" Owl 1280 from Raptor Photonics - for astronomical observations. Our primary objective is to validate the cameras performance through a tiered testing approach, supporting its eventual implementation on the POET (Photometric Observations of Exoplanet Transits) space mission. The validation process involves two phases. (1) Ground-based assessment: we benchmarked on-sky sensitivity and photometric stability using the Colibri Telescope Array across the operable 400\,nm - 1700\,nm spectrum of the camera. (2) Near-space testing: we plan to deploy the camera on a high-altitude stratospheric balloon to evaluate sensor performance in a near-space environment. We present initial results from our ground-based tests, which demonstrate the viability of low-cost InGaAs SWIR cameras for science on 0.5 m telescopes.

\hypertarget{poster-abstract-talk-306}{}\label{poster-abstract-page-talk-306}
\subsection*{Poster 4. Characterizing Exoplanet Host Stars with SATLAS: A Case Study of 55 Cancri}
\textbf{Presenter:} Tahere Parto\\
\textbf{Poster number:} 4\\
\textbf{Field:} Exoplanets

Characterizing exoplanet host stars, particularly their radii, is essential for accurately determining planetary properties. Since the radius of a transiting exoplanet is derived relative to the radius of its host star, uncertainties in stellar parameters directly propagate into uncertainties in planetary measurements. In this work, we use the spherical stellar atmosphere code SATLAS to estimate the radii of host stars and compare these estimates with values obtained using other methods. SATLAS is the spherical implementation of the ATLAS stellar atmosphere model and utilizes precomputed model grids along with extensive atomic and molecular line opacity databases. Compared to plane-parallel models, spherical atmospheres provide a more realistic treatment of limb darkening, an important factor in transit depth calculations that directly affect exoplanet radius estimates. As a case study, we compute a synthetic spectrum of the well-studied G-type host star 55 Cancri using SATLAS and SYNTHE. We compare our model predictions with high-resolution observations obtained with ESPaDOnS to evaluate how well spherical atmosphere models reproduce the observed stellar spectrum. We present our preliminary results and discuss their implications for improving host-star characterization and, consequently, exoplanet parameter estimates.

\hypertarget{poster-abstract-talk-88}{}\label{poster-abstract-page-talk-88}
\subsection*{Poster 5. Rotation Periods of 25 L and T Dwarfs from 20-Hour Spitzer Monitoring}
\textbf{Presenter:} Mykola Posternak\\
\textbf{Poster number:} 5\\
\textbf{Field:} Exoplanets

Recent time-resolved photometric and spectroscopic observations of brown dwarfs suggest that the viewing geometry may play a significant role in shaping their observed properties, including infrared colours and variability amplitudes. In this context, we present an update on rotation periods for 25 L and T dwarfs continuously observed for 10 hours with the Spitzer Space Telescope in the 3.6\,$\mu$m and 4.5\,$\mu$m bands. These observations provide some of the longest uninterrupted, high-precision mid-infrared light curves of brown dwarfs to date, enabling robust and precise period determinations. We employed Lomb-Scargle periodogram analysis to identify variable sources and obtain initial period estimates, which were subsequently refined using a Markov Chain Monte Carlo method. Out of the 25 targets, 15 are identified as variable at a 95\% confidence level, with rotation periods ranging from 1 to 7 hours and variability amplitudes between 0.18\% and 2.12\%. These results will be combined with follow-up spectroscopic observations, enabling the determination of spin-axis inclination angles for the targets. Together with 24 previously published inclination measurements, they will constitute a statistically significant and homogeneous sample of brown dwarfs with well-constrained spin-axis inclinations, allowing us to test the correlation between brown dwarf colours and their spin-axis inclinations.

\hypertarget{poster-abstract-talk-253}{}\label{poster-abstract-page-talk-253}
\subsection*{Poster 6. Understanding the Impact of Cold Gas Giants on the Formation of Super-Earths and Sub-Neptunes with Astrometry}
\textbf{Presenter:} Alexandra Rochon\\
\textbf{Poster number:} 6\\
\textbf{Field:} Exoplanets

The upcoming Gaia Data Release 4 (DR4) is expected to detect $\sim$10,000 new exoplanets via the astrometry method, with particular sensitivity to gas giants (GGs) at large orbital separations, a currently undersampled planet population. With masses $>$100x that of the Earth, GGs dynamically impact the distribution of planetary building blocks throughout the protoplanetary disk and are therefore expected to influence the formation and compositions of other planets in the system. We investigate planet formation trends for two of the most common types of exoplanets, super-Earths (SEs), rocky planets up to $\sim$1.6x the size of the Earth, and sub-Neptunes (SNs), larger planets enveloped in a low-density material such as hydrogen, helium, and/or water. Cold GGs carve disk gaps that may hinder the inward drift of icy pebbles, thus inhibiting the formation of massive solid cores that can grow into H/He-rich SNs, favouring the presence of low-mass SEs that form in a volatile-poor environment. Demographics-level radial velocity studies have shown a tentative positive correlation between GGs and small planet occurrence, but limited sample sizes have hindered establishing trends for SEs and SNs separately. In anticipation of DR4 (expected Dec 2026), we explore the utility of a calibrated Gaia DR2, DR3 and Hipparcos astrometric catalog (G23H) and modeling framework to extract the maximum information from existing, non-time-series, astrometric observations. We aim to measure the conditional occurrence of cold GGs in systems with known SEs and/or SNs to test the dominant formation processes of these planets around FGK, and possibly M dwarfs.

\hypertarget{poster-abstract-talk-91}{}\label{poster-abstract-page-talk-91}
\subsection*{Poster 7. Unraveling the Mystery of the Peculiar and Young Hot Jupiter CoRoT-2b I: H2O and CO Detection from Dayside Observations with Gemini-S/IGRINS}
\textbf{Presenter:} Ying Shu\\
\textbf{Poster number:} 7\\
\textbf{Field:} Exoplanets

Observations of the hot Jupiter CoRoT-2b have revealed a peculiar westward offset of its atmospheric hotspot, which may be explained by multiple physical scenarios. The first step toward understanding this phenomenon is to characterize the planets atmospheric composition. In this talk, we present ground-based, high-resolution spectroscopic pre-eclipse observations of CoRoT-2b obtained with the IGRINS spectrograph on Gemini South. Using a cross-correlation analysis, we detect the Doppler-shifted signature of the planets thermal emission with a signal-to-noise ratio of 4.32. Two fully independent data reduction and atmospheric retrieval pipelines confirm the detection and constrain the abundances of H$_2$O and CO in the atmosphere of CoRoT-2b. No significant detections of CH$_4$, CO$_2$, TiO, or VO are found. While the cross-correlation analysis tentatively suggests the presence of HCN and OH, retrieval analysis does not robustly confirm these molecules. The detected H$_2$O and CO features indicate that CoRoT-2bs dayside spectrum is not featureless, contrasting with previously inferred from HST/WFC3 observations, but instead reveals a complex atmospheric structure. Interestingly, we find a lack of significant molecular features at wavelengths shorter than 1.7\,$\mu$m, potentially due to high-altitude absorbers such as H-, clouds, or observational systematics. From our retrieved abundances of CO and H$_2$O, we constrain a super-solar C/O ratio and a slightly sub-solar metallicity. This study provides the first high-resolution constraints on the atmospheric composition of CoRoT-2b. Additional observations with CRIRES+ are analyzed to further probe the planets atmospheric dynamics, although more phase-resolved observations will be required to robustly test the proposed scenarios.

\hypertarget{poster-abstract-talk-46}{}\label{poster-abstract-page-talk-46}
\subsection*{Poster 8. Uncovering ultra-short period rocky worlds in the near infrared}
\textbf{Presenter:} Avidaan Srivastava\\
\textbf{Poster number:} 8\\
\textbf{Field:} Exoplanets

Ultra-short period rocky worlds represent a class of planets for which there exists no analogue in the Solar System, making them targets of high interest from a formation and evolution point of view. Only eleven such worlds orbiting M dwarfs have a characterised mass and radius. TOI-4552b (Mp = 1.87 Me, Rp = 1.10 Re) is the newest of this kind and is the only one to be detected using near-infrared radial velocities, owing to the new NIRPS instrument. Although NIRPS operates simultaneously with HARPS, the radial-velocity detection of TOI-4552b is primarily driven by the near infrared data, highlighting the advantage of NIRPS over visible wavelength spectrographs for planet detection around faint M dwarfs. The combined visible and near infrared coverage however enables a more robust characterisation of stellar activity and refractory element abundances. Using these stellar abundances, detailed interior structure models can be constructed, yielding a well constrained core mass fraction (CMF). TOI-4552b has a marginally high CMF of 0.55, which, coupled with its short orbital period (0.3 days) and high emission spectroscopic metric (ESM = 19.5), makes it an exceptionally valuable target for surface and atmospheric characterisation with JWST.

\hypertarget{poster-abstract-talk-298}{}\label{poster-abstract-page-talk-298}
\subsection*{Poster 9. HAT-P-70b through the Eyes of MAROON-X: Constraining Elemental Abundances of Metals and Insights on Atmosphere Dynamics}
\textbf{Presenter:} Shi Lin Sun\\
\textbf{Poster number:} 9\\
\textbf{Field:} Exoplanets

Ultra-hot Jupiters (UHJs) are exceptional laboratories for studying planetary atmospheres under extreme irradiation conditions. With close-in tidally locked orbits, these planets can have daysides hot enough for metals to be significantly ionized while still maintaining nightsides cold enough for refractory species to potentially condense. We present an analysis of the ultra-hot Jupiter HAT-P-70b taken with the MAROON-X high-resolution spectrograph. Using cross-correlations, we detect 14 neutral and singly ionized species, including Fe I, Fe II, Ti I, Ca I, Ca II, Cr I, Na I, V I, Mn I, Ni I, Mg I, Ba II, O I, and Sr I, with tentative evidence for H I, Co I, and K I. The absorption signals exhibit blueshifts on the order of a few km s$^{-1}$, consistent with day-to-night winds. We further constrain relative abundances with atmospheric retrievals and demonstrate that some inferred elemental abundance ratios depend strongly on modeling assumptions. In particular, we show that a well-mixed retrieval approach neglecting ionization can strongly bias highly ionizable elements such as Ca and Ti. Accounting for the effects of equilibrium chemistry and thermal ionization generally results in inferred elemental abundance ratios that are closer to expectations for a solar-like composition, although not in all cases. Interestingly, we find a distinct nickel enrichment on HAT-P-70b, adding to the growing number of UHJ studies where the Ni abundance is seemingly enhanced. Our results underline the importance of considering physical and chemical atmospheric processes such as ionization when interpreting high-resolution transmission spectra of UHJs.

\hypertarget{poster-abstract-talk-236}{}\label{poster-abstract-page-talk-236}
\subsection*{Poster 10. A BESPOKE Approach to Exploring Galactic Star Formation}
\textbf{Presenter:} Emily Rock\\
\textbf{Poster number:} 10\\
\textbf{Field:} Galaxy

We present BESPOKE high-resolution simulations of NGC 5055 in order to evaluate current star formation models through a direct comparison with the observations. Nearby galaxy observations using the VLA (THINGS), ALMA (e.g. PHANGS) and Integrated Field Units provide high-resolution snapshots of galactic star formation and the ISM. Simulations are required for a detailed study of star formation evolving over time.

Our goal is to improve star formation models and other physical assumptions we make to maximize the predictive power of galaxy simulations, not just locally but out to high redshifts and in extreme environments. NGC 5055 is a well-studied, nearby spiral galaxy. Our simulations are tailored to match NGC 5055's rotation curve, stellar surface density, and gas surface density obtained from The HI Nearby Galaxy Survey (THINGS). We have run 19 different star formation models that bracket the range of star formation parameters in the literature. We found that model choice has a significant impact on the simulated galaxies' evolution. The models vary greatly in not just the rate at which the simulated galaxy consumes its gas to form stars, but also in the structure of the galaxy and whether it is a detailed match for the real thing. This type of close comparison is uniquely enabled by our simulations that target a specific galaxy.

\hypertarget{poster-abstract-talk-72}{}\label{poster-abstract-page-talk-72}
\subsection*{Poster 11. Young star-forming regions in NGC4449 as seen with SITELLE}
\textbf{Presenter:} Anne-Sophie Roy\\
\textbf{Poster number:} 11\\
\textbf{Field:} Galaxy

SIGNALS is studying thousands of HII regions in various galactic environments across more than thirty nearby galaxies, using SITELLE at the CFHT. While emission lines in the three SN filters allow for the detection and study of HII regions, the continuum maps from these filters highlight stellar clusters. A comprehensive catalog of more than 2000 emission regions is presented for the dwarf irregular galaxy NGC 4449. The properties and abundances of the HII regions are studied, with a luminosity function, BPT diagrams and a metallicity distribution. Furthermore, to model and understand the physical processes linking HII regions and star-forming clusters, a catalog of stellar clusters was also built. The physical properties (age, mass, reddening and metallicity) of each cluster are derived from stellar population synthesis methods to fit the spectra obtained with SITELLE. A comparison of the young clusters properties to those of the nearby HII region is also made. This project aims to assess SITELLEs ability to determine the physical properties of star clusters by developing a stellar population synthesis method adapted to its data, and to model the HII regions formed by these very same clusters.

\hypertarget{poster-abstract-talk-139}{}\label{poster-abstract-page-talk-139}
\subsection*{Poster 12. Probing Globular Clusters in the ECDFS using LSST's Data Preview 1}
\textbf{Presenter:} Yashpranav Sairam\\
\textbf{Poster number:} 12\\
\textbf{Fields:} Galaxy $\cdot$ Next generation observatories $\cdot$ Stars and Stellar cluster formation

The Extended Chandra Deep Field South (ECDFS) is a widely studied X-ray extragalactic field hosting a rich population of galaxies and, consequently, active galactic nuclei and quasars. Using Data Preview 1 (DP1) from the Legacy Survey of Space and Time (LSST), which aims to conduct the deepest and widest multi-wavelength ground-based optical survey, we explore the ECDFS, as this deep drilling field is the prime field candidate to identify and characterize the time-variable nature of extragalactic objects. Using cross-matched observations with NED, we explore the classification of objects in the ECDFS with no object label and the average variability of the different object classes. We further explore the use of a fractional variability metric as a diagnostic to evaluate the nature of time-variable analysis in LSST DP1.

\hypertarget{poster-abstract-talk-116}{}\label{poster-abstract-page-talk-116}
\subsection*{Poster 13. Faraday complexity in NGC 891}
\textbf{Presenter:} Jeroen Stil\\
\textbf{Poster number:} 13\\
\textbf{Fields:} Galaxy $\cdot$ Gas, dust and the ISM

NGC 891 has been observed in full Stokes with the Jansky Very Large Array as part of the CHANG-ES survey in C-band, S-band and L-band. We probe the structure of the diffuse magneto-ionic medium of this nearby Milky-Way analogue with Faraday Rotation Measure Synthesis applied to its polarized diffuse synchrotron emission. We find complex Faraday rotation associated with super shells located on the Earth-facing side of the disk.

\hypertarget{poster-abstract-talk-250}{}\label{poster-abstract-page-talk-250}
\subsection*{Poster 14. Cosmology and Astrophysics with 17,000 Clusters from UNIONS}
\textbf{Presenter:} James Taylor\\
\textbf{Poster number:} 14\\
\textbf{Fields:} Galaxy $\cdot$ Galaxy Clusters

As the largest dense, gravitationally bound structures in the Universe, galaxy clusters provide many important tests of astrophysics and cosmology. The current generation of optical surveys are producing cluster samples 10-100 times the size of those previously available. Thus, the precision and scope of cluster-based tests are changing dramatically. In this talk, we will present the UxD catalogue, a cross-match between catalogues derived from the DESI LEGACY imaging survey and UNIONS photometric and weak lensing data. The initial match includes more than 250,000 clusters, while quality cuts reduce this sample to 17,000 with good redshift and membership information. We classify the dynamical state of these clusters using several simple optical metrics. We show that the luminosity function and mean density profile of clusters vary systematically with dynamical state. We calibrate the mean mass and mass-richness scaling relations using UNIONS weak lensing data. Comparing to similaly selected samples in the IllustrisTNG simulation, we estimate the purity, completeness and mass bias of our sample, and show that these also vary with dynamical state. Finally, we discuss ongoing tests of astrophysics and cosmology with this new sample.

\hypertarget{poster-abstract-talk-85}{}\label{poster-abstract-page-talk-85}
\subsection*{Poster 15. Unveiling a 150 kpc Stripping Tail in SpARCS1051, a z = 1.034 Galaxy Cluster}
\textbf{Presenter:} Auriane Thilloy\\
\textbf{Poster number:} 15\\
\textbf{Fields:} Galaxy $\cdot$ Galaxy Clusters $\cdot$ Gas, dust and the ISM $\cdot$ High Energy Astrophysics $\cdot$ Supermassive black hole

We investigate the origin of a $\sim$150 kpc tail observed around the brightest cluster galaxy (BCG) of the high-redshift galaxy cluster SpARCS1051 (z $\sim$ 1.034). Three scenarios are explored. First, we test whether strong cooling of the intracluster medium, similar to that in the extreme cool-core cluster SpARCS1049, could produce this feature. However, the 190 ks Chandra observations show no detection of diffuse X-ray emission associated with a strong cool core, making this explanation unlikely.

Second, we examine whether tidal stripping from nearby cluster members could explain the tail. This hypothesis is motivated by four massive galaxies near the BCG that could contribute stripped material. To assess the cluster dynamics, we perform a Dressler-Shectman test using spectroscopically confirmed members from the GoGreen survey to search for dynamical substructure and mergers. The results show no clear evidence of significant merging activity in the core or of these galaxies merging with the BCG, although an undetected interaction cannot be excluded if a galaxy lies along the line of sight or if the relative motion occurs mainly in the plane of the sky.

Finally, we explore a ram-pressure stripping scenario. Multi-band GoGreen imaging reveals another red and dead galaxy east of the BCG with a significant velocity offset relative to the cluster redshift. We propose that this galaxy is moving rapidly through the intracluster medium and producing the $\sim$150 kpc tail via ram-pressure stripping. The tail appears blue and rich in star-forming gas in the GoGreen observations, suggesting a rare example of ram-pressure stripping in a distant cluster environment, providing new insight into the environmental processes shaping galaxies in high-redshift clusters.

\hypertarget{poster-abstract-talk-245}{}\label{poster-abstract-page-talk-245}
\subsection*{Poster 16. Determining Envirnoment-Driven Quenching Mechanisms with H$\alpha$ and Metallicity}
\textbf{Presenter:} David Tibben\\
\textbf{Poster number:} 16\\
\textbf{Fields:} Galaxy $\cdot$ Galaxy Clusters

Star formation rate and metallicity are fundamental diagnostics for understanding galaxy evolution. The objective of my research is to follow up a sample of cluster galaxies at 0.8$<$z$<$1 from the GOGREEN and GCLASS samples, with Keck/Gemini NIR spectroscopy to measure H$\alpha$ and [N II] emission lines. This allows us to infer star formation rates and metallicities for $>$500 galaxies. Comparing metallicity, SFR and stellar mass has the potential to distinguish between different environment-driven quenching mechanisms. We compare our cluster sample with field galaxies at a similar redshift observed with MOSDEF and 3D-HST. We find reduced star formation rates compared to the field, at fixed stellar mass. However, cluster galaxies show similar gas-phase metallicities as field galaxies of the same stellar mass and SFR. This suggests the dominant quenching mechanism may be a rapid process.

\hypertarget{poster-abstract-talk-61}{}\label{poster-abstract-page-talk-61}
\subsection*{Poster 17. Luminosity and Mass Functions for the Virgo Cluster}
\textbf{Presenter:} Eleanore Todd\\
\textbf{Poster number:} 17\\
\textbf{Fields:} Galaxy $\cdot$ Galaxy Clusters

The Virgo Cluster, at a distance of 16.5 Mpc, provides an ideal laboratory for testing models of galaxy formation in a dense and dynamic environment. Using data from the final data release of the Next Generation Virgo Cluster Survey (NGVS), we measure the luminosity function of 3689 Virgo member galaxies using a Bayesian fit Schechter function. For the 1409 galaxies in the vicinity of M87 and above the 50\% completeness limit, we find a faint-end slope of $\alpha$ = -1.31 0.01 in the g band, in agreement with earlier measurements based on smaller samples. We report on variations in the faint-end slope between the primary Virgo substructures and lower-density regions, and compare the luminosity functions for galaxies of differing morphological types and star formation activities. We also provide a comparison with the Illustris TNG-Cluster simulations, as well as a catalogue of galaxies in the Local Group.

\hypertarget{poster-abstract-talk-92}{}\label{poster-abstract-page-talk-92}
\subsection*{Poster 18. A Low Surface Brightness Analog to the Local Group}
\textbf{Presenter:} Jason Young\\
\textbf{Poster number:} 18\\
\textbf{Field:} Galaxy

Low surface brightness (LSB) spiral galaxies present a paradox: They are massive and HI rich, yet harbor little star formation. What keeps this gas in suspense? Isolation may play a role, since LSB spirals tend to be comparatively isolated. Here, we present a case study answering the question: What would happen if an LSB spiral were a member of a group? UGC 12695 is a Milky Way-mass LSB spiral, with a Milky Way-mass high surface brightness companion 294 kpc away and an LMC-mass LSB dwarf 139 kpc away. Thus, the composition of this group is broadly similar to the environment of the Milky Way, with the massive companion playing the role of the Andromeda Galaxy, making UGC 12695 a true LSB analog to the Milky Way. Using a multi-instrument campaign that includes IFU optical spectra and space-based near-IR imaging, we have determined the non-parametric star-formation histories of all three galaxies in a spatially resolved fashion. Here, we present those histories, and implications for past interactions between these galaxies.

\hypertarget{poster-abstract-talk-238}{}\label{poster-abstract-page-talk-238}
\subsection*{Poster 19. Canadian Gemini News}
\textbf{Presenter:} Eric Steinbring\\
\textbf{Poster number:} 19\\
\textbf{Fields:} Instrumentation $\cdot$ Next generation observatories

We provide updates on Gemini operations over the last year, with some statistics for Canadian use of the telescopes in recent semesters. This last year there was particularly big demand for high-resolution spectroscopy with GHOST, MAROON-X and IGRINS-2, for example. And did you know that last semester was probably the last time proposers will use the Phase I Tool (PIT)? The all-new all-in-one Gemini Program Platform (GPP) proposal/observing software is coming! Come check out more future plans, upcoming instruments and ongoing upgrades. The CASCA 2026 schedule has opportunities to meet with Gemini experts to learn about what's new and answer all your questions about Phase I and Phase II, or on getting help with your data reductions.

\hypertarget{poster-abstract-talk-64}{}\label{poster-abstract-page-talk-64}
\subsection*{Poster 20. Using Balloon-borne VLBI to Improve Event Horizon Scale Black Hole Imaging}
\textbf{Presenter:} Felix Thiel\\
\textbf{Poster number:} 20\\
\textbf{Fields:} Instrumentation $\cdot$ Supermassive black hole

This talk will focus on how balloon-borne telescopes can improve high-frequency Very Long Baseline Interferometry (VLBI) arrays. VLBI is a technique in radio astronomy which combines simultaneous observations from radio telescopes across the Earth, effectively giving the resolution of an Earth-sized telescope. Most recently, this technique has been employed by the Event Horizon Telescope (EHT) to observe the black holes at the centre of M87 and the Milky Way at a frequency of 230\,GHz on event horizon scales. The EHT in its current configuration is limited in two ways, on one hand in resolution by the diameter of the Earth and on the other in dynamic range (i.e. spatial scale coverage) by the lack of ground-based sites that are suitable for high frequency VLBI observations. Balloon-borne telescopes operate above 99.5\% of the Earths atmosphere and therefore have access to observing frequencies that are not accessible from the ground as well as locations that in ground-based astronomy would be inaccessible such as oceans and high-alpine mountain ranges. In this talk, using synthetic VLBI observations, I will discuss the impact that a millimetre and submillimetre balloon-borne VLBI station on a long-duration balloon flight could have on black hole science (such as jet-launching physics and photon ring observations) when paired with the EHT or next-generation EHT (ngEHT). I will further discuss the broad design requirements such as dichroic receiver systems that a balloon-borne telescope would need for millimetre and submillimetre VLBI. I will wrap up the talk by motivating the implications on black hole imaging that entire submillimetre balloon-borne VLBI arrays have when paired with a single ground-based site such as ALMA or the SPT.

\hypertarget{poster-abstract-talk-170}{}\label{poster-abstract-page-talk-170}
\subsection*{Poster 21. The Camera PAnoramique Proche InfraRouge (CPAPIR) Is Getting an Upgrade !}
\textbf{Presenter:} Philippe Vallée\\
\textbf{Poster number:} 21\\
\textbf{Field:} Instrumentation

CPAPIR is a wide-field infrared camera that has been in use at the Observatoire du Mont-Mégantic (OMM) since 2004. This workhorse instrument, for which a quarter of the detector is now non-functional, is currently being upgraded with a new detector. In addition, in order to significantly reduce our consumption of liquid nitrogen, CPAPIR is being refurbished with a Sunpower cryocooler. This poster presents the finite element analysis and the early test results carried out for this project.

\hypertarget{poster-abstract-talk-125}{}\label{poster-abstract-page-talk-125}
\subsection*{Poster 22. POET-STRATOS: Validating a SWIR Camera For Space Operations}
\textbf{Presenter:} Noel Wajnblum\\
\textbf{Poster number:} 22\\
\textbf{Field:} Instrumentation

The POET-STRATOS project aims to validate the performance of a short-wavelength infrared (SWIR; 6001700 nm) camera - an OWL 1280 from Raptor Photonics, Inc. - in a near-space environment as a technology demonstration for the proposed POET mission (Rowe et al. 2022). The camera is sold in flight-ready aluminum shielding, with which it already has commercial space heritage. However, its usability for space-based astronomy has yet to be validated. We characterize the detectors photometric stability, temperature control performance, detector background behavior, and radiation tolerance under near-space conditions, demonstrating that commercial SWIR sensors can achieve repeatable flux measurements suitable for precision astronomical photometry. During a stratospheric balloon flight, the detector maintained a stable operating temperature of -27.7 C and dark current of $\sim$66 e-/sec over a range (-20 to 5 C) of ambient temperatures. Subsequent proton irradiation testing at TRIUMF showed that the camera remained operational up to a cumulative proton dose of 15 krad, although long-exposure capability was reduced beyond 9 krad - equivalent to about 3 years of anticipated irradiation in low-Earth orbit. These results demonstrate that commercially available SWIR detectors can provide stable photometric performance in near-space environments and are viable instruments for future astronomical observations.

\hypertarget{poster-abstract-talk-219}{}\label{poster-abstract-page-talk-241}
\subsection*{Poster 23. Extending NIRPS to the K-Band: Optical and Opto-Mechanical Design of a Cryogenic High-Resolution Spectrograph.}
\textbf{Presenter:} Denis Brousseau\\
\textbf{Poster number:} 23\\
\textbf{Field:} Instrumentation

We present an optical and opto-mechanical design concept of a cryogenic precise (m/s) spectrograph operating in the K band (19502450 nm), designed as an extension to NIRPS high resolution spectrograph, already in operation on the ESO 3.6 m telescope. The instrument adopts an R4 echelle grating achieving a resolving power of R $\sim$ 100,000, within a compact, fiber-fed architecture optimized for thermal stability. This addition complements the existing YJH coverage and provides access to the fundamental rovibrational CO bandhead, a key diagnostic for exoplanet atmospheric characterization. This K-band extension substantially increases the scientific reach of NIRPS by opening a spectral window essential for probing the chemistry, structure, and formation history of exoplanet atmospheres.


\hypertarget{poster-abstract-talk-337}{}\label{poster-abstract-page-talk-337}
\subsection*{Poster 24. Maria-based Atmospheric Simulations for CCAT Observations from 220-850 GHz}
\textbf{Presenter:} Nicholas Zaparniuk\\
\textbf{Poster number:} 24\\
\textbf{Fields:} Instrumentation $\cdot$ Next generation observatories

The CCAT Observatory, with its 6-m wide-field Fred Young Submillimeter Telescope, is located 5600\,m above sea-level on Cerro Chajnantor in northern Chile and designed to enable large-area surveys of the millimetre and submillimetre sky. Prime-Cam, its multi-wavelength camera and spectrometer, will provide broadband and spectroscopic observations across multiple frequency bands from 220 and 850\,GHz to address science questions ranging from cosmology through Galactic star formation. We use the Maria simulation framework to construct a simplified model of the Prime-Cam instrument and investigate how atmospheric opacity and detector loading will affect the observations. Using simulated observations across the Prime-Cam frequency range, we characterize the dependence of atmospheric temperature and inferred zenith opacity on precipitable water vapour (PWV) and elevation, and use these trends to estimate changes in detector loading during large scan motions. We find that the simulated atmosphere produces an approximately linear relation between PWV and inferred zenith opacity, with higher-frequency bands showing stronger atmospheric sensitivity and more restrictive observing conditions. These results provide an initial framework for evaluating the feasibility of wide-elevation scans and time-domain observations with Prime-Cam, including long-term monitoring of bright nearby star-forming regions to investigate episodic mass accretion onto deeply embedded protostars.

\hypertarget{poster-abstract-talk-254}{}\label{poster-abstract-page-talk-254}
\subsection*{Poster 25. Evaluating the Performance of Composite Material in Space and Stratospheric Environments}
\textbf{Presenter:} Seyedmahdiar Ziaolhagh\\
\textbf{Poster number:} 25\\
\textbf{Fields:} Instrumentation $\cdot$ Next generation observatories

Lightweight and thermally stable structural materials are essential for modern space-based astronomical instrumentation operating in cryogenic environments. Carbon fibre reinforced polymers (CFRP) are promising candidates for spacecraft and telescope structures due to their high stiffness-to-weight ratio, low thermal conductivity, and strong mechanical performance compared to traditional metallic alloys. The cryogenic material properties of CFRP must therefore be accurately characterized for space missions. An overall theme of this work is to understand the mechanical and thermal performance of CFRP composites for use in space instrumentation. Early stages for discussion now include: (1) measurement of Youngs modulus of a CFRP cantilever beam at room and cryogenic temperatures, (2) examination of the outgassing properties of CFRP material under vacuum using a residual gas analyzer, and (3) determination of the coefficient of thermal expansion of CFRP using a frequency-modulated laser metrology system during thermal cycling between room temperature and cryogenic conditions. These measurements provide constraints on the structural stability, thermal contraction, and contamination behaviour of CFRP materials in vacuum environments relevant to space-based instrumentation and observatories. The results will contribute to the design and reliability assessment of next-generation cryogenic astronomical instruments and space telescope structures and components.

\hypertarget{poster-abstract-talk-114}{}\label{poster-abstract-page-talk-114}
\subsection*{Poster 26. Driving Mechanism of Protostellar Outflows in a 3D Non-Ideal MHD Simulation.}
\textbf{Presenter:} Majd Noel\\
\textbf{Poster number:} 26\\
\textbf{Field:} Stars and Stellar cluster formation

Magnetic fields play a pivotal role in extracting angular momentum from the disk and driving jets and outflows in protostellar systems, yet the full launching mechanism remains incompletely understood. This work examines the launching mechanism of wide-angle outflows in class 0 phase of star formation. Surprisingly, the outflow is launched from an outer pseudo-disk that has good magnetic coupling, rather than from the inner Keplerian disk where the magnetic coupling is poor. The wide-angle wind carries away a significant fraction of the mass that would otherwise have fallen onto the star-disk system. By tracing the velocity streamlines, calculating the forces and the conserved parameters, we compare our model with steady-state axisymmetric MHD wind models, such as the magnetocentrifugal wind theory and magnetic tower flow. This work is crucial for enhancing our understanding of the disk wind theory and the star formation process in general. These results will guide future observational tests of disk wind models with facilities such as ALMA and JWST.

\hypertarget{poster-abstract-talk-153}{}\label{poster-abstract-page-talk-153}
\subsection*{Poster 27. Uncovering the Structure and Star Formation History of the SCYA-89 Stellar Association}
\textbf{Presenter:} Mayesha Rahaman\\
\textbf{Poster number:} 27\\
\textbf{Field:} Stars and Stellar cluster formation

SCYA-89 is a nearby ($\sim$490\,pc) young stellar association identified in Gaia surveys of pre-main-sequence stars and later cataloged as part of the SPYGLASS (Stars with Photometrically Young Gaia Luminosities Around the Solar System) program. Despite its clear youth and coherent structure, it has not previously been the focus of a dedicated study. This project presents the first detailed analysis of SCYA-89, with the goal of understanding how star formation has unfolded within it. We identify distinct subgroups through clustering based on shared spatial and kinematic structure, characterize their relative ages, and determine whether any age variations trace a coherent gradient or reflect localized formation events. Establishing the age structure of such populations helps reconstruct the formation history of nearby stellar associations and provides important age constraints for objects such as planetary systems that form within them. To improve age estimates, we develop a mixture-model sampling framework to better account for field contamination when determining isochronal ages. We also investigate the regions internal dynamics through expansion trends and dynamical traceback analyses, and measure population demographics for the association and its subgroups while correcting for low-mass incompleteness and unresolved multiplicity. Combining ages and dynamics, we assess whether the structure of SCYA-89 reflects preserved primordial structure or sequential star formation driven by stellar feedback.

\hypertarget{poster-abstract-talk-240}{}\label{poster-abstract-page-talk-240}
\subsection*{Poster 28. Bound-free continuum in a strongly magnetic white dwarf's atmosphere.}
\textbf{Presenter:} Jonathan Roussy\\
\textbf{Poster number:} 28\\
\textbf{Field:} Stars and Stellar cluster formation

In the presence of strong magnetic fields, the structure of the hydrogen atom is profoundly modified. While the splitting and shifting of bound energy levels are well known, the continuum states are also strongly affected, becoming structured and shifted due to magnetic quantization. These changes can significantly alter the boundfree opacity in stellar atmospheres. Recent quantum-mechanical calculations provide accurate descriptions of hydrogen boundfree spectra in strong magnetic fields, but their implementation in astrophysical models remains limited. Building on the recent work of Gmez et al. (2024), which presents detailed calculations of the boundfree spectrum using advanced numerical techniques, we incorporate these results into a magnetic stellar atmosphere code. This development enables improved modeling of the continuum opacity in strongly magnetic white dwarf atmospheres and provides a more realistic framework for interpreting their observed spectra.

\hypertarget{poster-abstract-talk-48}{}\label{poster-abstract-page-talk-48}
\subsection*{Poster 29. Mining the ALMA Archive: Building a Catalog of Structured Young Protostellar Disks Preliminary Results!}
\textbf{Presenter:} Lance Schonberg\\
\textbf{Poster number:} 29\\
\textbf{Field:} Stars and Stellar cluster formation

Observations of substructure in the form of rings and gaps in the disks of young stellar objects (YSOs) are a potential indication of future or even early current planetary formation. To date, most research and observations have been done on older, unembedded disks, but structure and possibly planet formation may already be happening in the youngest protostellar disks, those still embedded in their envelopes. Some previous research has shown structure in a small number of these younger YSOs.

Specialized code packages allow us to examine individual disks using their visibilities, achieving sub-beam resolution and giving the ability to consider lower resolution data for processing. Selecting archival ALMA projects for resolution and sensitivity, we are building a catalog of young disks with and without substructure.

At this time, there are on the order of ten younger YSOs with confirmed substructure, but my catalog has improved on this number by more than an order of magnitude, making it the largest collection of young disks with substructure to date. I will present preliminary results covering the approximately 200 objects that meet selection criteria out of 2000 surveyed across more than a dozen previous ALMA projects spanning more than seven clouds. I find that roughly 50\% of the selected objects appear to have substructure, perhaps indicating it may be common in these young disks. Additionally, I have found objects that do not fit my current analysis models, including binary objects and case where the disks have significant extended structure. The final catalog, including these less regular items along with the categorized YSOs will be a valuable resource for the star formation community.

\hypertarget{poster-abstract-talk-57}{}\label{poster-abstract-page-talk-57}
\subsection*{Poster 30. Mergers, mass transfer, and magnetic fields (oh my!): new evidence for stellar magnetic field generation during binary evolution}
\textbf{Presenter:} Gregg Wade\\
\textbf{Poster number:} 30\\
\textbf{Field:} Stars and Stellar cluster formation

Recent observations and modeling of massive stars lead to the inescapable conclusion that the classical theory of stellar evolution - envisioning single stars evolving in an isolated fashion - is fundamentally wrong in many, and quite possibly most, cases. A majority of massive stars are known to form in binary or higher-order multiple systems, and as these stars evolve they expand to fill their Roche Lobes, resulting in mass and angular momentum exchange. The resultant stripped star mass donors and rejuvenated mass recipients are characterized by masses, rotational properties, and surface chemistry that are entirely at odds with evolution models of single stars.

The rapid and violent evolution resulting from mass transfer and merger events also provide fertile grounds for the generation of powerful magnetic fields. In this presentation we review several examples of binary evolution implicating magnetic fields. HD 45166 (Shenar et al. 2023), a system consisting of a main sequence B-type primary and a stripped star companion, was recently discovered to host an extraordinarily strong magnetic field, likely a consequence of the merger of two intermediate-mass stars. Similarly, Plaskett's star (Wade et al. 2026), a massive O-type binary, comprises a magnetic component and a (non-magnetic) stripped-star companion that has recently undergone mass transfer. Finally, we present first results for LB-1 and HR 6819 (Budz et al., in prep), two systems hypothesized to have undergone evolution similar to Plasketts system. These case studies highlight the intricate interplay between binary interactions and magnetic fields, shedding light on the diverse evolutionary pathways of massive stars and the role of magnetic fields in shaping their observable properties.

\hypertarget{poster-abstract-talk-25}{}\label{poster-abstract-page-talk-25}
\subsection*{Poster 31. Practically perfect pixelated PSF posterior samples}
\textbf{Presenter:} Connor Stone\\
\textbf{Poster number:} 31\\
\textbf{Field:} Astrostatistics, Astroinformatics, and AI/Machine Learning

Point Spread Function (PSF) modelling is critical for precision astronomical analyses including: photometry, astrometry, weak gravitational lensing, and strong gravitational lensing. A key challenge to properly characterizing the PSF is that it changes with time and detector position, and we may only access noisy and low-resolution realizations (images of stars). To avoid overfitting noise, traditional methods tend to rely on simplified parametric models, or ad-hoc regularization. In this talk, I present a fully Bayesian approach that leverages a prior over PSF models learned from high resolution template images, and a simple Gaussian likelihood. Our posterior samples yield exquisite detail, and are inherently robust to noise and masking. Crucially, the residuals from our posterior samples are statistically indistinguishable from noise while the traditional methods produce substantial and structured residuals. Further, we may produce many realizations of our PSF posterior model to allow marginalization in downstream analysis tasks. I will conclude with examples of the PSF effect on photometric accuracy and strong gravitational lensing source reconstruction.

\hypertarget{poster-abstract-talk-229}{}\label{poster-abstract-page-talk-229}
\subsection*{Poster 32. Can Simulation-Based Inference Reshape the Search for Structured Protostellar Disks?}
\textbf{Presenter:} Callista Sullivan\\
\textbf{Poster number:} 32\\
\textbf{Fields:} Astrostatistics, Astroinformatics, and AI/Machine Learning $\cdot$ Stars and Stellar cluster formation

For over a decade, the Atacama Large Millimeter Array (ALMA) has been building an archive of protostellar disk observations that now numbers in the thousands. As an interferometer, however, ALMAs observations consist of large data tables rather than human-interpretable images. While quick analyses of metrics such as u-v visibilities have yielded some scientific insights, detailed disk modelling currently requires long computation times. In particular, identifying structure such as gaps or rings (i.e., regions of diminished or excess emission, respectively) has been prohibitively slow. This is a problem, as these key features are the best laboratories in which to study the competing influences of the various forces and physical processes that dictate how disks and planets evolve.

My work aims to enable the study of structure in ALMAs entire observational backlog by creating a neural network that can quickly interpret these observations and identify the structure in each disk using Simulation-Based Inference. Once preliminary testing has been completed, it will be trained on tens of thousands of simplified geometric models of both smooth and structured disks, alongside a set of pre-analysed observations that have already had their structure identified. After training, it will be tested against the remainder of all publicly available pre-analysed disks to validate its accuracy. At that point it will require only a few minutes to identify the structure in each real disk, compared to the current standard of several hours or longer.

With this drastic increase in speed, my model aims to shift the current paradigm from one of analysing individual disks to one where it is possible to measure the presence, form, and statistics of structure among large disk populations.

\hypertarget{poster-abstract-talk-336}{}\label{poster-abstract-page-talk-336}
\subsection*{Poster 33. Transformer-Based Radio Interferometric Imaging Algorithm for Next-Generation Telescopes}
\textbf{Presenter:} Qitong Anabel Tan\\
\textbf{Poster number:} 33\\
\textbf{Fields:} Astrostatistics, Astroinformatics, and AI/Machine Learning $\cdot$ Galaxy $\cdot$ Galaxy Clusters $\cdot$ Next generation observatories

Interferometric imaging plays a critical role in the measurement and interpretation of galaxy clusters in the radio wave band, where diffuse emission traces large-scale physical processes such as feedback from active galactic nuclei. However, the fidelity of reconstructed astronomical images depends on the imaging and deconvolution algorithms used to recover the sky brightness from incomplete Fourier sampling. Traditional approaches, such as CLEAN and its variants, often struggle to accurately reconstruct extended structures. More recent methods based on compressed sensing or machine learning can significantly improve image quality, but they are often computationally inefficient for the high-throughput data flows expected from next-generation radio telescopes. Canada committed to becoming a full member of the Square Kilometre Array (SKA) Observatory in 2023. Without algorithms capable of handling SKA-scale data, astronomers may not be able to achieve the expected scientific return from such next-generation telescope in a timely manner. Therefore, we develop a novel Transformer-based imaging algorithm to improve the robustness and fidelity of interferometric reconstructions of extended sources on large scales. The method is trained and validated using both simulated and real datasets, allowing the model to learn structural features of extended celestial sources. Applied to observations of the Perseus galaxy cluster, we benchmark its performance against widely used deconvolution algorithms using field-specific quality assessors. The algorithm provides a promising pathway toward high-fidelity, large-scale imaging with upcoming radio telescopes, enabling more accurate studies of galaxy clusters and other extended astrophysical systems.

\hypertarget{poster-abstract-talk-275}{}\label{poster-abstract-page-talk-275}
\subsection*{Poster 34. Automatically discovering optimal analytic dark matter halo profiles with deep reinforcement learning}
\textbf{Presenter:} Wassim Tenachi\\
\textbf{Poster number:} 34\\
\textbf{Fields:} Astrostatistics, Astroinformatics, and AI/Machine Learning $\cdot$ Galaxy

The equations describing dark matter (DM) halo profiles are essential for modeling and constraining DM behavior in both simulations and observations. While numerous profiles exist - such as the widely used NFW, Einasto, and Dekel-Zhao models - many lack key analytic properties necessary for broad applicability. For instance, the NFW profile has a non-converging enclosed mass at infinity, and many models do not provide closed-form expressions for essential quantities like circular velocity, radial velocity dispersion, or gravitational potential, complicating their use in rotation curve analysis, gravitational lensing studies, and semi-analytical galaxy formation and evolution simulations.

To address these limitations, I introduce a symbolic regression (SR) framework based on the PhySO code, which directly searches for analytic profile equations through trial and error using deep reinforcement learning. Unlike traditional black-box machine learning approaches, which often lack interpretability, our method optimizes symbolic expressions not only for accuracy but also for essential analytic properties, such as derivations of closed-form key physical quantities and overall physical consistency. Moreover, we enforce parsimony, favoring expressions with minimal parameters, leading to a threefold symbolic optimization objective: accuracy, tractability, and low complexity.

Applying this approach to the NIHAO simulation suite, we explore halo profiles in both hydrodynamical and dark-matter-only scenarios. Preliminary results suggest a Pareto-optimal alternative to existing models, providing a more versatile and physically grounded framework for DM halo studies.

\hypertarget{poster-abstract-talk-323}{}\label{poster-abstract-page-talk-323}
\subsection*{Poster 35. Probing the faintest galaxies below the confusion limit}
\textbf{Presenter:} Yunting Wang\\
\textbf{Poster number:} 35\\
\textbf{Fields:} Astrostatistics, Astroinformatics, and AI/Machine Learning $\cdot$ Galaxy

The far-infrared to submillimetre sky is composed of emission from dust-enshrouded star-forming galaxies. The number counts of these galaxies provide key constraints on models of galaxy evolution in cosmological simulations. To obtain the strongest constraints possible, the faintest galaxies below the confusion flux limit can be statistically leveraged using the "probability of deflection" (or P(D)) method. However, the correlation between galaxies, i.e., large-scale clustering of galaxies, complicates the analysis and has been shown to biases the results of galaxy number counts near the confusion limit. We outline a method that simultaneously measures and corrects for galaxy clustering based on a P(D) framework, yielding unbiased constraints on the distribution of faint galaxies below the confusion limit. Applying this method to Herschel-SPIRE observations of the GOODS-N field, we demonstrate that our correction recovers robust sub-mJy number counts. We find that at 500$\mu$m clustering increases the apparent counts by about a factor of 1.6 around 10mJy and slightly decreases the faintest counts, with smaller effects at 350$\mu$m and 250$\mu$m.This methodology is broadly applicable to other confusion-limited data sets where beams and noise are well characterized, including SCUBA-2, CCAT, and relevant for future facilities like PRIMA and AtLAST where source blending will remain a challenge, making it possible to fully exploit the maximum information from current and future far-infrared/submillimetre surveys.

\hypertarget{poster-abstract-talk-205}{}\label{poster-abstract-page-talk-205}
\subsection*{Poster 36. Detection of Gaseous Components of the Epsilon Eridani Debris Disk with SITELLE}
\textbf{Presenter:} Mehar Sahota\\
\textbf{Poster number:} 36\\
\textbf{Field:} Gas, dust and the ISM

Debris disks are circumstellar disks found around main sequence stars, sustained by ongoing collisions between planetesimals remaining from planet formation processes. These collisions produce gas and dust, rendering asteroid and cometary belts observable as the planetesimals collide. Recent detections of gas emission in some young debris disks motivated the use of the SITELLE, an optical imaging Fourier transform spectrometer, currently on the Canada France Hawaii Telescope (CFHT), to test whether this instrument can reliably detect gas in debris disks. In this pilot project, the nearby ($\sim$3\,pc) star Epsilon Eridani was observed in the Fe I emission line. The debris disk around Epsilon Eridani is well-characterized and spatially resolved by observations from the Atacama Large Millimeter Array (ALMA), providing a good baseline for modeling the disk in SITELLE data. Because SITELLE does not have a coronagraph, we must rely on PSF subtraction to remove the stars emission to allow detection of the relatively faint disk. Eps Eri was chosen for its rather large extent on the sky, and its face-on orientation, which guarantees that the gas will be observed in emission. Epsilon Eridanis disk is also known to be at least 1 Gyr old, implying that any gas detected in the disk cannot be residual from the systems protoplanetary phase, and instead must be generated from the debris disk itself. We aim to determine if the presence of the disk can be reliably inferred in the SITELLE data, and will present our findings.

\hypertarget{poster-abstract-talk-15}{}\label{poster-abstract-page-talk-15}
\subsection*{Poster 37. Multi-wavelength study of the quasar J1000+1242}
\textbf{Presenter:} Mélissa St-Pierre\\
\textbf{Poster number:} 37\\
\textbf{Fields:} Gas, dust and the ISM $\cdot$ Supermassive black hole

Quasars are extremely luminous active galactic nuclei characterized by high accretion rates. Located at a redshift of z=0.148, J1000+1242 is a type 2 quasar, defined by a dusty torus whose orientation obscures the accretion disk along our line of sight. I will present two sets of new hyperspectral data for this object. The first one was obtained with the SITELLE imaging Fourier transform spectrometer (CFHT) covering the [OIII] $\lambda$4959 and [OIII] $\lambda$5007 emission lines. The second one was obtained with the MUSE integral field spectrograph (VLT), providing higher spectral and spatial resolution, but a smaller field of view than SITELLE. These observations reveal a large amount of ionized gas extending over a total of 165 kpc on either side of the nucleus and, for the first time, allow the full extent of the quasars outflows to be imaged. Observations with the VLA and e-MERLIN have revealed a faint radio jet extending over approximately 25 kpc. ALMA has also detected molecular gas filaments extending from 5 to 11 kpc. These structures will be compared with the ionized gas as detected by MUSE and SITELLE. This comparison, in addition to the study of the gas dynamics and its excitation mechanisms, will help clarify the interactions between the supermassive black hole and its host galaxy (feedback).J1000+1242 highlights the importance of studying moderately powerful jets due to their significant impact on their surroundings.

\hypertarget{poster-abstract-talk-242}{}\label{poster-abstract-page-talk-242}
\subsection*{Poster 38. Ghost Nebulae of the Triangulum Galaxy: Studying He II and O III Emission from Nebulae without Detected Ionizing Sources}
\textbf{Presenter:} Arina Tseragotin\\
\textbf{Poster number:} 38\\
\textbf{Fields:} Gas, dust and the ISM $\cdot$ Local Group

Ionized nebulae are a key laboratory for understanding stellar evolution and the properties of the interstellar medium. A new class of nebulae, however, have been found to exhibit extended ionized emission line regions, yet do not appear to be associated with any known source of ionizing radiation. Two of these objects, HBW673 and BCLMP651 in the nearby Triangulum galaxy (M33), exhibit particularly strong He II 4686 emission, associated with ionization by very hot and luminous objects. This may point to a rare phenomenon of photoionization caused by supersoft X-ray sources (SSS) or a long-gone fossil x-ray nebula. Using the University of Manitobas Glenlea Research Observatory, our goal is to record the surface brightness of both targets using a He II 4686 filter. Therefore, we will be able to compare the surface brightnesses with CAL83, the only known nebula that indicates ionization by SSS. Additionally, we plan to apply an upper x-ray constraint from different x-ray catalogues, such as XMM-Newton, of HBW673 and BCLMP651 to pick up low-emitting x-ray sources to further confirm that the emission is caused by SSS. Lastly, we aim to use the line flux ratios derived using the photoionization code Cloudy, to determine the best fitting model for the photoionization source. Time permitting, we will also be able to identify the elapsed time since the photoionized source has shut off to reveal if the mystery behind these ghost nebulae may actually indicate a fossil x-ray nebula.

\hypertarget{poster-abstract-talk-314}{}\label{poster-abstract-page-talk-314}
\subsection*{Poster 39. Tracing Fullerene-Linked Dust Variation in Tc 1 with JWST}
\textbf{Presenter:} Simon Van Schuylenbergh\\
\textbf{Poster number:} 39\\
\textbf{Field:} Gas, dust and the ISM

We have observed the archetype fullerene-rich planetary nebula Tc 1 with the MIRI/MRS on JWST. We study the dust environment in detail, and link dust variation to fullerene formation processes. Understanding the dust processing in sources like Tc 1 is crucial to developing our understanding of the formation of aromatic molecules in space. The mid-infrared spectra of this spatially resolved source reveal that the fullerene emission is not co-spatial with the dust emission. Two broad plateau-like features can be distinguished in the mid-IR spectrum: one at 69 um and one at 913 um. Spectral variations of the 913 um plateau reveal at least three distinct components, peaking at 9.5 um, 11 um and 12 um respectively. SiC produces most likely at least one of those components and we find evidence that a second component is likely produced by SiC grains too. The variation in relative strength of these two components across the field of view hints at dust processing taking place in Tc 1. The third component of the 913 um plateau peaking at 9.5 um has a strong correlation with the SiC component(s) and is thus likely produced by similar material. The 69 um plateau does not change shape spectrally across the entire field of view, indicating it is produced by stochastically heated material. We discuss the ramifications of these results for the formation of aromatic molecules.

\hypertarget{poster-abstract-talk-165}{}\label{poster-abstract-page-talk-165}
\subsection*{Poster 40. Identifying Velocity-Coherent Gas Structures in KEYSTONE Survey Data}
\textbf{Presenter:} Sean Yin\\
\textbf{Poster number:} 40\\
\textbf{Fields:} Gas, dust and the ISM $\cdot$ Stars and Stellar cluster formation

Stars form in large clouds of gas and dust called molecular clouds, which naturally form filamentary structures. Studies have shown that filaments play an important role in the star formation process, and their dynamics are governed by a mix of gravity, turbulence, and magnetic fields. To study high-mass star formation, we must probe giant molecular clouds (GMCs), which are the most efficient sites of star formation. The KEYSTONE survey project on the Green Bank Telescope studied 11 GMCs in ammonia (NH3) emission, and recent work has found multiple velocity components in the NH3 gas. My aim is to identify independent regions of the clouds with smoothly varying kinematic properties, known as velocity-coherent structures (VCSs). By identifying VCSs, I can accurately measure the gas flows in these GMCs and analyze correlations to other properties such as the gravitational field and the local magnetic field orientation. Nevertheless, finding the VCSs when there are multiple velocity components is non-trivial, and requires the use of robust clustering algorithms. In this presentation, I share preliminary results of VCS identification of the KEYSTONE data using the clustering algorithm HDBSCAN. Future work will involve following standard procedures to identify filament spines and examine the relationships between velocity gradients, gravity gradients, and the magnetic field. My research will provide the most accurate picture of the underlying gas structures in the GMCs in the KEYSTONE survey, and help improve our greater understanding of high-mass star formation.

\hypertarget{poster-abstract-talk-209}{}\label{poster-abstract-page-talk-209}
\subsection*{Poster 41. Unveiling the Unseen Pulsar Population - CHAMPSS Science Highlights}
\textbf{Presenter:} Robert Main\\
\textbf{Poster number:} 41\\
\textbf{Field:} Compact Objects

The CHIME All-Sky Multiday Pulsar Stacking Search (CHAMPSS) is a massive, Canadian-based pulsar survey which is beginning to operate at scale. This survey leverages CHIMEs unique cylindrical design, probing a largely unexplored parameter space when contrasted to one-pass surveys using traditional parabolic dishes. Through repeated daily observing of the northern sky, we are sensitive to exotic highly intermittent pulsars. Through stacking, we probe deeper than full-sky surveys to date, particularly to faint pulsars away from the Galactic plane.

The survey began at scale in January 2026, and pulsar discoveries have begun pouring in. In this talk, I will detail scientific highlights from the discoveries in the first half-year of large-scale operation. A surprising number ($\sim$30\%) of our new pulsars reveal an electron column in excess of the expectation from state-of-the art Galactic electron models, revealing either a population of highly luminous pulsars in the Galactic halo, or showing the limitations of these models on under-constrained sightlines. We have discovered a faint pulsar within the local bubble, likely the least luminous radio pulsar to date. We have discovered a binary pulsar in a $\sim$170 day orbit, likely associated with a red-dwarf companion observed by Gaia, the first pulsar binary of its kind. Alongside CHIME/FRB, we have co-discovered a pulsar which is highly luminous ($>$3 kJy peak), yet highly intermittent (duty cycle of $<$0.5\%); despite its brightness, this source would be almost certainly missed in a one-pass survey, hinting at a massive undiscovered population of highly intermittent sources.

\hypertarget{poster-abstract-talk-174}{}\label{poster-abstract-page-talk-174}
\subsection*{Poster 42. Evaluating magnetars as the progenitors of fast radio bursts}
\textbf{Presenter:} Vishwangi Shah\\
\textbf{Poster number:} 42\\
\textbf{Fields:} Compact Objects $\cdot$ High Energy Astrophysics $\cdot$ Transients

Fast radio bursts (FRBs) are micro-to-millisecond-duration radio transients originating from extragalactic distances. While thousands of FRBs have been detected, their origins remain unknown. Many FRB theories invoke highly magnetized neutron stars, called magnetars, as FRB progenitors. In this talk, I will evaluate the magnetar progenitor scenario using recent results from the CHIME/FRB telescope. The most energetic FRBs provide crucial insights into their emission mechanisms and progenitor models. I will present constraints on the maximum energy of FRBs derived from the largest sample of $\sim$3,000 one-off bursts in the second CHIME/FRB catalog. Our results indicate that the isotropic energies of the most energetic FRBs are collectively limited to $\sim 10^{42}$ erg, thereby constraining the energy reservoir of their sources. Under standard assumptions for the radiative efficiency factor between radio and high-energy emission, this energy limit is consistent with the energies of the most powerful magnetar giant flares. Furthermore, the occurrence of most FRB sources inside star-forming galaxies supports an origin in magnetars formed via core-collapse supernovae. However, I will present the first ever localization of a repeating FRB far outside a quiescent galaxy, which challenges this theory. We propose that this source resides in a globular cluster, where the magnetar progenitor scenario is only viable if the magnetar formed via exotic formation channels such as compact-object mergers.

\hypertarget{poster-abstract-talk-249}{}\label{poster-abstract-page-talk-249}
\subsection*{Poster 43. An Overview of Current LIGO-Virgo-KAGRA Continuous Wave Results and Detector Characterization Efforts}
\textbf{Presenter:} Taylor Starkman\\
\textbf{Poster number:} 43\\
\textbf{Fields:} Compact Objects $\cdot$ Gravitational waves $\cdot$ High Energy Astrophysics

The current generation of ground-based gravitational wave detectors have detected hundreds of mergers of compact objects, but as of yet, no sources of continuous gravitational waves (CWs). Most likely produced by non-axisymmetric neutron stars, these signals would provide key insights into neutron star equations of state inaccessible to experiments on Earth. These signals are predicted to be extremely weak, only detectable over long periods of time, and current upper limits place them at less one-thousand times the strength of currently detectable gravitational wave events. Improving detector sensitivity is critical to increasing the chances of making a detection. Specifically, we must characterize and mitigate as many narrow spectral artifacts as possible. Colloquially known as lines, these artifacts could mimic or obscure a CW signal. In this presentation, I will provide an overview of current LIGO-Virgo-KAGRA continuous wave results and the detector characterization work that makes these results possible. I will also discuss challenges for this analysis, including work to increase the speed at which we are able to mitigate lines and significantly improve data quality in future observing runs.

\hypertarget{poster-abstract-talk-296}{}\label{poster-abstract-page-talk-296}
\subsection*{Poster 44. Is GW231123 a strongly lensed event?}
\textbf{Presenter:} Mariah Zeroug\\
\textbf{Poster number:} 44\\
\textbf{Fields:} Compact Objects $\cdot$ Cosmology and Early Universe $\cdot$ Gravitational waves

GW231123 is an exceptional binary black hole merger whose unusually large component masses challenge standard stellar-origin expectations. Additionally, its secondary mass (the less massive one) is in tension with the observed gravitational-wave population. Strong gravitational lensing provides a potential explanation, as magnification biases the inferred source parameters, making a high-redshift, lower-mass system appear much closer and more massive. In this talk, I will present the results of an investigation into whether or not gravitational lensing can plausibly account for the observed properties of GW231123. We find that reconciling its secondary mass with the observed population would require an extreme magnification. The probability of observing such a magnification is extremely low given current gravitational wave observations, strongly disfavoring the lensing hypothesis.

\hypertarget{poster-abstract-talk-325}{}\label{poster-abstract-page-talk-325}
\subsection*{Poster 45. Inside the CFHT Control Room, during a Shadow the Scientists Event}
\textbf{Presenter:} Viraja Khatu\\
\textbf{Poster number:} 45\\
\textbf{Field:} Education and Public Outreach

Pyjama-mode observing observing with telescopes from the comfort of our homes emerged as a popular setup in a pandemic-struck world back in 2020. Amidst lockdown, researchers faced issues with conducting scientific work at their pre-pandemic pace. However, students at educational institutions were in greater trouble with no access to outdoor activities and no active learning settings. The Creating Equity in Science, Technology, Engineering, Art, and Math (CrEST) team at the University of California Santa Cruz took this difficult hour as an opportunity to transform the pyjama-mode observing into Shadow the Scientists (StS) a one-of-a-kind educational outreach program. StS offers a way for students to witness STEAM scientists and professionals working live on their projects. Participants can interact, first-hand, with the Scientist being shadowed and engage in conversations related but not limited to the topic of the event. One of the most successful virtual programs, StS aims to inspire the younger generation in pursuing STEAM careers as well as build fundamental skills required to pursue a career in those fields. Equally popular amongst the public, StS strives to reach the underprivileged communities where people have restricted or no access to open, interactive learning. Canada-France-Hawaii Telescope (CFHT) has entered its second year of StS with seven successful events completed thus far. At CFHT, we share our behind-the-scenes preparation for an observing night, the observing setup as our Remote Observers collect telescope data, and more all from our control room at CFHT Headquarters. In my talk, I will cover CFHT becoming a part of StS, content and stories from our StS events, and opportunities for the CASCA audience to participate or collaborate with us.

\hypertarget{poster-abstract-talk-17}{}\label{poster-abstract-page-talk-17}
\subsection*{Poster 46. The ALMA Primer Instructional Video Series}
\textbf{Presenter:} Brenda Matthews\\
\textbf{Poster number:} 46\\
\textbf{Field:} Education and Public Outreach

The Atacama Large Millimeter/submillimeter Array (ALMA) observatory, consisting of 66 telescopes operating together at 5000m elevation in Chile, is one of the most complex astronomical facilities in the world. The Millimetre Astronomy Group at NRC-HAA leads the ALMA Primer Instructional Video Series in collaboration with the North American ALMA Science Center, with contributions from astronomers across North America. Undergraduate co-op students have played a major role in developing the videos, with contributions thus far from six students across 5 different institutions. This series is designed to provide a basic introduction to radio interferometry, calibration, imaging, and other topics in short (5-10 minute), easy-to-digest segments, and is targeted toward professional astronomers who may be unfamiliar with the concepts and who want to use ALMA for their science. The video series can also be used in the classroom to teach basic concepts of radio interferometry at a graduate or undergraduate level. The videos can be accessed through the ALMA Science Portal (https://almascience.org/tools/alma-primer-videos) where theyre described and categorized, or through our YouTube channel (https://youtube.com/@almaprimer920). This is a work in progress, with new videos released periodically. As of this writing we have published 20 videos, with $\sim$20,000 views total. In this presentation well discuss the philosophy behind the series, the process by which the videos are created, and show brief snippets of some of our offerings.

\hypertarget{poster-abstract-talk-210}{}\label{poster-abstract-page-talk-210}
\subsection*{Poster 47. Progress on the Westar Program}
\textbf{Presenter:} Dan Morrone\\
\textbf{Poster number:} 47\\
\textbf{Fields:} Education and Public Outreach $\cdot$ Equity, Diversity and Inclusion $\cdot$ Indigenous Engagement

The Westar Program, offered by the Westar Subcommittee, part of CASCAs Education and Public Outreach (EPO) Committee, is a revamped version of its predecessor, the Westar Lectureship. The Westar Program is a collaborative initiative to bring astronomy to remote and Indigenous Canadian communities and foster long term and reciprocal relationships with these communities.

The revamp of this program started in 2022, incorporating feedback from 12 individuals and consultations from Two Worlds Consulting. Since then, the Westar Coordinator, Dan Morrone, was hired to develop the program and compile training, hands-on activities, and other resources for the program.

Last year, the Westar Subcommittee put out a call for astronomers across Canada and received 8 applications from the now Westar Astronomers who continue to undergo training in EDI and Indigenous awareness. Moreover, the Westar Subcommittee began a partnership with Connected North, an initiative to host workshops and sessions virtually at Indigenous schools across northern Canada, now with three activities offered through the Westar Program listed on their catalogue. Collaborations with other Canadian STEM outreach initiatives are underway.

Over the last year, we hosted four virtual sessions at remote, Indigenous schools through Connected North, and in collaboration with the Yukon Astronomical Society, we organized a visit to the Old Crow community in northern Yukon. We continue to take requests from the Connected North catalogue, develop workshops with CASCA's Indigenous Engagement Committee for CASCA members and the broader astronomy community, establish and support connections through the Westar astronomers, and develop new connections through collaborations with outreach initiatives and Indigenous schools.

\hypertarget{poster-abstract-talk-133}{}\label{poster-abstract-page-talk-133}
\subsection*{Poster 48. Procedural Star Generation in SpaceEngine: The Universe Simulator}
\textbf{Presenter:} Megan Tannock\\
\textbf{Poster number:} 48\\
\textbf{Field:} Education and Public Outreach

SpaceEngine is an interactive virtual universe that allows users to explore space on their home computer. The software uses realistic physics and effects to create an accurate, high-definition, full-scale representation of objects in space. In this poster, we introduce the software, highlight its tools for learning about space, and outline the process for populating galaxies in the SpaceEngine universe with billions of stars. We describe our star generation procedure which includes deriving positions for stars in 3D space, calculating physical parameters (luminosity, mass, radius, etc.), and filling in missing data using procedural generation.

To depict space as understood by modern science, our process integrates hundreds of thousands of real objects from the astronomy literature with procedurally-generated data based on established astrophysical relationships. We also describe how we create and render a universe on-the-fly within the constraints of graphics and performance, which enables the software to run on a typical gaming computer. The resulting visualizations are used by educators, science communicators, and professional astronomy organizations, such as ESO, to support presentations, outreach activities, and public engagement.

\hypertarget{poster-abstract-talk-188}{}\label{poster-abstract-page-talk-188}
\subsection*{Poster 49. Determining Variability in High Velocity Quasar Outflows Using SimBAL}
\textbf{Presenter:} James Smith\\
\textbf{Poster number:} 49\\
\textbf{Field:} Supermassive black hole

Quasars are some of the highest energy objects we regularly observe in the universe. Extremely hot ionized gasses around actively accreting supermassive black holes (SMBH) glow brightly enough to be seen across massive distances, and their spectra can hold important details to their growth and evolution. The presence of broad absorption lines (BALs) in a quasars spectrum is evidence for outflowing gasses often travelling at tens of thousands of kilometres per second. Variation in BALs can be seen over only a few years in the observed frame, a remarkably brief timespan for extragalactic phenomena. Studying this variation can provide insight into the growth of SMBHs, and so we have been using a spectral fitting package SimBAL to model BALs and better understand their physical properties. SimBAL fits parameters to quasar spectra including the velocity of the gas, the column density, and covering fraction, that allow us to physically characterize the variability we observe in BALs. We identified a sample of quasars observed in the Sloan Digital Sky Survey that were found to have BALs which varied over multiple epochs. We present case studies of individual quasars that illustrate the utility of SimBAL for revealing the cause of absorption-line variability.

\hypertarget{poster-abstract-talk-183}{}\label{poster-abstract-page-talk-183}
\subsection*{Poster 50. LISA Massive Black Hole Merger Host Galaxies Display a Past Starburst}
\textbf{Presenter:} Alys-Ann Surprenant-Coache\\
\textbf{Poster number:} 50\\
\textbf{Field:} Supermassive black hole

The Laser Interferometer Space Antenna (LISA) mission will detect mergers of massive black hole binary (MBHB) systems at the centres of relatively low-mass galaxies for the first time, but identifying their host galaxies within the large gravitational-wave localization volume presents a major challenge. Since these lower-mass host galaxies are likelier to be gas-rich, we investigate whether the stellar populations of MBHB merger host galaxies detectable by LISA contain signatures of a recent starburst, which may be expected if the MBHB resulted from a major galaxy merger. We use the Romulus25 cosmological simulation to select samples of MBHB mergers, and produce synthetic optical spectra of their host galaxies through stellar population synthesis and dust radiative transfer. We find that the MBHB merger host galaxies often display a starburst $\sim$2 to 3 Gyrs before the MBHB merger, and systematically low star-formation rates at the time of the MBHB merger, in contrast to a mass- and redshift-matched control sample of simulated galaxies without recent MBHB mergers. Our results suggest that observationally, rest-frame optical spectra of galaxies within LISA localization volumes can aid in identifying the exact MBHB host galaxy, by probing the rapidly-declining star-formation rate over the past $\sim$1 Gyr through various spectral features.

\hypertarget{poster-abstract-talk-313}{}\label{poster-abstract-page-talk-313}
\subsection*{Poster 51. An Investigation on the Breakout Properties of Aspherical Relativistic Supernovae}
\textbf{Presenter:} Ari Pollak\\
\textbf{Poster number:} 51\\
\textbf{Fields:} High Energy Astrophysics $\cdot$ Transients

A relativistic blast wave driven by a bipolar jet at the centre of a stripped-envelope star will break out of the surface radially, emitting high-energy photons and the fastest ejecta, setting the stage for possible X-Ray Flash and even Gamma Ray Burst emission. As the blast front slows to trans- or non-relativistic velocities, the flow will partially become tangential, allowing the front to wrap around the progenitor star and collide with ejecta from the opposing side, creating another high-energy transient that we may observe. Through 2D axisymmetric SRMHD simulations with the Athena++ code, we will explore these relativistic aspherical supernovae, and by defining parameters of our blast wave, we can control how it evolves as it approaches the star's surface and classify the shock breakout. We may then compare their features with those of high-energy transients such as SN2008/XRT 080109 and others which have since been observed by the Einstein Probe.

\hypertarget{poster-abstract-talk-100}{}\label{poster-abstract-page-talk-100}
\subsection*{Poster 52. Joint Chandra and NuSTAR Analysis of M87 During the 2021-2025 Event Horizon Telescope Multi-Wavelength Campaigns}
\textbf{Presenter:} Christopher Sheridan\\
\textbf{Poster number:} 52\\
\textbf{Fields:} High Energy Astrophysics $\cdot$ Supermassive black hole

The nearby elliptical galaxy M87 is a mainstay in the study of active galactic nuclei and the evolution and dynamics of supermassive black holes (SMBH) and their outflows. The proximity of M87 and size of its prominent central SMBH, M87*, place it as one of only two SMBHs with large enough angular diameters to be imaged by the Event Horizon Telescope (EHT), alongside Sagittarius A* in our Milky Way. Since 2017, the EHT has mounted near-annual campaigns, accompanied by substantial contemporaneous multi-wavelength (MWL) observation, to study the year-to-year variability of M87, its SMBH, and the various physical properties near and far from the event horizon. I have conducted joint X-ray spectral analysis of Chandra and NuSTAR observations during the 2021-2025 MWL campaigns and will present the results of my analysis, including spectral and X-ray flux variability of the core, prominent knot HST-1, and the outer jet. I will also discuss the implications of these results in both a MWL context and a historical context, with archival Chandra observations of M87 over the last $\sim$two decades.

\hypertarget{poster-abstract-talk-30}{}\label{poster-abstract-page-talk-30}
\subsection*{Poster 53. Investigating the Lithium Plateau with Gaia DR4 and Bayesian Regression}
\textbf{Presenter:} Russell Huang\\
\textbf{Poster number:} 53\\
\textbf{Fields:} Cosmology and Early Universe $\cdot$ Stars and Stellar cluster formation

The cosmological lithium problem is one of the biggest discrepancies between Big Bang nucleosynthesis predictions and stellar observations. In this paper, we measure the first Spite Plateau data with the latest GALAH DR4 survey. After we apply cuts to the data to isolate warm, metal-poor dwarf stars, we analyze lithium abundance compared to metallicity using the Bayesian regression code ROXY, which also accounts for uncertainties in Spite Plateau data. This allows us to calculate the slope, intercept, and scatter of Spite Plateau stars across varying metallicity cuts. Across all metallicity cuts to [Fe/H] $<$ -2.0, the slopes remain at zero within uncertainties, confirming that the Spite Plateau has a flat relationship at low metallicity. These findings demonstrate that even with modern surveys and advanced Bayesian regression tools, there is no explanation for the Spite Plateau. Even with modern tools, the oldest stars appear to have near constant lithium with a strong discrepancy with the value expected from Big Bang Nucleosynthesis, which continues to be an open problem in cosmology and stellar astrophysics. Index Terms: Big Bang Nucleosynthesis, Spite Plateau, Lithium Abundance, Bayesian Regression, Stellar Astrophysics

\hypertarget{poster-abstract-talk-214}{}\label{poster-abstract-page-talk-214}
\subsection*{Poster 54. The Mapper of the IGM Spin Temperature (MIST) Experiment and Soil Characterization in the Canadian High Arctic}
\textbf{Presenter:} Francis McGee\\
\textbf{Poster number:} 54\\
\textbf{Fields:} Cosmology and Early Universe $\cdot$ Instrumentation

The Mapper of the IGM Spin Temperature experiment (MIST) is a cosmology experiment in the Canadian high Arctic measuring the 21-cm global signal produced from neutral hydrogen in the very early universe. A previous global 21-cm experiment called EDGES detected an anomalous flattened absorption feature centered at 78 MHz with an amplitude of around 500 mK which has sometimes been attributed to systematics related to the ground plane beneath the antenna. To remedy this source of systematic error, we perform our experiment without a ground plane, requiring detailed knowledge of the soil beneath the antenna.

We are designing a custom open-source ground penetrating radar (GPR) system to determine the soil parameters below the antenna. The GPR system is drone-mounted due to the rough terrain of the Arctic where conventional summer GPR measurements are difficult without a snowpack to place equipment on. The GPR system uses an RFSoC 4x2 FPGA and custom-designed antennas alongside custom firmware to provide complete waveform capture.

To facilitate these measurements and inform the design of our drone, we perform a series of experiments using the FEKO electromagnetic simulation package to determine the directivity of the antenna above various heterogeneous soil layouts simulating realistic Arctic conditions. We discuss our findings from these simulations and how our ability to determine the global 21-cm signal is limited by our knowledge of the surrounding soil. Informed by the results of these simulations, we outline version 2 of our ground penetrating radar drone.

\hypertarget{poster-abstract-talk-235}{}\label{poster-abstract-page-talk-235}
\subsection*{Poster 55. Insights from the Largest Radio Polarization Catalog of Fast Radio Bursts from CHIME}
\textbf{Presenter:} Mason Ng\\
\textbf{Poster number:} 55\\
\textbf{Field:} Transients

Fast radio bursts (FRBs) are very brief (millisecond-duration), yet bright flashes of radio emission from all over the sky with yet unknown physical origins. Polarimetry provides us with a powerful tool to probe the emission mechanisms and local magneto-ionic environments of FRBs. I will present an overview of the polarisation properties of approximately 2,100 FRBs detected by the Canadian Hydrogen Intensity Mapping Experiment (CHIME) between 9 December 2018 and 15 September 2023, operating at 400-800 MHz. These FRBs originate from over 1,500 unique sources and represent the largest catalog of FRB radio polarimetric measurements from a single instrument to date. I will highlight some key results from the catalog, including a comparison of polarimetric properties between repeating and apparently non-repeating FRBs, the diversity of polarisation angle swings in the bursts, as well as FRBs with high Faraday rotation measures ($>$ 1000 rad/m$^2$) that suggest extreme local environments with strong magnetic fields and/or high electron number densities. I will conclude with a discussion of the implications for FRB emission models and their environments.

\hypertarget{poster-abstract-talk-266}{}\label{poster-abstract-page-talk-266}
\subsection*{Poster 56. CanDIAPL Dashboard for LSST Progress}
\textbf{Presenter:} Anna Ordog\\
\textbf{Poster number:} 56\\
\textbf{Field:} Next generation observatories

The Canadian Data-Intensive Astrophysics PLatform (CanDIAPL) is an infrastructure project providing cloud-computing and specialised software for processing data from the Vera C. Rubin telescope and SKA precursors. The Rubin Observatorys Legacy Survey of Space and Time (LSST) will produce several hundred images of each field covering the entire southern sky over its 10-year observing plan, and will contribute invaluable data to a wide range of science questions. Tracking the progress of data collected over the course of the LSSTs complex observing schedule will be important for ensuring efficient access to the relevant information for many applications. While alert brokers have already been deployed for transient and variable sources detected in LSST difference images, astronomers may need to track accumulation of data on specific static targets or regions of the sky, or particular classes of objects. Some analysis applications may rely on individual visit images accumulated until a specific sensitivity is achieved rather than the coadds from yearly data releases, or require prompt follow-up and cross-matching with radio datasets. We present a prototype of an interactive web-based dashboard for users to monitor their selected targets that will also include an email alert system for notifying users when their targets have met selected thresholds. We invite feedback from the Canadian LSST community regarding features and use-cases that would be of interest to incorporate into this platform as we enter the exciting LSST era.
